%%%%%%%%%%%%%%%%%%%%%%%%%%%%%%%%%%%%%%%%%%%%%%%%%%%%%%%%%%%%%%%%%%%%
%  Algèbre Abstraite I — Groupes et Anneaux
%  Notes de Cours — Licence L2
%  Partie 1 : Chapitres 1–2
%%%%%%%%%%%%%%%%%%%%%%%%%%%%%%%%%%%%%%%%%%%%%%%%%%%%%%%%%%%%%%%%%%%%
\documentclass[12pt,a4paper,openany]{book}

%% ── Encodage et langue ──────────────────────────────────────────
\usepackage[utf8]{inputenc}
\usepackage[T1]{fontenc}
\usepackage[french]{babel}

%% ── Mathématiques ───────────────────────────────────────────────
\usepackage{amsmath,amssymb,amsthm,mathtools,mathrsfs}

%% ── Diagrammes et figures ───────────────────────────────────────
\usepackage{tikz,tikz-cd}
\usetikzlibrary{calc,positioning,arrows.meta,shapes.geometric,decorations.pathreplacing}

%% ── Boîtes colorées ────────────────────────────────────────────
\usepackage[most]{tcolorbox}
\tcbuselibrary{theorems,skins,breakable}

%% ── Mise en page ────────────────────────────────────────────────
\usepackage[margin=2.5cm]{geometry}
\usepackage{fancyhdr}
\usepackage{enumitem}
\usepackage{makeidx}
\usepackage{booktabs}
\usepackage{caption}
\usepackage{xcolor}
\usepackage[colorlinks=true,linkcolor=blue!60!black,citecolor=green!50!black,urlcolor=blue!70!black]{hyperref}

%% ── Couleurs ────────────────────────────────────────────────────
\definecolor{defblue}{RGB}{0,80,160}
\definecolor{thmred}{RGB}{180,30,30}
\definecolor{propgreen}{RGB}{20,120,50}
\definecolor{rembg}{RGB}{255,248,230}
\definecolor{exbg}{RGB}{235,245,255}

%% ── Environnements de théorèmes (amsthm) ───────────────────────
\theoremstyle{definition}
\newtheorem{definition}{Définition}[chapter]
\newtheorem{theoreme}{Théorème}[chapter]
\newtheorem{proposition}{Proposition}[chapter]
\newtheorem{lemme}{Lemme}[chapter]
\newtheorem{corollaire}{Corollaire}[chapter]
\newtheorem{remarque}{Remarque}[chapter]
\newtheorem{exemple}{Exemple}[chapter]
\newtheorem{exercice}{Exercice}[chapter]

%% ── Habillage tcolorbox des environnements ──────────────────────
\tcolorboxenvironment{definition}{%
  enhanced, breakable,
  colback=blue!5, colframe=defblue,
  boxrule=0.8pt, arc=2pt,
  left=6pt, right=6pt, top=6pt, bottom=6pt,
  fonttitle=\bfseries
}
\tcolorboxenvironment{theoreme}{%
  enhanced, breakable,
  colback=red!5, colframe=thmred,
  boxrule=0.8pt, arc=2pt,
  left=6pt, right=6pt, top=6pt, bottom=6pt,
  fonttitle=\bfseries
}
\tcolorboxenvironment{proposition}{%
  enhanced, breakable,
  colback=green!5, colframe=propgreen,
  boxrule=0.8pt, arc=2pt,
  left=6pt, right=6pt, top=6pt, bottom=6pt,
  fonttitle=\bfseries
}
\tcolorboxenvironment{lemme}{%
  enhanced, breakable,
  colback=yellow!8, colframe=orange!70!black,
  boxrule=0.8pt, arc=2pt,
  left=6pt, right=6pt, top=6pt, bottom=6pt,
  fonttitle=\bfseries
}
\tcolorboxenvironment{corollaire}{%
  enhanced, breakable,
  colback=green!3, colframe=propgreen!70,
  boxrule=0.6pt, arc=2pt,
  left=6pt, right=6pt, top=6pt, bottom=6pt,
  fonttitle=\bfseries
}
\tcolorboxenvironment{remarque}{%
  enhanced, breakable,
  colback=rembg, colframe=orange!40!black,
  boxrule=0.5pt, arc=2pt,
  left=6pt, right=6pt, top=4pt, bottom=4pt,
  fonttitle=\bfseries
}
\tcolorboxenvironment{exemple}{%
  enhanced, breakable,
  colback=exbg, colframe=blue!40!black,
  boxrule=0.5pt, arc=2pt,
  left=6pt, right=6pt, top=4pt, bottom=4pt,
  fonttitle=\bfseries
}
\tcolorboxenvironment{exercice}{%
  enhanced, breakable,
  colback=gray!5, colframe=gray!60!black,
  boxrule=0.5pt, arc=2pt,
  left=6pt, right=6pt, top=4pt, bottom=4pt,
  fonttitle=\bfseries
}

%% ── Commandes personnalisées ────────────────────────────────────
% Corps et ensembles
\newcommand{\N}{\mathbb{N}}
\newcommand{\Z}{\mathbb{Z}}
\newcommand{\Q}{\mathbb{Q}}
\newcommand{\R}{\mathbb{R}}
\newcommand{\C}{\mathbb{C}}
\newcommand{\F}{\mathbb{F}}

% Groupes classiques
\DeclareMathOperator{\Sym}{Sym}
\DeclareMathOperator{\Alt}{Alt}
\DeclareMathOperator{\Aut}{Aut}
\DeclareMathOperator{\Inn}{Inn}
\DeclareMathOperator{\Out}{Out}
\DeclareMathOperator{\End}{End}
\DeclareMathOperator{\Hom}{Hom}
\DeclareMathOperator{\Ker}{Ker}
\DeclareMathOperator{\Id}{Id}
\renewcommand{\Im}{\operatorname{Im}}
\DeclareMathOperator{\ord}{ord}
\DeclareMathOperator{\car}{car}
\DeclareMathOperator{\pgcd}{pgcd}
\DeclareMathOperator{\ppcm}{ppcm}
\DeclareMathOperator{\GL}{GL}
\DeclareMathOperator{\SL}{SL}

% Raccourcis
\newcommand{\iso}{\simeq}
\newcommand{\normal}{\trianglelefteq}
\newcommand{\snormal}{\triangleleft}
\newcommand{\acts}{\curvearrowright}
\newcommand{\abs}[1]{\lvert #1 \rvert}
\newcommand{\gen}[1]{\langle #1 \rangle}
\newcommand{\idx}[2]{[#1:#2]}

%% ── En-têtes et pieds de page ───────────────────────────────────
\pagestyle{fancy}
\fancyhf{}
\fancyhead[LE]{\slshape\leftmark}
\fancyhead[RO]{\slshape\rightmark}
\fancyfoot[C]{\thepage}
\renewcommand{\headrulewidth}{0.4pt}

%% ── Index ───────────────────────────────────────────────────────
\makeindex

%%%%%%%%%%%%%%%%%%%%%%%%%%%%%%%%%%%%%%%%%%%%%%%%%%%%%%%%%%%%%%%%%%%%
%  PAGE DE TITRE
%%%%%%%%%%%%%%%%%%%%%%%%%%%%%%%%%%%%%%%%%%%%%%%%%%%%%%%%%%%%%%%%%%%%
\begin{document}

\begin{titlepage}
\begin{tikzpicture}[remember picture, overlay]
  \fill[left color=blue!15!white, right color=blue!5!white]
    (current page.south west) rectangle (current page.north east);
  \fill[color=orange!70!yellow]
    ([yshift=-2cm]current page.north west) rectangle ([yshift=-2.3cm]current page.north east);
  \fill[color=blue!80!black, rounded corners=8pt]
    ([xshift=3cm, yshift=-4cm]current page.north west) rectangle
    ([xshift=-3cm, yshift=-10cm]current page.north east);
  \node[anchor=center, text=white, font=\Huge\bfseries]
    at ([yshift=-6cm]current page.north) {Alg\`ebre Abstraite I};
  \node[anchor=center, text=orange!80!yellow, font=\Large]
    at ([yshift=-7.5cm]current page.north) {Notes de Cours};
  \node[anchor=center, text=white!80, font=\large]
    at ([yshift=-8.8cm]current page.north) {Licence L2 --- 2025--2026};
  \node[anchor=center, text=blue!80!black, font=\Large\itshape]
    at ([yshift=-12cm]current page.north) {Ya\'e Ulrich Gaba};
  \draw[orange!70!yellow, line width=2pt]
    ([xshift=5cm, yshift=-13.5cm]current page.north west) --
    ([xshift=-5cm, yshift=-13.5cm]current page.north east);
  \node[anchor=center, text width=12cm, align=center, text=blue!60!black, font=\itshape]
    at ([yshift=-16cm]current page.north) {<<\,Les math\'ematiques sont un langage plus parfait que le langage ordinaire pour exprimer les relations entre les choses.\,>>\\[4pt]--- Nicolas Bourbaki};
  \node[anchor=south, text=gray, font=\small]
    at ([yshift=1.5cm]current page.south) {\today};
\end{tikzpicture}
\end{titlepage}
%%%%%%%%%%%%%%%%%%%%%%%%%%%%%%%%%%%%%%%%%%%%%%%%%%%%%%%%%%%%%%%%%%%%
%  TABLE DES MATIÈRES
%%%%%%%%%%%%%%%%%%%%%%%%%%%%%%%%%%%%%%%%%%%%%%%%%%%%%%%%%%%%%%%%%%%%
\tableofcontents
\newpage

%%%%%%%%%%%%%%%%%%%%%%%%%%%%%%%%%%%%%%%%%%%%%%%%%%%%%%%%%%%%%%%%%%%%
%  PRÉFACE
%%%%%%%%%%%%%%%%%%%%%%%%%%%%%%%%%%%%%%%%%%%%%%%%%%%%%%%%%%%%%%%%%%%%
\chapter*{Préface}
\addcontentsline{toc}{chapter}{Préface}
\markboth{Préface}{Préface}

L'algèbre abstraite, telle que nous la connaissons aujourd'hui, est le fruit d'une longue
maturation intellectuelle qui s'étend sur plus de deux siècles. Ses origines se trouvent
dans des questions très concrètes~: la résolution d'équations polynomiales par radicaux,
l'étude des symétries de figures géométriques, et l'arithmétique des nombres entiers.

C'est Joseph-Louis \textbf{Lagrange}\index{Lagrange} qui, dans ses \emph{Réflexions sur la
résolution algébrique des équations} (1770--1771), mit en évidence pour la première fois
le rôle des permutations des racines d'un polynôme dans la compréhension de sa
résolubilité. En étudiant systématiquement les \og résolvantes\fg{} associées aux
équations de degré~3 et~4, Lagrange identifia un phénomène fondamental~: le nombre de
valeurs distinctes prises par une fonction rationnelle des racines, lorsqu'on permute
celles-ci, divise le nombre total de permutations. Ce résultat, ancêtre du théorème qui
porte aujourd'hui son nom, allait devenir la pierre angulaire de la théorie des groupes.

Niels Henrik \textbf{Abel}\index{Abel} (1824) démontra l'impossibilité de résoudre par
radicaux l'équation générale de degré~$\geq 5$, puis Évariste \textbf{Galois}%
\index{Galois}, dans ses écrits tragiquement brefs (1831--1832), accomplit le pas
décisif~: il associa à toute équation polynomiale un \emph{groupe} de permutations ---
le groupe de Galois --- et montra que la résolubilité par radicaux équivaut à une
propriété purement algébrique de ce groupe (sa résolubilité). Ce faisant, Galois inventa
simultanément la notion de groupe, celle de sous-groupe distingué, et celle de groupe
quotient, anticipant de plusieurs décennies le développement formel de ces concepts.

Au cours du XIXe siècle, Augustin-Louis \textbf{Cauchy}\index{Cauchy} systématisa la
théorie des permutations, Arthur \textbf{Cayley}\index{Cayley} proposa la définition
abstraite d'un groupe (1854), et Camille \textbf{Jordan}\index{Jordan} publia son
\emph{Traité des substitutions et des équations algébriques} (1870), premier ouvrage de
synthèse sur la théorie des groupes. Parallèlement, Richard \textbf{Dedekind}%
\index{Dedekind} introduisit les notions d'idéal et de module, jetant les bases de
l'algèbre commutative moderne.

Le XXe siècle vit l'avènement de l'algèbre abstraite proprement dite. Emmy
\textbf{Noether}\index{Noether}, dans ses travaux fondateurs des années 1920, démontra
que la bonne façon d'étudier les structures algébriques consiste à se concentrer sur les
\emph{morphismes} --- les applications qui préservent la structure --- plutôt que sur les
éléments individuels. Son approche, systématisée par Emil \textbf{Artin}\index{Artin}
et Bartel Leendert \textbf{van der Waerden}\index{van der Waerden} dans le célèbre
traité \emph{Moderne Algebra} (1930--1931), donna naissance au point de vue qui domine
l'algèbre contemporaine~: l'étude des structures (groupes, anneaux, corps, modules) et
des applications qui les relient (homomorphismes, isomorphismes).

\medskip

\textbf{Pourquoi l'abstraction~?} Le lecteur pourrait légitimement se demander pourquoi
nous étudions des \og groupes\fg{} et des \og anneaux\fg{} en toute généralité, plutôt
que de nous contenter des exemples concrets ($\Z$, $\R$, les matrices, les
permutations). La réponse est triple~:
\begin{enumerate}[label=(\roman*)]
  \item \emph{Unification.} Un même théorème, démontré une fois dans le cadre abstrait,
    s'applique simultanément à une multitude de situations concrètes. Le théorème de
    Lagrange, par exemple, donne d'un seul coup le petit théorème de Fermat, le théorème
    d'Euler, et des contraintes sur l'ordre des éléments de tout groupe fini.
  \item \emph{Clarification.} L'abstraction permet d'identifier précisément quelles
    hypothèses sont nécessaires pour chaque résultat, séparant l'essentiel de
    l'accessoire.
  \item \emph{Découverte.} Le cadre abstrait suggère des constructions (produits,
    quotients, extensions) et des questions (classification, simplicité, représentation)
    qui n'auraient pas émergé naturellement dans un cadre concret.
\end{enumerate}

Le fil conducteur de ce cours est l'idée noethérienne que \emph{comprendre une structure
algébrique, c'est comprendre ses morphismes}. Nous étudierons chaque structure (groupe,
anneau, idéal, corps) en mettant systématiquement l'accent sur les homomorphismes, les
sous-structures, les quotients, et les théorèmes d'isomorphisme qui les relient. Les
diagrammes commutatifs, omniprésents dans ce texte, sont l'expression visuelle de cette
philosophie.

\medskip

Ce premier volume couvre la théorie des groupes (chapitres~1 à~5) et l'introduction aux
anneaux et corps (chapitres~6 à~8). Il est destiné aux étudiants de deuxième année de
licence (L2) et suppose acquises les notions de base de la théorie des ensembles, des
applications, et des relations d'équivalence.

%%%%%%%%%%%%%%%%%%%%%%%%%%%%%%%%%%%%%%%%%%%%%%%%%%%%%%%%%%%%%%%%%%%%
%  NOTATIONS
%%%%%%%%%%%%%%%%%%%%%%%%%%%%%%%%%%%%%%%%%%%%%%%%%%%%%%%%%%%%%%%%%%%%
\chapter*{Notations}
\addcontentsline{toc}{chapter}{Notations}
\markboth{Notations}{Notations}

\begin{tabular}{@{}ll@{}}
\toprule
\textbf{Notation} & \textbf{Signification} \\
\midrule
$\N$ & Ensemble des entiers naturels $\{0,1,2,\ldots\}$ \\
$\Z$ & Anneau des entiers relatifs \\
$\Q$ & Corps des nombres rationnels \\
$\R$ & Corps des nombres réels \\
$\C$ & Corps des nombres complexes \\
$\F_q$ & Corps fini à $q$ éléments \\
$\Z/n\Z$ & Anneau (ou groupe) des entiers modulo $n$ \\
$(\Z/n\Z)^\times$ & Groupe des éléments inversibles de $\Z/n\Z$ \\
$\abs{G}$ & Ordre (cardinal) du groupe $G$ \\
$\ord(g)$ & Ordre de l'élément $g$ dans un groupe \\
$\gen{g}$ & Sous-groupe engendré par $g$ \\
$\gen{S}$ & Sous-groupe engendré par le sous-ensemble $S$ \\
$H \leq G$ & $H$ est un sous-groupe de $G$ \\
$H \normal G$ & $H$ est un sous-groupe distingué de $G$ \\
$\idx{G}{H}$ & Indice de $H$ dans $G$ \\
$G/N$ & Groupe quotient de $G$ par $N$ \\
$G \times H$ & Produit direct de $G$ et $H$ \\
$G \iso H$ & $G$ est isomorphe à $H$ \\
$\GL_n(K)$ & Groupe linéaire général d'ordre $n$ sur $K$ \\
$\SL_n(K)$ & Groupe spécial linéaire d'ordre $n$ sur $K$ \\
$S_n$ & Groupe symétrique sur $n$ éléments \\
$A_n$ & Groupe alterné sur $n$ éléments \\
$D_n$ & Groupe diédral d'ordre $2n$ \\
$Q_8$ & Groupe des quaternions \\
$V_4$ & Groupe de Klein à quatre éléments \\
$Z(G)$ & Centre du groupe $G$ \\
$[G,G]$ & Sous-groupe dérivé (des commutateurs) de $G$ \\
$\Aut(G)$ & Groupe des automorphismes de $G$ \\
$\Inn(G)$ & Groupe des automorphismes intérieurs de $G$ \\
$\Hom(G,H)$ & Ensemble des homomorphismes de $G$ dans $H$ \\
$\Ker(f)$ & Noyau de l'homomorphisme $f$ \\
$\Im(f)$ & Image de l'homomorphisme $f$ \\
$\Id$ & Application identité \\
$\car(A)$ & Caractéristique de l'anneau $A$ \\
$\pgcd(a,b)$ & Plus grand commun diviseur \\
$\ppcm(a,b)$ & Plus petit commun multiple \\
\bottomrule
\end{tabular}


%%%%%%%%%%%%%%%%%%%%%%%%%%%%%%%%%%%%%%%%%%%%%%%%%%%%%%%%%%%%%%%%%%%%
%%%%%%%%%%%%%%%%%%%%%%%%%%%%%%%%%%%%%%%%%%%%%%%%%%%%%%%%%%%%%%%%%%%%
%%
%%  CHAPITRE 1 : GROUPES — DÉFINITIONS, EXEMPLES, PREMIÈRES PROPRIÉTÉS
%%
%%%%%%%%%%%%%%%%%%%%%%%%%%%%%%%%%%%%%%%%%%%%%%%%%%%%%%%%%%%%%%%%%%%%
%%%%%%%%%%%%%%%%%%%%%%%%%%%%%%%%%%%%%%%%%%%%%%%%%%%%%%%%%%%%%%%%%%%%

\chapter{Groupes~: définitions, exemples et premières propriétés}
\label{chap:groupes-definitions}

\section{Introduction historique}
\label{sec:ch1-intro-historique}

La notion de \emph{groupe}\index{groupe} est l'une des plus fécondes de l'ensemble des
mathématiques. Elle naît de l'observation que de nombreux objets mathématiques possèdent
des \emph{symétries}, et que l'ensemble de ces symétries, muni de la composition, obéit
à des lois algébriques remarquablement simples et universelles.

Historiquement, les groupes apparaissent sous trois formes~:
\begin{enumerate}[label=(\alph*)]
  \item \textbf{Les groupes de permutations.} Lagrange\index{Lagrange} (1770) étudia les
    permutations des racines d'un polynôme~; Cauchy\index{Cauchy} (1815, 1844)
    systématisa la théorie des substitutions~; Cayley\index{Cayley} (1854) montra que
    tout groupe (abstrait) est isomorphe à un sous-groupe d'un groupe de permutations
    (théorème de Cayley).
  \item \textbf{Les groupes de symétries géométriques.} Les symétries d'un polygone
    régulier forment un \emph{groupe diédral}~; celles d'un polyèdre régulier, un
    sous-groupe de $\GL_3(\R)$. Klein\index{Klein} (1872), dans son \emph{Programme
    d'Erlangen}, proposa de \emph{définir} une géométrie par son groupe de
    transformations.
  \item \textbf{Les groupes de nombres.} L'ensemble $\Z$ muni de l'addition, le groupe
    $(\Z/n\Z, +)$, le groupe multiplicatif $\R^* = \R \setminus\{0\}$~: autant
    d'exemples issus de l'arithmétique et de l'analyse.
\end{enumerate}

L'abstraction consiste à dégager les propriétés communes à tous ces exemples et à les
prendre comme \emph{axiomes}.

\section{Loi de composition interne}
\label{sec:ch1-lci}

\begin{definition}[Loi de composition interne]
\label{def:lci}
\index{loi de composition interne}
Soit $E$ un ensemble non vide. Une \emph{loi de composition interne} (ou
\emph{opération binaire interne}) sur $E$ est une application
\[
  \star : E \times E \longrightarrow E, \qquad (a,b) \longmapsto a \star b.
\]
\end{definition}

\begin{remarque}
La condition essentielle est que $a \star b \in E$ pour tous $a, b \in E$~: on dit que
$E$ est \emph{stable}\index{stabilité} (ou \emph{fermé}) pour~$\star$.
\end{remarque}

\begin{exemple}[Exemples de lois de composition interne]
\label{ex:lci-exemples}
\begin{enumerate}[label=(\roman*)]
  \item L'addition $+$ et la multiplication $\times$ sont des lois de composition
    internes sur $\Z$, $\Q$, $\R$, $\C$.
  \item La soustraction $-$ est une loi de composition interne sur $\Z$, mais
    \emph{pas} sur $\N$ (car $2 - 5 = -3 \notin \N$).
  \item La composition des applications $\circ$ est une loi de composition interne sur
    l'ensemble des bijections d'un ensemble $E$ dans lui-même.
  \item La multiplication matricielle est une loi de composition interne sur $M_n(K)$,
    l'ensemble des matrices carrées d'ordre $n$ à coefficients dans un corps $K$.
\end{enumerate}
\end{exemple}

\section{Définition d'un groupe}
\label{sec:ch1-def-groupe}

\begin{definition}[Groupe]
\label{def:groupe}
\index{groupe!définition}
Un \emph{groupe} est un couple $(G, \star)$ où $G$ est un ensemble non vide et $\star$
est une loi de composition interne sur $G$ vérifiant les axiomes suivants~:
\begin{enumerate}[label=\textbf{(G\arabic*)}]
  \item \textbf{Associativité.}\index{associativité} Pour tous $a, b, c \in G$~:
    \[
      (a \star b) \star c = a \star (b \star c).
    \]
  \item \textbf{Élément neutre.}\index{élément neutre} Il existe un élément $e \in G$
    tel que, pour tout $a \in G$~:
    \[
      e \star a = a \star e = a.
    \]
  \item \textbf{Symétrique.}\index{inverse}\index{symétrique} Pour tout $a \in G$, il
    existe un élément $a' \in G$ tel que~:
    \[
      a \star a' = a' \star a = e.
    \]
\end{enumerate}
On note usuellement $a' = a^{-1}$ (notation multiplicative) ou $a' = -a$ (notation
additive).
\end{definition}

\begin{remarque}
\label{rem:quatre-axiomes}
Certains auteurs incluent explicitement la stabilité (le fait que $\star$ soit une loi de
composition interne sur $G$) comme axiome \textbf{(G0)}, ce qui donne quatre axiomes.
Dans notre formulation, la stabilité fait partie de la définition de «~loi de composition
interne~».
\end{remarque}

\section{Premières propriétés}
\label{sec:ch1-premieres-proprietes}

\begin{theoreme}[Unicité de l'élément neutre]
\label{thm:unicite-neutre}
\index{élément neutre!unicité}
Soit $(G, \star)$ un groupe. L'élément neutre est unique.
\end{theoreme}

\begin{proof}
Supposons que $e$ et $e'$ soient deux éléments neutres de $G$. Alors~:
\[
  e = e \star e' = e',
\]
la première égalité provenant du fait que $e'$ est neutre (donc $e \star e' = e$) et la
seconde du fait que $e$ est neutre (donc $e \star e' = e'$). Ainsi $e = e'$.
\end{proof}

\begin{theoreme}[Unicité de l'inverse]
\label{thm:unicite-inverse}
\index{inverse!unicité}
Soit $(G, \star)$ un groupe et $a \in G$. L'inverse de $a$ est unique.
\end{theoreme}

\begin{proof}
Supposons que $b$ et $c$ soient deux inverses de $a$, c'est-à-dire $a \star b = b \star
a = e$ et $a \star c = c \star a = e$. Alors~:
\[
  b = b \star e = b \star (a \star c) = (b \star a) \star c = e \star c = c.
\]
Donc $b = c$.
\end{proof}

\begin{theoreme}[Inverse d'un produit]
\label{thm:inverse-produit}
\index{inverse!d'un produit}
Soit $(G, \star)$ un groupe et $a, b \in G$. Alors~:
\[
  (a \star b)^{-1} = b^{-1} \star a^{-1}.
\]
\end{theoreme}

\begin{proof}
Il suffit de vérifier que $b^{-1} \star a^{-1}$ satisfait la définition de l'inverse de
$a \star b$. On calcule~:
\[
  (a \star b) \star (b^{-1} \star a^{-1})
  = a \star (b \star b^{-1}) \star a^{-1}
  = a \star e \star a^{-1}
  = a \star a^{-1}
  = e.
\]
De même~:
\[
  (b^{-1} \star a^{-1}) \star (a \star b)
  = b^{-1} \star (a^{-1} \star a) \star b
  = b^{-1} \star e \star b
  = b^{-1} \star b
  = e.
\]
Par unicité de l'inverse (théorème~\ref{thm:unicite-inverse}), on conclut
$(a \star b)^{-1} = b^{-1} \star a^{-1}$.
\end{proof}

\begin{proposition}[Régularité]
\label{prop:regularite}
\index{régularité}
Soit $(G, \star)$ un groupe et $a, b, c \in G$.
\begin{enumerate}[label=(\roman*)]
  \item Si $a \star b = a \star c$, alors $b = c$ \textup{(régularité à gauche)}.
  \item Si $b \star a = c \star a$, alors $b = c$ \textup{(régularité à droite)}.
\end{enumerate}
\end{proposition}

\begin{proof}
(i) Si $a \star b = a \star c$, on compose à gauche par $a^{-1}$~:
\[
  a^{-1} \star (a \star b) = a^{-1} \star (a \star c)
  \implies (a^{-1} \star a) \star b = (a^{-1} \star a) \star c
  \implies e \star b = e \star c
  \implies b = c.
\]
La preuve de (ii) est analogue, en composant à droite par $a^{-1}$.
\end{proof}

\begin{proposition}[Involutivité de l'inversion]
\label{prop:involution-inverse}
Soit $(G, \star)$ un groupe et $a \in G$. Alors $(a^{-1})^{-1} = a$.
\end{proposition}

\begin{proof}
Par définition, $a^{-1} \star a = a \star a^{-1} = e$. Cela signifie que $a$ est un
inverse de $a^{-1}$. Par unicité de l'inverse, $(a^{-1})^{-1} = a$.
\end{proof}

\section{Groupes abéliens}
\label{sec:ch1-abelien}

\begin{definition}[Groupe abélien]
\label{def:abelien}
\index{groupe!abélien}
\index{commutatif|see{groupe abélien}}
Un groupe $(G, \star)$ est dit \emph{abélien} (ou \emph{commutatif}) si, pour tous $a,
b \in G$~:
\[
  a \star b = b \star a.
\]
\end{definition}

\begin{remarque}
Le nom «~abélien~» est un hommage à Niels Henrik Abel\index{Abel}. Par convention, les
groupes abéliens sont souvent notés additivement $(G, +)$, l'élément neutre étant noté
$0$ et l'inverse de $a$ étant noté $-a$.
\end{remarque}

\section{Exemples fondamentaux}
\label{sec:ch1-exemples}

\begin{exemple}[Le groupe $(\Z, +)$]
\label{ex:Z-additif}
\index{Z@$\Z$}
L'ensemble $\Z$ des entiers relatifs, muni de l'addition, est un groupe abélien.
L'élément neutre est $0$~; l'inverse de $n \in \Z$ est $-n$.
\end{exemple}

\begin{exemple}[Le groupe $(\Z/n\Z, +)$]
\label{ex:ZnZ-additif}
\index{Z/nZ@$\Z/n\Z$}
Pour tout entier $n \geq 1$, l'ensemble $\Z/n\Z = \{\bar{0}, \bar{1}, \ldots,
\overline{n-1}\}$ muni de l'addition modulo $n$ est un groupe abélien fini d'ordre $n$.
L'élément neutre est $\bar{0}$~; l'inverse de $\bar{k}$ est $\overline{n-k}$.
\end{exemple}

\begin{exemple}[Le groupe $(\R^*, \times)$]
\label{ex:R-etoile}
\index{R*@$\R^*$}
L'ensemble $\R^* = \R \setminus \{0\}$ muni de la multiplication est un groupe abélien.
L'élément neutre est $1$~; l'inverse de $a$ est $1/a$. On note que $\R$ tout entier ne
forme \emph{pas} un groupe pour la multiplication, car $0$ n'a pas d'inverse.
\end{exemple}

\begin{exemple}[Le groupe linéaire $\GL_n(\R)$]
\label{ex:GLn}
\index{GL@$\GL_n$}
L'ensemble $\GL_n(\R) = \{ A \in M_n(\R) \mid \det(A) \neq 0 \}$ des matrices
inversibles d'ordre $n$, muni du produit matriciel, est un groupe. L'élément neutre est
la matrice identité $I_n$~; l'inverse de $A$ est la matrice inverse $A^{-1}$. Ce groupe
est \textbf{non abélien} dès que $n \geq 2$.
\end{exemple}

\begin{exemple}[Le groupe spécial linéaire $\SL_n(\R)$]
\label{ex:SLn}
\index{SL@$\SL_n$}
L'ensemble $\SL_n(\R) = \{ A \in \GL_n(\R) \mid \det(A) = 1 \}$ est un groupe pour le
produit matriciel. En effet, si $\det(A) = \det(B) = 1$, alors $\det(AB) =
\det(A)\det(B) = 1$ et $\det(A^{-1}) = 1/\det(A) = 1$.
\end{exemple}

\begin{exemple}[Le groupe symétrique $S_n$]
\label{ex:Sn}
\index{groupe!symétrique}\index{S@$S_n$}
Soit $E = \{1, 2, \ldots, n\}$. L'ensemble $S_n$ de toutes les bijections de $E$ dans
$E$ (appelées \emph{permutations}\index{permutation}) muni de la composition est un
groupe d'ordre $n!$. L'élément neutre est l'identité $\Id_E$~; l'inverse d'une
permutation $\sigma$ est la permutation réciproque $\sigma^{-1}$. Le groupe $S_n$ est
non abélien dès que $n \geq 3$.
\end{exemple}

\begin{exemple}[Le groupe diédral $D_n$]
\label{ex:Dn}
\index{groupe!diédral}\index{D@$D_n$}
Le \emph{groupe diédral} $D_n$ est le groupe des isométries du plan qui préservent un
polygone régulier à $n$ côtés. Il contient $n$ rotations (d'angles $2k\pi/n$, pour $k =
0, 1, \ldots, n-1$) et $n$ réflexions, soit $2n$ éléments au total. On écrit~:
\[
  D_n = \gen{r, s \mid r^n = s^2 = e, \; srs = r^{-1}},
\]
où $r$ désigne la rotation d'angle $2\pi/n$ et $s$ une réflexion. Le groupe $D_n$ est
non abélien pour $n \geq 3$.
\end{exemple}

\begin{exemple}[Le groupe des quaternions $Q_8$]
\label{ex:Q8}
\index{quaternions}\index{Q8@$Q_8$}
Le \emph{groupe des quaternions} est le groupe d'ordre~$8$~:
\[
  Q_8 = \{\pm 1, \pm i, \pm j, \pm k\}
\]
avec les relations $i^2 = j^2 = k^2 = -1$, $ij = k$, $jk = i$, $ki = j$,
$ji = -k$, $kj = -i$, $ik = -j$. C'est un groupe non abélien. Il fournit un
exemple important d'un groupe dont tous les sous-groupes sont distingués (on dit que
$Q_8$ est un groupe \emph{hamiltonien}\index{groupe!hamiltonien}).
\end{exemple}

\begin{exemple}[Le groupe de Klein $V_4$]
\label{ex:V4}
\index{Klein (groupe de)}\index{V4@$V_4$}
Le \emph{groupe de Klein} (ou \emph{Vierergruppe}) est le groupe abélien d'ordre~$4$~:
\[
  V_4 = \{e, a, b, c\} \quad \text{avec} \quad a^2 = b^2 = c^2 = e, \quad ab = c, \quad
  bc = a, \quad ac = b.
\]
On peut le réaliser comme $V_4 \iso \Z/2\Z \times \Z/2\Z$, ou comme le sous-groupe
$\{\Id, (1\;2)(3\;4), (1\;3)(2\;4), (1\;4)(2\;3)\}$ de $S_4$.
\end{exemple}

\section{Groupes cycliques}
\label{sec:ch1-cycliques}

\begin{definition}[Groupe monogène, groupe cyclique]
\label{def:cyclique}
\index{groupe!cyclique}\index{groupe!monogène}
Un groupe $G$ est dit \emph{monogène} s'il existe un élément $g \in G$ tel que $G =
\gen{g} = \{g^n \mid n \in \Z\}$. Un groupe monogène fini est appelé \emph{cyclique}.
Par abus courant, on qualifie aussi de cyclique tout groupe monogène (fini ou infini).
L'élément $g$ est appelé un \emph{générateur}\index{générateur} de $G$.
\end{definition}

\begin{exemple}
\label{ex:cycliques-exemples}
\begin{enumerate}[label=(\roman*)]
  \item $(\Z, +)$ est un groupe cyclique (infini) engendré par $1$ (ou par $-1$).
  \item $(\Z/n\Z, +)$ est un groupe cyclique d'ordre $n$, engendré par $\bar{1}$.
  \item Le groupe des racines $n$-ièmes de l'unité $\mu_n = \{e^{2ik\pi/n} \mid k =
    0, \ldots, n-1\} \subset \C^*$ est un groupe cyclique d'ordre $n$, engendré par
    $e^{2i\pi/n}$.
\end{enumerate}
\end{exemple}

\begin{theoreme}[Classification des groupes cycliques]
\label{thm:classification-cycliques}
\index{groupe!cyclique!classification}
Tout groupe cyclique est isomorphe soit à $\Z$ (s'il est infini), soit à $\Z/n\Z$ pour
un certain $n \geq 1$ (s'il est fini d'ordre $n$).
\end{theoreme}

\begin{proof}
Soit $G = \gen{g}$ un groupe cyclique et considérons l'application
\[
  \varphi : \Z \longrightarrow G, \qquad n \longmapsto g^n.
\]
C'est un homomorphisme de groupes~: $\varphi(m + n) = g^{m+n} = g^m \cdot g^n =
\varphi(m) \cdot \varphi(n)$. De plus, $\varphi$ est surjective puisque $G = \gen{g}$.

Posons $K = \Ker(\varphi) = \{n \in \Z \mid g^n = e\}$. C'est un sous-groupe de $\Z$.
Or, tout sous-groupe de $\Z$ est de la forme $d\Z$ pour un certain $d \geq 0$
(cf.\ exercice~\ref{exo:sg-Z}).

\textbf{Cas 1~:} $K = \{0\}$, c'est-à-dire $d = 0$. Alors $\varphi$ est injective,
donc bijective, et $G \iso \Z$.

\textbf{Cas 2~:} $K = d\Z$ avec $d \geq 1$. Par le théorème d'isomorphisme
(que nous admettrons ici et démontrerons au chapitre~3), on a~:
\[
  G \iso \Z / K = \Z / d\Z.
\]
De plus, $\abs{G} = \abs{\Z/d\Z} = d$, donc $d = n$ si $G$ est d'ordre $n$. Ainsi $G
\iso \Z/n\Z$.
\end{proof}

\section{Ordre d'un élément}
\label{sec:ch1-ordre}

\begin{definition}[Ordre d'un élément]
\label{def:ordre-element}
\index{ordre!d'un élément}
Soit $(G, \cdot)$ un groupe et $g \in G$. L'\emph{ordre} de $g$, noté $\ord(g)$, est le
plus petit entier $n \geq 1$ tel que $g^n = e$, s'il existe. Si aucun tel entier
n'existe, on dit que $g$ est d'\emph{ordre infini} et on pose $\ord(g) = +\infty$.
\end{definition}

\begin{definition}[Ordre d'un groupe]
\label{def:ordre-groupe}
\index{ordre!d'un groupe}
L'\emph{ordre} d'un groupe $G$, noté $\abs{G}$, est le cardinal de $G$.
\end{definition}

\begin{proposition}
\label{prop:ordre-gen}
Soit $G$ un groupe et $g \in G$. Alors $\ord(g) = \abs{\gen{g}}$.
\end{proposition}

\begin{proof}
Si $\ord(g) = +\infty$, les éléments $g^k$ ($k \in \Z$) sont tous distincts (sinon
$g^{k-l} = e$ pour $k \neq l$, contredisant l'hypothèse), donc $\gen{g}$ est infini.

Si $\ord(g) = n < +\infty$, montrons que $\gen{g} = \{e, g, g^2, \ldots, g^{n-1}\}$,
ce qui donnera $\abs{\gen{g}} = n$. Tout élément de $\gen{g}$ s'écrit $g^k$ pour un
certain $k \in \Z$. Par division euclidienne, $k = nq + r$ avec $0 \leq r < n$, d'où
$g^k = g^{nq+r} = (g^n)^q \cdot g^r = e^q \cdot g^r = g^r$. Donc $\gen{g} \subseteq
\{e, g, \ldots, g^{n-1}\}$. De plus, ces $n$ éléments sont distincts~: si $g^i = g^j$
avec $0 \leq i < j < n$, alors $g^{j-i} = e$ avec $0 < j - i < n$, contredisant la
minimalité de $n$.
\end{proof}

\begin{proposition}
\label{prop:gn-iff-divise}
Soit $g \in G$ un élément d'ordre fini $n$. Alors $g^k = e$ si et seulement si $n
\mid k$.
\end{proposition}

\begin{proof}
Si $n \mid k$, alors $k = nq$ et $g^k = (g^n)^q = e^q = e$. Réciproquement, si $g^k =
e$, écrivons $k = nq + r$ avec $0 \leq r < n$. Alors $g^r = g^{k - nq} = g^k \cdot
(g^n)^{-q} = e \cdot e = e$. Par minimalité de $n$, on a $r = 0$, donc $n \mid k$.
\end{proof}

\section{Tables de Cayley}
\label{sec:ch1-cayley}
\index{table de Cayley}

La \emph{table de Cayley}\index{Cayley!table de} d'un groupe fini $G = \{g_1, \ldots,
g_n\}$ est le tableau dont l'entrée à la ligne~$i$ et à la colonne~$j$ est le produit
$g_i \star g_j$.

\begin{exemple}[Table de Cayley de $\Z/4\Z$]
\label{ex:cayley-Z4}

\begin{center}
\begin{tikzpicture}[scale=0.85]
  \def\elems{{"0","1","2","3"}}
  \def\table{{"0","1","2","3","1","2","3","0","2","3","0","1","3","0","1","2"}}
  \foreach \i in {0,...,3} {
    \pgfmathparse{\elems[\i]}
    \node[font=\bfseries] at (-0.7, -\i) {$\overline{\pgfmathresult}$};
    \node[font=\bfseries] at (\i, 0.7) {$\overline{\pgfmathresult}$};
  }
  \node[font=\bfseries] at (-0.7, 0.7) {$+$};
  \draw[thick] (-0.3, 0.35) -- (3.5, 0.35);
  \draw[thick] (-0.3, 0.35) -- (-0.3, -3.5);
  \foreach \i in {0,...,3} {
    \foreach \j in {0,...,3} {
      \pgfmathtruncatemacro{\idx}{\i*4+\j}
      \pgfmathparse{\table[\idx]}
      \node at (\j, -\i) {$\overline{\pgfmathresult}$};
    }
  }
\end{tikzpicture}
\end{center}
\end{exemple}

\begin{exemple}[Table de Cayley du groupe de Klein $V_4$]
\label{ex:cayley-V4}

\begin{center}
\begin{tikzpicture}[scale=0.85]
  \def\names{{"e","a","b","c"}}
  \def\table{{"e","a","b","c","a","e","c","b","b","c","e","a","c","b","a","e"}}
  \foreach \i in {0,...,3} {
    \pgfmathparse{\names[\i]}
    \node[font=\bfseries] at (-0.7, -\i) {$\pgfmathresult$};
    \node[font=\bfseries] at (\i, 0.7) {$\pgfmathresult$};
  }
  \node[font=\bfseries] at (-0.7, 0.7) {$\cdot$};
  \draw[thick] (-0.3, 0.35) -- (3.5, 0.35);
  \draw[thick] (-0.3, 0.35) -- (-0.3, -3.5);
  \foreach \i in {0,...,3} {
    \foreach \j in {0,...,3} {
      \pgfmathtruncatemacro{\idx}{\i*4+\j}
      \pgfmathparse{\table[\idx]}
      \node at (\j, -\i) {$\pgfmathresult$};
    }
  }
\end{tikzpicture}
\end{center}
\end{exemple}

\section{Symétries des polygones}
\label{sec:ch1-symmetries-polygones}

\subsection{Symétries du triangle --- le groupe $D_3$}
\label{subsec:D3}

\begin{center}
\begin{tikzpicture}[scale=2]
  % Triangle
  \coordinate (A) at (90:1);
  \coordinate (B) at (210:1);
  \coordinate (C) at (330:1);
  \draw[thick] (A) -- (B) -- (C) -- cycle;
  \node[above] at (A) {$1$};
  \node[below left] at (B) {$2$};
  \node[below right] at (C) {$3$};

  % Axes de symétrie
  \draw[dashed, blue!60, thick] (A) -- ($(B)!0.5!(C)$);
  \draw[dashed, red!60, thick] (B) -- ($(A)!0.5!(C)$);
  \draw[dashed, green!60!black, thick] (C) -- ($(A)!0.5!(B)$);

  % Centre
  \fill (0,0) circle (1pt);

  % Légende
  \node[right] at (1.3, 0.5) {$r$~: rotation de $120°$};
  \node[right] at (1.3, 0) {$s$~: réflexion};
  \node[right] at (1.3, -0.5) {$D_3 = \{e, r, r^2, s, sr, sr^2\}$};
  \node[right] at (1.3, -0.9) {$\abs{D_3} = 6$};
\end{tikzpicture}
\end{center}

Les éléments de $D_3$ sont~:
\begin{itemize}
  \item $e = \Id$~: l'identité,
  \item $r$~: rotation de $2\pi/3$ (sens trigonométrique),
  \item $r^2$~: rotation de $4\pi/3$,
  \item $s$~: réflexion par rapport à l'axe passant par le sommet~$1$,
  \item $sr$~: réflexion par rapport à l'axe passant par le sommet~$3$,
  \item $sr^2$~: réflexion par rapport à l'axe passant par le sommet~$2$.
\end{itemize}
Les relations fondamentales sont $r^3 = e$, $s^2 = e$ et $srs = r^{-1}$ (c'est-à-dire
$sr = r^{-1}s = r^2 s$). On remarque que $D_3 \iso S_3$ (puisque le groupe a $6$
éléments et opère fidèlement sur les trois sommets).

\subsection{Symétries du carré --- le groupe $D_4$}
\label{subsec:D4}

\begin{center}
\begin{tikzpicture}[scale=1.8]
  % Carré
  \coordinate (A) at (1,1);
  \coordinate (B) at (-1,1);
  \coordinate (C) at (-1,-1);
  \coordinate (D) at (1,-1);
  \draw[thick] (A) -- (B) -- (C) -- (D) -- cycle;
  \node[above right] at (A) {$1$};
  \node[above left] at (B) {$2$};
  \node[below left] at (C) {$3$};
  \node[below right] at (D) {$4$};

  % Axes de symétrie
  \draw[dashed, blue!60, thick] (0,1.3) -- (0,-1.3);
  \draw[dashed, blue!60, thick] (-1.3,0) -- (1.3,0);
  \draw[dashed, red!60, thick] (1.2,1.2) -- (-1.2,-1.2);
  \draw[dashed, red!60, thick] (-1.2,1.2) -- (1.2,-1.2);

  % Centre
  \fill (0,0) circle (1pt);

  % Légende
  \node[right] at (1.8, 0.8) {$r$~: rotation de $90°$};
  \node[right] at (1.8, 0.3) {$4$ rotations~: $e, r, r^2, r^3$};
  \node[right] at (1.8, -0.2) {$4$ réflexions~: $s, sr, sr^2, sr^3$};
  \node[right] at (1.8, -0.7) {$\abs{D_4} = 8$};
\end{tikzpicture}
\end{center}

Le groupe $D_4 = \gen{r, s \mid r^4 = s^2 = e,\; srs = r^{-1}}$ est d'ordre~$8$. Il
est non abélien et possède une structure remarquablement riche pour un groupe si petit~:
il a $10$ sous-groupes, dont $3$ sont distingués.

\section{Diagrammes commutatifs}
\label{sec:ch1-diagrammes}

Les \emph{diagrammes commutatifs}\index{diagramme commutatif} constituent un outil
fondamental en algèbre. Un diagramme d'ensembles (ou de groupes) et d'applications (ou
d'homomorphismes) est dit \emph{commutatif} si, pour tout couple de chemins reliant deux
mêmes sommets, les applications composées correspondantes sont égales.

\begin{exemple}
Le diagramme suivant exprime que $h = g \circ f$~:
\[
\begin{tikzcd}
  A \arrow[r, "f"] \arrow[dr, "h"'] & B \arrow[d, "g"] \\
  & C
\end{tikzcd}
\]
\end{exemple}

\begin{exemple}
Le diagramme commutatif suivant résume la situation de la classification des groupes
cycliques. Si $G = \gen{g}$ est cyclique d'ordre $n$, alors~:
\[
\begin{tikzcd}
  \Z \arrow[r, "\varphi", two heads] \arrow[d, "\pi"'] & G \\
  \Z/n\Z \arrow[ur, "\bar{\varphi}"', "\sim"]
\end{tikzcd}
\]
où $\varphi(k) = g^k$, $\pi$ est la projection canonique, et $\bar{\varphi}$ est
l'isomorphisme induit.
\end{exemple}

\section{Produit direct de groupes}
\label{sec:ch1-produit-direct}

\begin{definition}[Produit direct]
\label{def:produit-direct}
\index{produit direct}
Soient $(G, \cdot_G)$ et $(H, \cdot_H)$ deux groupes. Le \emph{produit direct} $G
\times H$ est l'ensemble $\{(g, h) \mid g \in G, \; h \in H\}$ muni de la loi~:
\[
  (g_1, h_1) \cdot (g_2, h_2) = (g_1 \cdot_G g_2, \; h_1 \cdot_H h_2).
\]
\end{definition}

\begin{proposition}
\label{prop:produit-direct-groupe}
Le produit direct $G \times H$ est un groupe, d'élément neutre $(e_G, e_H)$ et
d'inverses $(g, h)^{-1} = (g^{-1}, h^{-1})$. Si $G$ et $H$ sont finis, alors $\abs{G
\times H} = \abs{G} \cdot \abs{H}$.
\end{proposition}

\begin{proof}
L'associativité découle composante par composante de celle de $G$ et de $H$. On vérifie
immédiatement que $(e_G, e_H)$ est neutre et que $(g^{-1}, h^{-1})$ est l'inverse de
$(g, h)$. Enfin, $\abs{G \times H} = \abs{G} \cdot \abs{H}$ résulte du principe de
dénombrement des paires.
\end{proof}

\begin{remarque}
La construction se généralise à un nombre fini quelconque de facteurs~: $G_1 \times G_2
\times \cdots \times G_k$.
\end{remarque}

\begin{exemple}
\label{ex:produit-direct-exemples}
\begin{enumerate}[label=(\roman*)]
  \item $\Z/2\Z \times \Z/2\Z \iso V_4$ (groupe de Klein).
  \item $\Z/2\Z \times \Z/3\Z \iso \Z/6\Z$ (car $\pgcd(2,3)=1$).
  \item $\Z/2\Z \times \Z/2\Z \not\iso \Z/4\Z$ (le premier n'a aucun élément d'ordre
    $4$).
\end{enumerate}
\end{exemple}

\section{Exercices}
\label{sec:ch1-exercices}

\begin{exercice}[Sous-groupes de $\Z$]
\label{exo:sg-Z}
Montrer que tout sous-groupe de $(\Z, +)$ est de la forme $n\Z$ pour un certain $n \in
\N$. [\emph{Indication~: considérer le plus petit élément strictement positif, s'il
existe.}]
\end{exercice}

\begin{exercice}[Ordre dans $\Z/n\Z$]
\label{exo:ordre-ZnZ}
Soit $n \geq 1$ et $\bar{k} \in \Z/n\Z$. Montrer que $\ord(\bar{k}) = n / \pgcd(k,n)$.
En déduire que $\bar{k}$ engendre $\Z/n\Z$ si et seulement si $\pgcd(k, n) = 1$.
\end{exercice}

\begin{exercice}[Groupe d'ordre $2$]
\label{exo:ordre-2-abelien}
Soit $G$ un groupe tel que $g^2 = e$ pour tout $g \in G$. Montrer que $G$ est abélien.
\end{exercice}

\begin{exercice}[Centre de $\GL_n(\R)$]
\label{exo:centre-GLn}
Montrer que $Z(\GL_n(\R)) = \{\lambda I_n \mid \lambda \in \R^*\}$, l'ensemble des
matrices scalaires inversibles.
\end{exercice}

\begin{exercice}[Le groupe $S_3$]
\label{exo:S3-detaille}
Écrire les $6$ éléments de $S_3$ en notation cyclique. Calculer la table de Cayley de
$S_3$. Identifier tous les sous-groupes de $S_3$ et dessiner le diagramme de Hasse
correspondant.
\end{exercice}

\begin{exercice}[Groupe des quaternions]
\label{exo:Q8-table}
Écrire la table de Cayley complète de $Q_8$. Déterminer le centre $Z(Q_8)$ et l'ordre de
chaque élément.
\end{exercice}

\begin{exercice}[Produit direct et cyclicité]
\label{exo:produit-cyclique}
Soient $m, n \geq 1$. Montrer que $\Z/m\Z \times \Z/n\Z$ est cyclique si et seulement
si $\pgcd(m, n) = 1$. [\emph{Indication~: considérer l'ordre de $(\bar{1}, \bar{1})$.}]
\end{exercice}

\begin{exercice}[Groupe diédral $D_n$]
\label{exo:Dn-proprietes}
Soit $n \geq 3$.
\begin{enumerate}[label=(\alph*)]
  \item Montrer que $D_n = \{e, r, r^2, \ldots, r^{n-1}, s, sr, sr^2, \ldots,
    sr^{n-1}\}$ et que $\abs{D_n} = 2n$.
  \item Montrer que $r^k s = s r^{-k}$ pour tout $k \in \Z$.
  \item Déterminer le centre $Z(D_n)$.
  \item Montrer que $D_n$ est isomorphe à un sous-groupe de $S_n$.
\end{enumerate}
\end{exercice}

\begin{exercice}[Sous-groupe engendré]
\label{exo:sg-engendre}
Soit $G$ un groupe et $S \subseteq G$ une partie non vide. Montrer que le sous-groupe
engendré par $S$ est~:
\[
  \gen{S} = \bigcap_{\substack{H \leq G \\ S \subseteq H}} H
  = \{ s_1^{\varepsilon_1} s_2^{\varepsilon_2} \cdots s_k^{\varepsilon_k} \mid k \geq
  0, \; s_i \in S, \; \varepsilon_i \in \{-1, 1\} \}.
\]
\end{exercice}

\begin{exercice}[Automorphismes de $\Z/n\Z$]
\label{exo:aut-ZnZ}
Montrer que $\Aut(\Z/n\Z) \iso (\Z/n\Z)^\times$. En déduire que $\abs{\Aut(\Z/n\Z)} =
\varphi(n)$, où $\varphi$ désigne la fonction indicatrice d'Euler.
\end{exercice}

\begin{exercice}[Application~: petit théorème de Fermat \emph{(difficile)}]
\label{exo:fermat}
Soit $p$ un nombre premier et $a \in \Z$ tel que $p \nmid a$.
\begin{enumerate}[label=(\alph*)]
  \item Montrer que $(\Z/p\Z)^\times$ est un groupe d'ordre $p - 1$.
  \item En utilisant le théorème de Lagrange (chapitre~2), montrer que $a^{p-1} \equiv 1
    \pmod{p}$.
  \item En déduire que $a^p \equiv a \pmod{p}$ pour tout $a \in \Z$.
\end{enumerate}
\end{exercice}

\begin{exercice}[Commutateur]
\label{exo:commutateur-intro}
Soient $G$ un groupe et $a, b \in G$. Le \emph{commutateur}\index{commutateur} de $a$
et $b$ est $[a,b] = aba^{-1}b^{-1}$.
\begin{enumerate}[label=(\alph*)]
  \item Montrer que $ab = ba[a,b]^{-1}$... Non~: montrer que $ab = ba \cdot [a,b]$
    si et seulement si $[a,b] = e$, c'est-à-dire si et seulement si $a$ et $b$ commutent.
  \item Montrer que $[a,b]^{-1} = [b,a]$.
  \item Calculer quelques commutateurs dans $S_3$ et dans $D_4$.
\end{enumerate}
\end{exercice}

\bigskip
\begin{center}
\fbox{\parbox{0.85\textwidth}{%
\textbf{Résumé du chapitre~1.}
Un \emph{groupe} $(G, \star)$ est un ensemble muni d'une loi associative possédant un
élément neutre et des inverses. Les exemples fondamentaux comprennent $\Z$, $\Z/n\Z$,
$\R^*$, $\GL_n$, $S_n$, $D_n$, $Q_8$, $V_4$. Tout groupe cyclique est isomorphe à $\Z$
ou à $\Z/n\Z$. L'ordre d'un élément $g$ est le cardinal du sous-groupe $\gen{g}$. Le
produit direct $G \times H$ est un groupe d'ordre $\abs{G}\cdot\abs{H}$.
}}
\end{center}


%%%%%%%%%%%%%%%%%%%%%%%%%%%%%%%%%%%%%%%%%%%%%%%%%%%%%%%%%%%%%%%%%%%%
%%%%%%%%%%%%%%%%%%%%%%%%%%%%%%%%%%%%%%%%%%%%%%%%%%%%%%%%%%%%%%%%%%%%
%%
%%  CHAPITRE 2 : SOUS-GROUPES, SOUS-GROUPES DISTINGUÉS, QUOTIENTS
%%
%%%%%%%%%%%%%%%%%%%%%%%%%%%%%%%%%%%%%%%%%%%%%%%%%%%%%%%%%%%%%%%%%%%%
%%%%%%%%%%%%%%%%%%%%%%%%%%%%%%%%%%%%%%%%%%%%%%%%%%%%%%%%%%%%%%%%%%%%

\chapter{Sous-groupes, sous-groupes distingués et quotients}
\label{chap:sous-groupes}

\section{Introduction}
\label{sec:ch2-intro}

L'étude d'un groupe $G$ passe de manière essentielle par celle de ses
\emph{sous-groupes}\index{sous-groupe}. Tout comme un espace vectoriel est analysé à
travers ses sous-espaces, un groupe est compris à travers ses sous-groupes et les
relations qu'ils entretiennent entre eux. Le théorème central de ce chapitre --- le
théorème de Lagrange --- relie l'ordre d'un sous-groupe à l'ordre du groupe ambiant et
constitue l'un des résultats les plus fondamentaux de toute la théorie des groupes.

Le passage au quotient, qui nécessite la notion de sous-groupe \emph{distingué},
fournit un mécanisme de construction de nouveaux groupes à partir d'anciens, et se
révèle indispensable pour la compréhension structurelle des groupes.

\section{Sous-groupes}
\label{sec:ch2-sous-groupes}

\begin{definition}[Sous-groupe]
\label{def:sous-groupe}
\index{sous-groupe!définition}
Soit $(G, \cdot)$ un groupe. Un sous-ensemble $H \subseteq G$ est un \emph{sous-groupe}
de $G$ si $(H, \cdot)$ est lui-même un groupe (pour la loi induite). On note $H \leq G$.
\end{definition}

\begin{theoreme}[Critère de sous-groupe --- forme condensée]
\label{thm:critere-sg}
\index{sous-groupe!critère}
Soit $G$ un groupe et $H \subseteq G$ une partie non vide. Les conditions suivantes sont
équivalentes~:
\begin{enumerate}[label=(\roman*)]
  \item $H$ est un sous-groupe de $G$.
  \item Pour tous $a, b \in H$, on a $ab^{-1} \in H$.
  \item $H$ est non vide, et pour tous $a, b \in H$, on a $ab \in H$ et $a^{-1} \in H$.
\end{enumerate}
\end{theoreme}

\begin{proof}
\textbf{(i) $\Rightarrow$ (ii).} Si $H \leq G$ et $a, b \in H$, alors $b^{-1} \in H$
(car $H$ est un groupe) et $ab^{-1} \in H$ (par stabilité).

\textbf{(ii) $\Rightarrow$ (iii).} Supposons (ii). Puisque $H \neq \varnothing$,
prenons un élément $a \in H$. Alors $e = aa^{-1} \in H$ (en prenant $b = a$). Pour tout
$b \in H$, on a $b^{-1} = eb^{-1} \in H$ (en prenant $a = e$). Enfin, pour $a, b \in
H$, comme $b^{-1} \in H$, on obtient $ab = a(b^{-1})^{-1} \in H$ (en appliquant (ii)
avec $a$ et $b^{-1}$).

\textbf{(iii) $\Rightarrow$ (i).} Si (iii) est vérifiée, alors $H$ est stable par la
loi de $G$ et par passage à l'inverse. L'associativité est héritée de $G$. L'existence
de l'élément neutre~: pour $a \in H$, on a $a^{-1} \in H$ puis $e = a \cdot a^{-1} \in
H$. Donc $(H, \cdot)$ est un groupe.
\end{proof}

\begin{remarque}[Sous-groupes triviaux]
\label{rem:sg-triviaux}
Tout groupe $G$ possède au moins deux sous-groupes~: le sous-groupe \emph{trivial}
$\{e\}$ et $G$ lui-même. Un sous-groupe $H$ tel que $\{e\} \subsetneq H \subsetneq G$
est dit \emph{propre}\index{sous-groupe!propre}.
\end{remarque}

\begin{corollaire}[Critère de sous-groupe fini]
\label{cor:critere-sg-fini}
Si $G$ est un groupe et $H \subseteq G$ est une partie \textbf{finie} et non vide,
stable par la loi de $G$ (c'est-à-dire $ab \in H$ pour tous $a, b \in H$), alors $H$
est un sous-groupe de $G$.
\end{corollaire}

\begin{proof}
Il suffit de vérifier que $H$ contient les inverses. Soit $a \in H$. Si $a = e$, alors
$a^{-1} = e \in H$. Sinon, considérons les éléments $a, a^2, a^3, \ldots$ Par stabilité,
tous ces éléments appartiennent à $H$. Puisque $H$ est fini, il existe $i < j$ tels que
$a^i = a^j$, d'où $a^{j-i} = e$ par régularité. Posons $n = j - i \geq 1$. Alors
$a^{-1} = a^{n-1} \in H$ (car $n - 1 \geq 0$ et $a^0 = e \in H$, $a^k \in H$ par
stabilité).
\end{proof}

\section{Exemples de sous-groupes}
\label{sec:ch2-exemples-sg}

\begin{exemple}[Sous-groupes de $(\Z, +)$]
\label{ex:sg-Z}
\index{Z@$\Z$!sous-groupes de}
Les sous-groupes de $(\Z, +)$ sont exactement les ensembles $n\Z = \{nk \mid k \in \Z\}$
pour $n \in \N$ (cf.\ exercice~\ref{exo:sg-Z}). En particulier, tout sous-groupe de
$\Z$ est cyclique.
\end{exemple}

\begin{exemple}[$\SL_n$ dans $\GL_n$]
\label{ex:SLn-dans-GLn}
Le groupe spécial linéaire $\SL_n(K)$ est un sous-groupe de $\GL_n(K)$~: il est non
vide ($I_n \in \SL_n(K)$), et pour $A, B \in \SL_n(K)$, on a $\det(AB^{-1}) =
\det(A)/\det(B) = 1/1 = 1$, donc $AB^{-1} \in \SL_n(K)$.
\end{exemple}

\begin{exemple}[$A_n$ dans $S_n$]
\label{ex:An-dans-Sn}
\index{A@$A_n$}\index{groupe!alterné}
Le \emph{groupe alterné} $A_n$ est l'ensemble des permutations paires de $S_n$,
c'est-à-dire les permutations pouvant s'écrire comme produit d'un nombre pair de
transpositions. C'est un sous-groupe de $S_n$ d'ordre $n!/2$ (pour $n \geq 2$).
\end{exemple}

\begin{exemple}[Intersection de sous-groupes]
\label{ex:intersection-sg}
Si $(H_i)_{i \in I}$ est une famille de sous-groupes de $G$, alors $\bigcap_{i \in I}
H_i$ est un sous-groupe de $G$. En effet, l'intersection est non vide (elle contient
$e$) et, pour $a, b \in \bigcap H_i$, on a $ab^{-1} \in H_i$ pour tout $i$ (car chaque
$H_i$ est un sous-groupe), donc $ab^{-1} \in \bigcap H_i$.
\end{exemple}

\begin{remarque}
La \emph{réunion} de deux sous-groupes n'est en général \emph{pas} un sous-groupe. Par
exemple, $2\Z \cup 3\Z$ n'est pas un sous-groupe de $\Z$ (car $2 + 3 = 5 \notin 2\Z
\cup 3\Z$).
\end{remarque}

\section{Théorème de Lagrange}
\label{sec:ch2-lagrange}

Le théorème de Lagrange\index{Lagrange!théorème de} est l'un des résultats les plus
importants de la théorie des groupes finis. Sa démonstration repose sur la notion de
\emph{classe latérale}.

\begin{definition}[Classes latérales]
\label{def:classes-laterales}
\index{classe latérale}\index{coset|see{classe latérale}}
Soient $G$ un groupe et $H \leq G$. Pour $g \in G$, on définit~:
\begin{itemize}
  \item la \emph{classe latérale à gauche} de $g$ modulo $H$~: $gH = \{gh \mid h \in
    H\}$,
  \item la \emph{classe latérale à droite} de $g$ modulo $H$~: $Hg = \{hg \mid h \in
    H\}$.
\end{itemize}
\end{definition}

\begin{lemme}
\label{lem:classes-laterales}
Soient $G$ un groupe et $H \leq G$. Pour tous $a, b \in G$~:
\begin{enumerate}[label=(\roman*)]
  \item $aH = bH$ si et seulement si $b^{-1}a \in H$.
  \item Les classes latérales à gauche forment une partition de $G$.
  \item Toutes les classes latérales à gauche ont le même cardinal~: $\abs{aH} =
    \abs{H}$.
\end{enumerate}
\end{lemme}

\begin{proof}
\textbf{(i)} On a $aH = bH$ si et seulement si $a \in bH$ (car $a = ae \in aH$), ce
qui équivaut à l'existence de $h \in H$ tel que $a = bh$, soit $b^{-1}a = h \in H$.

\textbf{(ii)} Définissons la relation $\sim$ sur $G$ par $a \sim b \Leftrightarrow aH =
bH$, c'est-à-dire $b^{-1}a \in H$ d'après~(i). Montrons que $\sim$ est une relation
d'équivalence~:
\begin{itemize}
  \item \emph{Réflexivité~:} $a^{-1}a = e \in H$.
  \item \emph{Symétrie~:} si $b^{-1}a \in H$, alors $a^{-1}b = (b^{-1}a)^{-1} \in H$.
  \item \emph{Transitivité~:} si $b^{-1}a \in H$ et $c^{-1}b \in H$, alors $c^{-1}a =
    (c^{-1}b)(b^{-1}a) \in H$.
\end{itemize}
Les classes latérales à gauche sont les classes d'équivalence de $\sim$, elles forment
donc une partition de $G$.

\textbf{(iii)} L'application $\psi : H \to aH$ définie par $\psi(h) = ah$ est bijective.
En effet, elle est surjective par définition de $aH$, et injective car $ah_1 = ah_2$
implique $h_1 = h_2$ (régularité à gauche). Donc $\abs{aH} = \abs{H}$.
\end{proof}

\begin{definition}[Indice]
\label{def:indice}
\index{indice d'un sous-groupe}
L'\emph{indice} de $H$ dans $G$, noté $\idx{G}{H}$, est le nombre de classes latérales
à gauche de $H$ dans $G$~:
\[
  \idx{G}{H} = \abs{G/H} \quad (\text{ici, $G/H$ désigne l'ensemble des classes $gH$}).
\]
\end{definition}

\begin{theoreme}[Théorème de Lagrange]
\label{thm:lagrange}
\index{Lagrange!théorème de}
Soit $G$ un groupe fini et $H \leq G$. Alors~:
\[
  \abs{G} = \idx{G}{H} \cdot \abs{H}.
\]
En particulier, $\abs{H}$ divise $\abs{G}$.
\end{theoreme}

\begin{proof}
Par le lemme~\ref{lem:classes-laterales}, les classes latérales à gauche $g_1 H, g_2 H,
\ldots, g_k H$ (où $k = \idx{G}{H}$) forment une partition de $G$ en $k$ parties
disjointes, chacune de cardinal $\abs{H}$. Par conséquent~:
\[
  \abs{G} = \sum_{i=1}^{k} \abs{g_i H} = \sum_{i=1}^{k} \abs{H} = k \cdot \abs{H}
  = \idx{G}{H} \cdot \abs{H}. \qedhere
\]
\end{proof}

\begin{corollaire}[L'ordre d'un élément divise l'ordre du groupe]
\label{cor:ordre-divise}
\index{ordre!divise le cardinal}
Soit $G$ un groupe fini et $g \in G$. Alors $\ord(g)$ divise $\abs{G}$. En particulier,
$g^{\abs{G}} = e$.
\end{corollaire}

\begin{proof}
On a $\ord(g) = \abs{\gen{g}}$ (proposition~\ref{prop:ordre-gen}) et $\gen{g} \leq G$.
Par le théorème de Lagrange, $\abs{\gen{g}}$ divise $\abs{G}$. Donc $\ord(g)$ divise
$\abs{G}$. De plus, $\abs{G} = \ord(g) \cdot q$ pour un certain $q$, d'où $g^{\abs{G}}
= (g^{\ord(g)})^q = e^q = e$.
\end{proof}

\begin{corollaire}[Groupes d'ordre premier]
\label{cor:ordre-premier-cyclique}
\index{groupe!d'ordre premier}
Si $G$ est un groupe d'ordre premier $p$, alors $G$ est cyclique, isomorphe à
$\Z/p\Z$.
\end{corollaire}

\begin{proof}
Puisque $\abs{G} = p > 1$, il existe $g \in G$ avec $g \neq e$. Alors $\ord(g) > 1$ et
$\ord(g)$ divise $p$ (corollaire~\ref{cor:ordre-divise}). Comme $p$ est premier, $\ord(g)
= p = \abs{G}$. Donc $\gen{g} = G$~: le groupe $G$ est cyclique d'ordre $p$, isomorphe
à $\Z/p\Z$ par le théorème~\ref{thm:classification-cycliques}.
\end{proof}

\begin{remarque}[Réciproque fausse]
\label{rem:lagrange-reciproque}
La réciproque du théorème de Lagrange est fausse~: si $d$ divise $\abs{G}$, il n'existe
pas nécessairement de sous-groupe d'ordre $d$. Par exemple, $A_4$ (d'ordre~$12$) ne
possède pas de sous-groupe d'ordre~$6$.
\end{remarque}

\section{Classes latérales~: visualisation}
\label{sec:ch2-visualisation-cosets}

La partition de $G$ en classes latérales de $H$ se visualise comme suit~:

\begin{center}
\begin{tikzpicture}[scale=0.9]
  % Grand rectangle pour G
  \draw[thick, rounded corners=5pt] (0,0) rectangle (12,4);
  \node[above] at (6,4) {$G$};

  % Classes latérales
  \draw[fill=blue!10, rounded corners=3pt] (0.3,0.3) rectangle (2.7,3.7);
  \node at (1.5, 2) {$eH = H$};

  \draw[fill=red!10, rounded corners=3pt] (3.1,0.3) rectangle (5.5,3.7);
  \node at (4.3, 2) {$g_1 H$};

  \draw[fill=green!10, rounded corners=3pt] (5.9,0.3) rectangle (8.3,3.7);
  \node at (7.1, 2) {$g_2 H$};

  \draw[fill=yellow!15, rounded corners=3pt] (8.7,0.3) rectangle (11.7,3.7);
  \node at (10.2, 2) {$\cdots$};
  \node at (10.2, 1.2) {$g_{k-1}H$};

  % Accolade en bas
  \draw[decorate, decoration={brace, amplitude=8pt, mirror}] (0.3,-0.1) -- (2.7,-0.1)
    node[midway, below=10pt] {$\abs{H}$ éléments};

  % Accolade en haut
  \draw[decorate, decoration={brace, amplitude=8pt}] (0.3,4.5) -- (11.7,4.5)
    node[midway, above=10pt] {$k = \idx{G}{H}$ classes};
\end{tikzpicture}
\end{center}

On a donc $\abs{G} = k \cdot \abs{H}$, la formule du théorème de Lagrange.

\section{Sous-groupes distingués}
\label{sec:ch2-distingues}

\begin{definition}[Sous-groupe distingué]
\label{def:distingue}
\index{sous-groupe!distingué}\index{sous-groupe!normal|see{distingué}}
Un sous-groupe $N$ de $G$ est dit \emph{distingué} (ou \emph{normal}, ou
\emph{invariant}) dans $G$, et on note $N \normal G$, si pour tout $g \in G$~:
\[
  gNg^{-1} = N,
\]
c'est-à-dire $gNg^{-1} \subseteq N$ (ce qui suffit, voir ci-dessous).
\end{definition}

\begin{theoreme}[Caractérisations des sous-groupes distingués]
\label{thm:caract-distingue}
\index{sous-groupe!distingué!caractérisations}
Soit $G$ un groupe et $N \leq G$. Les conditions suivantes sont équivalentes~:
\begin{enumerate}[label=(\roman*)]
  \item $N \normal G$.
  \item Pour tout $g \in G$, $gNg^{-1} \subseteq N$.
  \item Pour tout $g \in G$, $gN = Ng$ \textup{(les classes à gauche et à droite
    coïncident)}.
  \item Pour tout $g \in G$ et tout $n \in N$, on a $gng^{-1} \in N$.
\end{enumerate}
\end{theoreme}

\begin{proof}
\textbf{(i) $\Leftrightarrow$ (ii).} L'implication (i) $\Rightarrow$ (ii) est triviale.
Pour la réciproque, supposons $gNg^{-1} \subseteq N$ pour tout $g$. En appliquant cela
à $g^{-1}$, on obtient $g^{-1}Ng \subseteq N$, soit $N \subseteq gNg^{-1}$.
D'où $gNg^{-1} = N$.

\textbf{(i) $\Leftrightarrow$ (iii).} On a $gN = Ng$ si et seulement si
$gNg^{-1} = N$ (en multipliant à droite par $g^{-1}$, respectivement par $g$).

\textbf{(ii) $\Leftrightarrow$ (iv).} La condition $gNg^{-1} \subseteq N$ signifie
exactement que pour tout $n \in N$, $gng^{-1} \in N$.
\end{proof}

\begin{remarque}
\label{rem:abelien-distingue}
Si $G$ est abélien, alors \emph{tout} sous-groupe de $G$ est distingué, car $gng^{-1} =
n \in N$ pour tout $g \in G$ et $n \in N$.
\end{remarque}

\begin{exemple}[Sous-groupes distingués]
\label{ex:sg-distingues}
\begin{enumerate}[label=(\roman*)]
  \item $\{e\}$ et $G$ sont toujours distingués dans $G$.
  \item $n\Z \normal \Z$ (puisque $\Z$ est abélien).
  \item $\SL_n(K) \normal \GL_n(K)$, car pour $A \in \GL_n(K)$ et $M \in \SL_n(K)$,
    $\det(AMA^{-1}) = \det(M) = 1$.
  \item $A_n \normal S_n$, car $A_n$ est le noyau du morphisme signature $\varepsilon :
    S_n \to \{-1, +1\}$.
  \item Dans $S_3$, le sous-groupe $\gen{(1\;2\;3)} = \{e, (1\;2\;3), (1\;3\;2)\}$
    (d'ordre~$3$) est distingué, mais les sous-groupes d'ordre~$2$ (engendrés par les
    transpositions) ne le sont pas.
\end{enumerate}
\end{exemple}

\section{Groupe quotient}
\label{sec:ch2-quotient}

\begin{definition}[Ensemble quotient]
\label{def:ensemble-quotient}
Soit $G$ un groupe et $N \normal G$. L'\emph{ensemble quotient} $G/N$ est l'ensemble des
classes latérales à gauche~:
\[
  G/N = \{ gN \mid g \in G \}.
\]
\end{definition}

\begin{theoreme}[Groupe quotient]
\label{thm:groupe-quotient}
\index{groupe!quotient}\index{quotient}
Soit $G$ un groupe et $N \normal G$. La loi~:
\[
  (aN) \cdot (bN) = (ab)N
\]
est bien définie sur $G/N$ et munit $G/N$ d'une structure de groupe. L'élément neutre
est $eN = N$ et l'inverse de $aN$ est $a^{-1}N$. L'ordre de $G/N$ est $\idx{G}{N}$.
\end{theoreme}

\begin{proof}
\textbf{Bonne définition.} C'est le point crucial. Supposons $aN = a'N$ et $bN = b'N$~;
il faut montrer que $(ab)N = (a'b')N$. Par hypothèse, il existe $n_1, n_2 \in N$ tels
que $a' = an_1$ et $b' = bn_2$. Alors~:
\[
  a'b' = an_1 \cdot bn_2 = a \cdot (n_1 b) \cdot n_2.
\]
Or $n_1 b = b(b^{-1}n_1 b)$. Puisque $N$ est distingué dans $G$, on a $b^{-1}n_1 b \in
N$~; notons $n_3 = b^{-1}n_1 b \in N$. Donc~:
\[
  a'b' = ab \cdot n_3 n_2,
\]
et $n_3 n_2 \in N$ car $N$ est un sous-groupe. Ainsi $a'b' \in (ab)N$, c'est-à-dire
$(a'b')N = (ab)N$.

\textbf{Associativité.} Pour $a, b, c \in G$~:
\[
  ((aN)(bN))(cN) = (ab)N \cdot cN = ((ab)c)N = (a(bc))N = aN \cdot (bc)N = (aN)((bN)(cN)).
\]

\textbf{Élément neutre.} $N \cdot aN = (ea)N = aN$ et $aN \cdot N = (ae)N = aN$, donc
$N = eN$ est neutre.

\textbf{Inverses.} $(aN)(a^{-1}N) = (aa^{-1})N = eN = N$ et de même
$(a^{-1}N)(aN) = N$. Donc l'inverse de $aN$ est $a^{-1}N$.
\end{proof}

\begin{remarque}
La condition $N \normal G$ est \textbf{indispensable} pour que la loi sur $G/N$ soit
bien définie. Si $H \leq G$ n'est pas distingué, l'ensemble $G/H$ n'a pas de structure
naturelle de groupe.
\end{remarque}

\begin{definition}[Projection canonique]
\label{def:projection-canonique}
\index{projection canonique}
L'application
\[
  \pi : G \longrightarrow G/N, \qquad g \longmapsto gN
\]
est un homomorphisme surjectif de groupes, appelé \emph{projection canonique} (ou
\emph{surjection canonique}). Son noyau est $\Ker(\pi) = N$.
\end{definition}

Le diagramme commutatif suivant résume la situation~:
\[
\begin{tikzcd}
  G \arrow[r, "\pi", two heads] & G/N
\end{tikzcd}
\]

\begin{exemple}[Quotients classiques]
\label{ex:quotients-classiques}
\begin{enumerate}[label=(\roman*)]
  \item $\Z / n\Z$ est le quotient du groupe $(\Z, +)$ par le sous-groupe $n\Z$. Les
    classes sont $\bar{0}, \bar{1}, \ldots, \overline{n-1}$.
  \item $\GL_n(K) / \SL_n(K) \iso K^*$. En effet, le morphisme $\det : \GL_n(K) \to
    K^*$ est surjectif et a pour noyau $\SL_n(K)$.
  \item $S_n / A_n \iso \Z/2\Z$. Le morphisme signature $\varepsilon : S_n \to \{-1,
    +1\} \iso \Z/2\Z$ est surjectif de noyau $A_n$.
\end{enumerate}
\end{exemple}

\section{Centre d'un groupe}
\label{sec:ch2-centre}

\begin{definition}[Centre]
\label{def:centre}
\index{centre d'un groupe}
Le \emph{centre} d'un groupe $G$, noté $Z(G)$, est l'ensemble des éléments qui commutent
avec tous les éléments de $G$~:
\[
  Z(G) = \{ z \in G \mid \forall g \in G, \; zg = gz \}.
\]
\end{definition}

\begin{proposition}
\label{prop:centre-distingue}
Le centre $Z(G)$ est un sous-groupe distingué abélien de $G$. De plus, $G$ est abélien
si et seulement si $Z(G) = G$.
\end{proposition}

\begin{proof}
\textbf{Sous-groupe~:} $e \in Z(G)$ (donc $Z(G) \neq \varnothing$). Si $a, b \in Z(G)$,
alors pour tout $g \in G$, $(ab^{-1})g = a(b^{-1}g) = a(gb^{-1}) = (ag)b^{-1} =
(ga)b^{-1} = g(ab^{-1})$. (On utilise le fait que $b^{-1}$ commute avec $g$ si $b$ le
fait~: $b^{-1}g = b^{-1}g(bb^{-1}) = b^{-1}(gb)b^{-1} = b^{-1}(bg)b^{-1} =
gb^{-1}$.) Donc $ab^{-1} \in Z(G)$.

\textbf{Distingué~:} pour $g \in G$ et $z \in Z(G)$, on a $gzg^{-1} = zgg^{-1} = z \in
Z(G)$.

\textbf{Abélien~:} si $a, b \in Z(G)$, alors $ab = ba$ par définition.

Enfin, $Z(G) = G$ si et seulement si tout élément commute avec tout autre, soit $G$
abélien.
\end{proof}

\begin{exemple}
\label{ex:centres}
\begin{enumerate}[label=(\roman*)]
  \item $Z(S_n) = \{e\}$ pour $n \geq 3$.
  \item $Z(\GL_n(K)) = \{\lambda I_n \mid \lambda \in K^*\}$ (les matrices scalaires).
  \item $Z(D_n) = \{e\}$ si $n$ est impair, $Z(D_n) = \{e, r^{n/2}\}$ si $n$ est pair.
  \item $Z(Q_8) = \{1, -1\}$.
\end{enumerate}
\end{exemple}

\section{Sous-groupe dérivé}
\label{sec:ch2-derive}

\begin{definition}[Commutateur, sous-groupe dérivé]
\label{def:derive}
\index{commutateur}\index{sous-groupe!dérivé}
Soit $G$ un groupe. Le \emph{commutateur} de $a, b \in G$ est $[a, b] = aba^{-1}b^{-1}$.
Le \emph{sous-groupe dérivé} (ou \emph{groupe des commutateurs}) de $G$ est~:
\[
  [G, G] = \gen{\, [a,b] \mid a, b \in G \,},
\]
le sous-groupe engendré par l'ensemble de tous les commutateurs.
\end{definition}

\begin{remarque}
L'ensemble $\{[a,b] \mid a, b \in G\}$ n'est en général \emph{pas} un sous-groupe~; il
faut prendre le sous-groupe engendré.
\end{remarque}

\begin{proposition}
\label{prop:derive-distingue}
Le sous-groupe dérivé $[G, G]$ est un sous-groupe distingué de $G$, et le quotient
$G/[G,G]$ est abélien. De plus, si $N \normal G$ et $G/N$ est abélien, alors $[G,G]
\subseteq N$.
\end{proposition}

\begin{proof}
\textbf{Distingué~:} pour $g \in G$ et $a, b \in G$, on a~:
\[
  g[a,b]g^{-1} = g(aba^{-1}b^{-1})g^{-1} = (gag^{-1})(gbg^{-1})(ga^{-1}g^{-1})
  (gb^{-1}g^{-1}) = [gag^{-1}, gbg^{-1}].
\]
Donc $g[a,b]g^{-1}$ est un commutateur, et par conséquent $g[G,G]g^{-1} \subseteq [G,G]$.

\textbf{$G/[G,G]$ est abélien~:} soient $a, b \in G$. Dans $G/[G,G]$, on a~:
\[
  (a[G,G])(b[G,G]) = ab[G,G] \quad \text{et} \quad (b[G,G])(a[G,G]) = ba[G,G].
\]
Or $ab(ba)^{-1} = aba^{-1}b^{-1} = [a,b] \in [G,G]$, donc $ab[G,G] = ba[G,G]$.

\textbf{Minimalité~:} si $N \normal G$ et $G/N$ est abélien, alors pour tous $a, b \in
G$, $aN \cdot bN = bN \cdot aN$, soit $(ab)N = (ba)N$, d'où $a^{-1}b^{-1}ab \in N$,
c'est-à-dire $[a^{-1}, b^{-1}] \in N$. Comme $N$ contient tous les commutateurs et est
un sous-groupe, $[G,G] \subseteq N$.
\end{proof}

\section{Groupes simples}
\label{sec:ch2-simples}

\begin{definition}[Groupe simple]
\label{def:simple}
\index{groupe!simple}
Un groupe $G$ est dit \emph{simple} s'il est non trivial ($\abs{G} > 1$) et si ses seuls
sous-groupes distingués sont $\{e\}$ et $G$.
\end{definition}

\begin{exemple}
\label{ex:simples}
\begin{enumerate}[label=(\roman*)]
  \item Tout groupe d'ordre premier $p$ est simple (car il n'a aucun sous-groupe propre,
    par Lagrange).
  \item Le groupe alterné $A_n$ est simple pour $n \geq 5$ (théorème classique).
  \item $A_4$ n'est pas simple~: il contient le sous-groupe distingué $V_4 = \{e,
    (1\;2)(3\;4), (1\;3)(2\;4), (1\;4)(2\;3)\}$.
\end{enumerate}
\end{exemple}

\begin{theoreme}[Simplicité de $A_n$ pour $n \geq 5$]
\label{thm:An-simple}
\index{A@$A_n$!simplicité}
Pour tout $n \geq 5$, le groupe alterné $A_n$ est simple.
\end{theoreme}

La démonstration de ce théorème sera donnée au chapitre~4, après l'étude des actions de
groupes.

\section{Treillis de sous-groupes}
\label{sec:ch2-treillis}

Le \emph{treillis}\index{treillis de sous-groupes} (ou \emph{diagramme de Hasse}) des
sous-groupes d'un groupe $G$ est le diagramme dans lequel chaque sous-groupe est un
sommet, et on relie $H$ à $K$ par un trait si $H \subsetneq K$ et s'il n'existe pas de
sous-groupe strictement intermédiaire.

\begin{exemple}[Treillis des sous-groupes de $\Z/12\Z$]
\label{ex:treillis-Z12}
\begin{center}
\begin{tikzpicture}[scale=1.2, every node/.style={draw, circle, inner sep=2pt,
  minimum size=20pt, font=\small}]
  \node (Z12) at (0,4) {$\Z/12\Z$};
  \node (6Z) at (-2,3) {$\gen{\bar{2}}$};
  \node (4Z) at (0,3) {$\gen{\bar{3}}$};
  \node (3Z) at (2,3) {$\gen{\bar{4}}$};
  \node (2Z) at (-1,2) {$\gen{\bar{6}}$};
  \node (Z3) at (1,2) {$\gen{\bar{6}}$};

  % On corrige pour le vrai treillis
  \node[draw=none] at (0,1.3) {};
\end{tikzpicture}
\end{center}
Précisons~: les sous-groupes de $\Z/12\Z$ sont $\gen{\bar{d}}$ pour $d \mid 12$, soit
$d \in \{1, 2, 3, 4, 6, 12\}$~:
\begin{center}
\begin{tikzpicture}[
  every node/.style={draw, rounded corners=2pt, inner sep=3pt, font=\small},
  level distance=1.2cm,
  scale=1
]
  \node (G) at (0,4) {$\Z/12\Z$};
  \node (2) at (-2,2.7) {$\gen{\bar{2}}$};
  \node (3) at (0,2.7) {$\gen{\bar{3}}$};
  \node (4) at (2,2.7) {$\gen{\bar{4}}$};
  \node (6) at (-1,1.4) {$\gen{\bar{6}}$};
  \node (0) at (0,0) {$\{\bar{0}\}$};

  \draw (G) -- (2);
  \draw (G) -- (3);
  \draw (G) -- (4);
  \draw (2) -- (6);
  \draw (3) -- (6);
  \draw (2) -- (4);
  \draw (4) -- (0);
  \draw (6) -- (0);
\end{tikzpicture}
\end{center}
où $\gen{\bar{2}} = \{\bar{0}, \bar{2}, \bar{4}, \bar{6}, \bar{8}, \bar{10}\}$ (ordre
$6$), $\gen{\bar{3}} = \{\bar{0}, \bar{3}, \bar{6}, \bar{9}\}$ (ordre $4$),
$\gen{\bar{4}} = \{\bar{0}, \bar{4}, \bar{8}\}$ (ordre $3$), $\gen{\bar{6}} =
\{\bar{0}, \bar{6}\}$ (ordre $2$). Tous sont distingués puisque $\Z/12\Z$ est abélien.
\end{exemple}

\begin{exemple}[Treillis des sous-groupes de $S_3$]
\label{ex:treillis-S3}
\begin{center}
\begin{tikzpicture}[
  every node/.style={draw, rounded corners=2pt, inner sep=3pt, font=\small},
  scale=1
]
  \node (G) at (0,3) {$S_3$};
  \node (A3) at (-2.5,1.5) {$A_3 = \gen{(1\;2\;3)}$};
  \node (s1) at (-0.5,1.5) {$\gen{(1\;2)}$};
  \node (s2) at (1.5,1.5) {$\gen{(1\;3)}$};
  \node (s3) at (3.5,1.5) {$\gen{(2\;3)}$};
  \node (e) at (0,0) {$\{e\}$};

  \draw (G) -- (A3);
  \draw (G) -- (s1);
  \draw (G) -- (s2);
  \draw (G) -- (s3);
  \draw (A3) -- (e);
  \draw (s1) -- (e);
  \draw (s2) -- (e);
  \draw (s3) -- (e);
\end{tikzpicture}
\end{center}
Le sous-groupe $A_3$ est le seul sous-groupe distingué propre de $S_3$.
\end{exemple}

\section{Exercices}
\label{sec:ch2-exercices}

\begin{exercice}[Sous-groupes de $\Z/n\Z$]
\label{exo:sg-ZnZ}
Montrer que les sous-groupes de $\Z/n\Z$ sont exactement les $\gen{\bar{d}}$ pour $d$
diviseur de $n$. En déduire que $\Z/n\Z$ a exactement autant de sous-groupes que de
diviseurs de $n$.
\end{exercice}

\begin{exercice}[Sous-groupes de $D_4$]
\label{exo:sg-D4}
Déterminer tous les sous-groupes de $D_4$. Lesquels sont distingués~? Dessiner le treillis
des sous-groupes.
\end{exercice}

\begin{exercice}[Indice $2$ implique distingué]
\label{exo:indice-2}
Montrer que tout sous-groupe d'indice~$2$ est distingué. [\emph{Indication~: il n'y a que
deux classes latérales.}]
\end{exercice}

\begin{exercice}[Centre et quotient]
\label{exo:G-sur-ZG-cyclique}
Soit $G$ un groupe. Montrer que si $G/Z(G)$ est cyclique, alors $G$ est abélien.
[\emph{Indication~: écrire $G/Z(G) = \gen{gZ(G)}$ et montrer que deux éléments
quelconques de $G$ commutent.}]
\end{exercice}

\begin{exercice}[Commutateur dans $S_n$]
\label{exo:commutateur-Sn}
Montrer que $[S_n, S_n] = A_n$ pour $n \geq 2$. [\emph{Indication~: montrer que tout
$3$-cycle est un commutateur.}]
\end{exercice}

\begin{exercice}[Sous-groupes du groupe de Klein]
\label{exo:sg-V4}
Déterminer tous les sous-groupes de $V_4$. Montrer qu'ils sont tous distingués. Calculer
les groupes quotients correspondants.
\end{exercice}

\begin{exercice}[Intersection de sous-groupes distingués]
\label{exo:intersection-distingues}
Soient $N_1, N_2 \normal G$. Montrer que $N_1 \cap N_2 \normal G$.
\end{exercice}

\begin{exercice}[Produit de sous-groupes]
\label{exo:produit-sg}
Soient $H, K \leq G$. On pose $HK = \{hk \mid h \in H, k \in K\}$.
\begin{enumerate}[label=(\alph*)]
  \item Montrer que $HK$ est un sous-groupe de $G$ si et seulement si $HK = KH$.
  \item Montrer que si $H \normal G$ ou $K \normal G$, alors $HK$ est un sous-groupe
    de $G$.
  \item Si $H$ et $K$ sont finis, montrer que $\abs{HK} = \abs{H}\abs{K}/\abs{H \cap
    K}$.
\end{enumerate}
\end{exercice}

\begin{exercice}[Noyau et sous-groupe distingué]
\label{exo:noyau-distingue}
Soit $f : G \to H$ un homomorphisme de groupes. Montrer que $\Ker(f) \normal G$.
\end{exercice}

\begin{exercice}[Quotients de $\Z$]
\label{exo:quotients-Z}
\begin{enumerate}[label=(\alph*)]
  \item Montrer que $\Z/n\Z$ est un groupe cyclique d'ordre $n$ en utilisant la
    définition du groupe quotient.
  \item Déterminer tous les homomorphismes de groupes $\Z/12\Z \to \Z/4\Z$.
\end{enumerate}
\end{exercice}

\begin{exercice}[Correspondance de sous-groupes \emph{(difficile)}]
\label{exo:correspondance}
Soit $N \normal G$ et $\pi : G \to G/N$ la projection canonique. Montrer qu'il existe
une bijection croissante~:
\[
  \{ H \leq G \mid N \subseteq H \} \longleftrightarrow \{ K \leq G/N \}
\]
donnée par $H \mapsto H/N = \pi(H)$ et $K \mapsto \pi^{-1}(K)$. De plus, $H \normal G$
si et seulement si $H/N \normal G/N$.
\end{exercice}

\begin{exercice}[Groupes d'ordre $p^2$ \emph{(difficile)}]
\label{exo:ordre-p2}
Soit $p$ un nombre premier. Montrer que tout groupe d'ordre $p^2$ est abélien.
[\emph{Indication~: montrer d'abord que $Z(G) \neq \{e\}$ en utilisant l'équation des
classes (chapitre~4), puis appliquer l'exercice~\ref{exo:G-sur-ZG-cyclique}.}]
\end{exercice}

\bigskip
\begin{center}
\fbox{\parbox{0.85\textwidth}{%
\textbf{Résumé du chapitre~2.}
Un \emph{sous-groupe} $H \leq G$ satisfait le critère $ab^{-1} \in H$ pour tous $a, b
\in H$. Le \emph{théorème de Lagrange} affirme que $\abs{H}$ divise $\abs{G}$ (en
particulier, $\ord(g)$ divise $\abs{G}$). Un sous-groupe est \emph{distingué} ($N
\normal G$) si $gNg^{-1} = N$ pour tout $g$, ce qui permet de construire le \emph{groupe
quotient} $G/N$ avec la loi $(aN)(bN) = (ab)N$. Le \emph{centre} $Z(G)$ et le
\emph{sous-groupe dérivé} $[G,G]$ sont des sous-groupes distingués fondamentaux. Un
groupe est \emph{simple} s'il n'a pas de sous-groupe distingué propre non trivial~;
$A_n$ est simple pour $n \geq 5$.
}}
\end{center}
% chapters3_4.tex — Algèbre Abstraite I, Chapitres 3–4
% Chapitre 3 : Homomorphismes et Théorèmes d'Isomorphisme
% Chapitre 4 : Actions de Groupes et Théorèmes de Sylow

% ============================================================
%
%  CHAPITRE 3 : HOMOMORPHISMES ET THÉORÈMES D'ISOMORPHISME
%
% ============================================================

\chapter{Homomorphismes et théorèmes d'isomorphisme}\label{ch:homomorphismes}

\section*{Introduction}
\addcontentsline{toc}{section}{Introduction}

Dans les chapitres précédents, nous avons construit un répertoire de
groupes et de sous-groupes, et nous avons appris à fabriquer de nouveaux
groupes par quotient. Mais comment \emph{comparer} deux groupes~? Quand
peut-on affirmer que deux groupes, présentés de manières très
différentes, sont en réalité <<~le même~>> objet algébrique~?

\medskip

La réponse réside dans la notion d'\emph{homomorphisme}\index{homomorphisme}~:
une application entre deux groupes qui \emph{respecte la structure}, c'est-à-dire
qui envoie le produit sur le produit. Les homomorphismes sont les <<~bonnes~>>
applications en théorie des groupes, de la même façon que les applications
linéaires sont les <<~bonnes~>> applications en algèbre linéaire, et les
applications continues en topologie.

\medskip

Ce chapitre culmine avec les \emph{théorèmes d'isomorphisme}, qui
relient noyaux, images, sous-groupes distingués et quotients dans un
cadre à la fois élégant et puissant. Ces théorèmes constituent l'outil
fondamental pour l'étude structurelle des groupes.

% ============================================================
\section{Homomorphismes de groupes}\label{sec:homomorphismes}
% ============================================================

\subsection{Définition et premières propriétés}\label{subsec:hom-def}

\begin{definition}[Homomorphisme de groupes]\label{def:homomorphisme}
  Soient $(G,\cdot)$ et $(H,\ast)$ deux groupes. Un
  \emph{homomorphisme de groupes}\index{homomorphisme!de groupes}
  est une application $f\colon G\to H$ telle que
  \[
    \forall\, a,b\in G,\quad f(a\cdot b) = f(a)\ast f(b).
  \]
\end{definition}

\begin{proposition}[Propriétés élémentaires]\label{prop:hom-proprietes}
  Soit $f\colon G\to H$ un homomorphisme de groupes. Alors~:
  \begin{enumerate}[label=(\roman*)]
    \item $f(e_G) = e_H$\,;
    \item $\forall\, a\in G,\quad f(a^{-1}) = f(a)^{-1}$\,;
    \item $\forall\, a\in G,\;\forall\, n\in\Z,\quad f(a^n) = f(a)^n$.
  \end{enumerate}
\end{proposition}

\begin{proof}
  \textbf{(i)} On a $f(e_G) = f(e_G\cdot e_G) = f(e_G)\ast f(e_G)$.
  En multipliant les deux membres à gauche par $f(e_G)^{-1}$, on
  obtient $e_H = f(e_G)$.

  \medskip

  \textbf{(ii)} Soit $a\in G$. Alors
  \[
    f(a)\ast f(a^{-1}) = f(a\cdot a^{-1}) = f(e_G) = e_H.
  \]
  De même, $f(a^{-1})\ast f(a) = e_H$. Par unicité de l'inverse dans
  $H$, on conclut $f(a^{-1}) = f(a)^{-1}$.

  \medskip

  \textbf{(iii)} Pour $n\geq 1$, on procède par récurrence. Le cas
  $n=0$ découle de (i) et le cas $n<0$ se ramène au cas positif par
  (ii).
\end{proof}

\begin{remarque}\label{rem:hom-composition}
  La composée de deux homomorphismes de groupes est un homomorphisme de
  groupes~: si $f\colon G\to H$ et $g\colon H\to K$ sont des
  homomorphismes, alors $g\circ f\colon G\to K$ est un homomorphisme.
  En effet, pour tous $a,b\in G$,
  \[
    (g\circ f)(ab) = g(f(ab)) = g(f(a)f(b)) = g(f(a))\,g(f(b))
    = (g\circ f)(a)\,(g\circ f)(b).
  \]
\end{remarque}

% ============================================================
\subsection{Noyau et image}\label{subsec:ker-im}
% ============================================================

\begin{definition}[Noyau et image]\label{def:ker-im}
  Soit $f\colon G\to H$ un homomorphisme de groupes. On définit~:
  \begin{enumerate}[label=(\roman*)]
    \item le \emph{noyau}\index{noyau} de $f$~:
      $\Ker f = \{g\in G \mid f(g) = e_H\}$\,;
    \item l'\emph{image}\index{image!d'un homomorphisme} de $f$~:
      $\Im f = \{f(g) \mid g\in G\} = f(G)$.
  \end{enumerate}
\end{definition}

\begin{theoreme}[Le noyau est distingué]\label{thm:ker-distingue}
  Soit $f\colon G\to H$ un homomorphisme de groupes. Alors
  $\Ker f$ est un sous-groupe distingué de $G$, c'est-à-dire
  $\Ker f \trianglelefteq G$.
\end{theoreme}

\begin{proof}
  \emph{Sous-groupe.} L'élément $e_G$ appartient à $\Ker f$ car
  $f(e_G)=e_H$ (proposition~\ref{prop:hom-proprietes}). Si
  $a,b\in\Ker f$, alors
  \[
    f(ab^{-1}) = f(a)\,f(b^{-1}) = f(a)\,f(b)^{-1} = e_H\cdot e_H^{-1} = e_H,
  \]
  donc $ab^{-1}\in\Ker f$. Par le critère du sous-groupe,
  $\Ker f\leq G$.

  \medskip

  \emph{Distingué.} Soit $g\in G$ et $k\in\Ker f$. Alors
  \[
    f(gkg^{-1}) = f(g)\,f(k)\,f(g^{-1}) = f(g)\cdot e_H\cdot f(g)^{-1} = e_H.
  \]
  Donc $gkg^{-1}\in\Ker f$ pour tout $g\in G$, ce qui montre que
  $\Ker f\trianglelefteq G$.
\end{proof}

\begin{theoreme}[L'image est un sous-groupe]\label{thm:im-sous-groupe}
  Soit $f\colon G\to H$ un homomorphisme de groupes. Alors
  $\Im f$ est un sous-groupe de $H$.
\end{theoreme}

\begin{proof}
  L'image est non vide car $e_H = f(e_G)\in\Im f$. Soient
  $h_1,h_2\in\Im f$; il existe $g_1,g_2\in G$ tels que $h_i=f(g_i)$.
  Alors
  \[
    h_1 h_2^{-1} = f(g_1)\,f(g_2)^{-1} = f(g_1)\,f(g_2^{-1}) = f(g_1 g_2^{-1})\in\Im f.
  \]
  Par le critère du sous-groupe, $\Im f\leq H$.
\end{proof}

\begin{proposition}[Injectivité et noyau]\label{prop:hom-injectif}
  Soit $f\colon G\to H$ un homomorphisme. Alors
  \[
    f \text{ est injectif} \iff \Ker f = \{e_G\}.
  \]
\end{proposition}

\begin{proof}
  ($\Rightarrow$) Si $f$ est injectif et $g\in\Ker f$, alors
  $f(g)=e_H=f(e_G)$, d'où $g=e_G$.

  ($\Leftarrow$) Si $\Ker f = \{e_G\}$ et $f(a)=f(b)$, alors
  $f(ab^{-1})=f(a)f(b)^{-1}=e_H$, d'où $ab^{-1}\in\Ker f=\{e_G\}$,
  soit $a=b$.
\end{proof}

% ============================================================
\subsection{Types d'homomorphismes}\label{subsec:types-hom}
% ============================================================

\begin{definition}[Types d'homomorphismes]\label{def:types-hom}
  Soit $f\colon G\to H$ un homomorphisme de groupes.
  \begin{itemize}
    \item $f$ est un \emph{monomorphisme}\index{monomorphisme}
      si $f$ est injectif.
    \item $f$ est un \emph{épimorphisme}\index{epimorphisme@épimorphisme}
      si $f$ est surjectif.
    \item $f$ est un \emph{isomorphisme}\index{isomorphisme}
      si $f$ est bijectif.
      On écrit alors $G\cong H$.
    \item $f$ est un \emph{endomorphisme}\index{endomorphisme}
      si $H=G$ (homomorphisme d'un groupe dans lui-même).
    \item $f$ est un \emph{automorphisme}\index{automorphisme}
      si $f$ est un endomorphisme bijectif.
  \end{itemize}
\end{definition}

\begin{remarque}\label{rem:iso-inverse}
  Si $f\colon G\to H$ est un isomorphisme, alors $f^{-1}\colon H\to G$
  est aussi un isomorphisme. En effet, si $h_1,h_2\in H$, posons
  $g_i=f^{-1}(h_i)$; alors $f(g_1g_2)=h_1h_2$, d'où
  $f^{-1}(h_1h_2)=g_1g_2=f^{-1}(h_1)f^{-1}(h_2)$.
\end{remarque}

% ============================================================
\subsection{Exemples fondamentaux}\label{subsec:hom-exemples}
% ============================================================

\begin{exemple}[Déterminant]\label{ex:hom-det}
  L'application $\det\colon\mathrm{GL}_n(\R)\to\R^*$\index{determinant@déterminant}
  est un homomorphisme de groupes (le groupe
  $\mathrm{GL}_n(\R)$ étant muni de la multiplication matricielle et
  $\R^*$ de la multiplication usuelle). En effet,
  $\det(AB)=\det(A)\det(B)$. On a
  \[
    \Ker(\det) = \{A\in\mathrm{GL}_n(\R)\mid \det(A)=1\} = \mathrm{SL}_n(\R),
  \]
  et $\Im(\det) = \R^*$ (l'application est surjective). Le
  théorème~\ref{thm:ker-distingue} redonne le fait classique que
  $\mathrm{SL}_n(\R)\trianglelefteq\mathrm{GL}_n(\R)$.
\end{exemple}

\begin{exemple}[Signature]\label{ex:hom-sign}
  L'application signature $\varepsilon\colon\Sym_n\to\{+1,-1\}$\index{signature}
  est un homomorphisme de groupes surjectif (pour $n\geq 2$).
  Son noyau est le groupe alterné~:
  $\Ker\varepsilon = \Alt_n \trianglelefteq \Sym_n$.
  Par le premier théorème d'isomorphisme (théorème~\ref{thm:premier-iso}),
  $\Sym_n/\Alt_n\cong\{+1,-1\}\cong\Z/2\Z$.
\end{exemple}

\begin{exemple}[Exponentielle]\label{ex:hom-exp}
  L'application $\exp\colon(\R,+)\to(\R_+^*,\times)$\index{exponentielle}
  est un isomorphisme de groupes. En effet,
  $\exp(x+y)=\exp(x)\cdot\exp(y)$ et $\exp$ est bijective de $\R$ sur
  $\R_+^*$ (son inverse étant le logarithme népérien).
\end{exemple}

\begin{exemple}[Projection canonique]\label{ex:hom-projection}
  Soit $N\trianglelefteq G$. La \emph{projection
  canonique}\index{projection canonique}
  $\pi\colon G\to G/N$ définie par $\pi(g)=gN$ est un épimorphisme de
  groupes. En effet,
  \[
    \pi(ab) = (ab)N = (aN)(bN) = \pi(a)\pi(b),
  \]
  et $\pi$ est clairement surjective. De plus, $\Ker\pi = N$.
\end{exemple}

% ============================================================
\section{Théorèmes d'isomorphisme}\label{sec:thm-iso}
% ============================================================

Les théorèmes d'isomorphisme constituent le c\oe ur de l'algèbre des
groupes. Ils établissent des liens profonds entre noyaux, images,
sous-groupes distingués et quotients.

\subsection{Premier théorème d'isomorphisme}\label{subsec:premier-iso}

\begin{theoreme}[Premier théorème d'isomorphisme]\label{thm:premier-iso}
  Soit $f\colon G\to H$ un homomorphisme de groupes.\index{theoreme
  d'isomorphisme@théorème d'isomorphisme!premier}
  Alors il existe un unique isomorphisme
  $\bar{f}\colon G/\Ker f \xrightarrow{\;\sim\;} \Im f$
  tel que $\bar{f}(g\,\Ker f) = f(g)$ pour tout $g\in G$.
  Autrement dit,
  \[
    G/\Ker f \;\cong\; \Im f.
  \]
\end{theoreme}

Le diagramme commutatif suivant résume la situation~:
\[
\begin{tikzcd}
  G \arrow[r,"f"] \arrow[d,"\pi"',twoheadrightarrow]
    & H \\
  G/\Ker f \arrow[ru,"\bar{f}"',dashed,"\sim"']
\end{tikzcd}
\]
où $\pi\colon G\to G/\Ker f$ désigne la projection canonique.

\begin{proof}
  Posons $N=\Ker f$ et définissons $\bar{f}\colon G/N\to H$ par
  $\bar{f}(gN) = f(g)$.

  \medskip

  \emph{Bonne définition.} Si $gN = g'N$, alors $g'^{-1}g\in N=\Ker f$,
  d'où $f(g'^{-1}g) = e_H$, c'est-à-dire $f(g') = f(g)$.
  Donc $\bar{f}(gN) = \bar{f}(g'N)$.

  \medskip

  \emph{Homomorphisme.} Pour $gN,g'N\in G/N$,
  \[
    \bar{f}(gN\cdot g'N) = \bar{f}((gg')N) = f(gg') = f(g)f(g')
    = \bar{f}(gN)\,\bar{f}(g'N).
  \]

  \emph{Injectivité.} Si $\bar{f}(gN) = e_H$, alors $f(g)=e_H$, d'où
  $g\in\Ker f = N$, soit $gN = N = e_{G/N}$.
  Donc $\Ker\bar{f} = \{N\}$, et $\bar{f}$ est injectif par la
  proposition~\ref{prop:hom-injectif}.

  \medskip

  \emph{Surjectivité sur $\Im f$.} Par construction,
  $\Im\bar{f} = \{f(g)\mid g\in G\} = \Im f$.

  \medskip

  \emph{Unicité.} Si $\tilde{f}\colon G/N\to\Im f$ est un autre
  isomorphisme tel que $\tilde{f}\circ\pi = f$, alors pour tout $g\in G$,
  $\tilde{f}(gN) = \tilde{f}(\pi(g)) = f(g) = \bar{f}(gN)$, donc
  $\tilde{f}=\bar{f}$.
\end{proof}

\begin{corollaire}\label{cor:premier-iso-surjectif}
  Si $f\colon G\to H$ est un épimorphisme (homomorphisme surjectif),
  alors $G/\Ker f\cong H$.
\end{corollaire}

\begin{proof}
  Si $f$ est surjectif, alors $\Im f = H$, et le
  théorème~\ref{thm:premier-iso} donne $G/\Ker f\cong H$.
\end{proof}

\begin{exemple}\label{ex:iso-Zn}
  L'application $f\colon\Z\to\Z/n\Z$ définie par $f(k) = \bar{k}$ est
  un épimorphisme de $(\Z,+)$ sur $(\Z/n\Z,+)$ dont le noyau est
  $n\Z$. Le premier théorème d'isomorphisme redonne
  $\Z/n\Z \cong \Z/n\Z$ (tautologiquement, mais cela illustre le
  mécanisme).
\end{exemple}

% ============================================================
\subsection{Deuxième théorème d'isomorphisme}\label{subsec:deuxieme-iso}
% ============================================================

\begin{theoreme}[Deuxième théorème d'isomorphisme]\label{thm:deuxieme-iso}
  Soient $G$ un groupe, $H\leq G$ un sous-groupe et $N\trianglelefteq G$
  un sous-groupe distingué.\index{theoreme d'isomorphisme@théorème
  d'isomorphisme!deuxième}
  Alors~:
  \begin{enumerate}[label=(\roman*)]
    \item $HN = \{hn\mid h\in H,\; n\in N\}$ est un sous-groupe de $G$\,;
    \item $H\cap N \trianglelefteq H$\,;
    \item $N\trianglelefteq HN$\,;
    \item $HN/N \;\cong\; H/(H\cap N)$.
  \end{enumerate}
\end{theoreme}

Ce théorème est parfois appelé le \emph{théorème du
parallélogramme}\index{theoreme du parallelogramme@théorème du
parallélogramme} ou \emph{du diamant}, en raison du diagramme suivant~:
\[
\begin{tikzcd}[row sep=1.5em, column sep=1.5em]
  & HN \arrow[dl,dash] \arrow[dr,dash] & \\
  H \arrow[dr,dash] & & N \arrow[dl,dash] \\
  & H\cap N &
\end{tikzcd}
\]

\begin{proof}
  \textbf{(i)} Puisque $N\trianglelefteq G$, pour tout $h\in H$ et
  $n\in N$, on a $hn = h\cdot n\cdot h^{-1}\cdot h = n' h$ pour un
  certain $n'\in N$ (car $hNh^{-1}=N$). Plus précisément, $HN$ est
  non vide (il contient $e$), et si $h_1n_1,\, h_2n_2\in HN$, alors
  \[
    (h_1n_1)(h_2n_2)^{-1} = h_1n_1 n_2^{-1}h_2^{-1}
    = h_1 h_2^{-1}\cdot(h_2 n_1 n_2^{-1} h_2^{-1}).
  \]
  Or $h_1h_2^{-1}\in H$ et $h_2 n_1n_2^{-1}h_2^{-1}\in N$ (car
  $N\trianglelefteq G$), donc $(h_1n_1)(h_2n_2)^{-1}\in HN$.
  Par le critère du sous-groupe, $HN\leq G$.

  \medskip

  \textbf{(ii)} Soit $h\in H$ et $x\in H\cap N$. Alors
  $hxh^{-1}\in H$ (car $H$ est un sous-groupe) et $hxh^{-1}\in N$
  (car $N\trianglelefteq G$). Donc $hxh^{-1}\in H\cap N$, ce qui
  montre $H\cap N\trianglelefteq H$.

  \medskip

  \textbf{(iii)} Puisque $N\trianglelefteq G$ et $N\subseteq HN\leq G$,
  on a $N\trianglelefteq HN$.

  \medskip

  \textbf{(iv)} Définissons $\varphi\colon H\to HN/N$ par
  $\varphi(h) = hN$. Vérifions que $\varphi$ est un homomorphisme
  surjectif de noyau $H\cap N$.

  \emph{Homomorphisme.} $\varphi(h_1h_2) = (h_1h_2)N =
  (h_1N)(h_2N) = \varphi(h_1)\varphi(h_2)$.

  \emph{Surjectivité.} Tout élément de $HN/N$ est de la forme
  $(hn)N = hN = \varphi(h)$ pour $h\in H$, $n\in N$.

  \emph{Noyau.} $h\in\Ker\varphi \iff hN = N \iff h\in N$.
  Puisque $h\in H$, on obtient $\Ker\varphi = H\cap N$.

  Par le premier théorème d'isomorphisme
  (théorème~\ref{thm:premier-iso}),
  \[
    H/(H\cap N) = H/\Ker\varphi \;\cong\; \Im\varphi = HN/N. \qedhere
  \]
\end{proof}

% ============================================================
\subsection{Troisième théorème d'isomorphisme}\label{subsec:troisieme-iso}
% ============================================================

\begin{theoreme}[Troisième théorème d'isomorphisme]\label{thm:troisieme-iso}
  Soient $N\trianglelefteq G$ et $H\trianglelefteq G$ avec $N\subseteq H$.
  \index{theoreme d'isomorphisme@théorème d'isomorphisme!troisième}
  Alors $H/N\trianglelefteq G/N$ et
  \[
    (G/N)\big/(H/N) \;\cong\; G/H.
  \]
\end{theoreme}

\[
\begin{tikzcd}
  G \arrow[r,"\pi_H",twoheadrightarrow] \arrow[d,"\pi_N"',twoheadrightarrow]
    & G/H \\
  G/N \arrow[r,"\bar{\pi}"',twoheadrightarrow]
    & (G/N)/(H/N) \arrow[u,"\sim"',dashed]
\end{tikzcd}
\]

\begin{proof}
  \emph{$H/N\trianglelefteq G/N$.} Soit $gN\in G/N$ et $hN\in H/N$.
  Alors
  \[
    (gN)(hN)(gN)^{-1} = (ghg^{-1})N.
  \]
  Puisque $H\trianglelefteq G$, on a $ghg^{-1}\in H$, donc
  $(ghg^{-1})N\in H/N$. Ainsi $H/N\trianglelefteq G/N$.

  \medskip

  \emph{Isomorphisme.} Définissons $\varphi\colon G/N\to G/H$
  par $\varphi(gN) = gH$.

  \emph{Bonne définition.} Si $gN = g'N$, alors $g'^{-1}g\in N\subseteq H$,
  d'où $gH = g'H$.

  \emph{Homomorphisme.}
  $\varphi(gN\cdot g'N) = \varphi((gg')N) = (gg')H = (gH)(g'H)
  = \varphi(gN)\,\varphi(g'N)$.

  \emph{Surjectivité.} Pour tout $gH\in G/H$, on a $\varphi(gN)=gH$.

  \emph{Noyau.}
  $gN\in\Ker\varphi \iff gH = H \iff g\in H \iff gN\in H/N$.
  Donc $\Ker\varphi = H/N$.

  Par le premier théorème d'isomorphisme,
  $(G/N)/(H/N) = (G/N)/\Ker\varphi \cong \Im\varphi = G/H$.
\end{proof}

% ============================================================
\subsection{Théorème de correspondance}\label{subsec:correspondance}
% ============================================================

\begin{theoreme}[Théorème de correspondance]\label{thm:correspondance}
  Soit $N\trianglelefteq G$ et $\pi\colon G\to G/N$ la projection
  canonique.\index{theoreme de correspondance@théorème de correspondance}
  L'application
  \[
    \Phi\colon \{H\leq G \mid N\subseteq H\} \to \{K\leq G/N\},
    \qquad H\mapsto H/N = \pi(H),
  \]
  est une bijection, dont l'inverse est $K\mapsto\pi^{-1}(K)$.
  De plus~:
  \begin{enumerate}[label=(\roman*)]
    \item $H_1\subseteq H_2 \iff H_1/N \subseteq H_2/N$\,;
    \item $H\trianglelefteq G \iff H/N\trianglelefteq G/N$\,;
    \item $[G:H] = [G/N : H/N]$.
  \end{enumerate}
\end{theoreme}

\[
\begin{tikzcd}[row sep=2.5em, column sep=3em]
  \left\{\substack{H\leq G \\ N\subseteq H}\right\}
    \arrow[r,shift left=0.5ex,"\Phi"]
  & \{K\leq G/N\}
    \arrow[l,shift left=0.5ex,"\pi^{-1}"]
\end{tikzcd}
\]

\begin{proof}
  \emph{$\Phi$ est bien définie.}
  Si $N\subseteq H\leq G$, alors $H/N = \pi(H)$ est un sous-groupe de
  $G/N$ car $\pi$ est un homomorphisme et l'image d'un sous-groupe par
  un homomorphisme est un sous-groupe.

  \medskip

  \emph{$\pi^{-1}$ est l'inverse.}
  Pour $K\leq G/N$, posons $H=\pi^{-1}(K)$. Alors $H$ est un
  sous-groupe de $G$ contenant $N$ (car $\pi(n)=N\in K$ pour tout
  $n\in N$). On a $\pi(H)=K$ (surjectivité de $\pi$ sur $K$), donc
  $\Phi(\pi^{-1}(K))=K$.

  Inversement, si $N\subseteq H\leq G$, montrons que
  $\pi^{-1}(\pi(H))=H$. L'inclusion $H\subseteq\pi^{-1}(\pi(H))$ est
  évidente. Pour l'autre sens, si $g\in\pi^{-1}(\pi(H))$, alors
  $gN=hN$ pour un $h\in H$, d'où $g\in hN\subseteq H$ (car
  $N\subseteq H$).

  \medskip

  \textbf{(i)} $H_1\subseteq H_2 \Rightarrow H_1/N\subseteq H_2/N$
  est clair. Réciproquement, si $H_1/N\subseteq H_2/N$ et $h\in H_1$,
  alors $hN\in H_1/N\subseteq H_2/N$, donc $h\in\pi^{-1}(H_2/N)=H_2$.

  \medskip

  \textbf{(ii)} On a vu dans la preuve du
  théorème~\ref{thm:troisieme-iso} que si $H\trianglelefteq G$ et
  $N\subseteq H$, alors $H/N\trianglelefteq G/N$. Réciproquement, si
  $H/N\trianglelefteq G/N$, alors pour tout $g\in G$ et $h\in H$,
  $(gN)(hN)(gN)^{-1} = (ghg^{-1})N\in H/N$, d'où $ghg^{-1}\in H$,
  soit $H\trianglelefteq G$.

  \medskip

  \textbf{(iii)} La bijection $\Phi$ induit une bijection entre les
  classes à gauche de $H$ dans $G$ et celles de $H/N$ dans $G/N$
  par $gH\mapsto (gN)(H/N)$.
\end{proof}

% ============================================================
\section{Automorphismes}\label{sec:automorphismes}
% ============================================================

\begin{definition}[Groupe des automorphismes]\label{def:aut}
  Soit $G$ un groupe. L'ensemble
  $\Aut(G)$\index{automorphisme!groupe des automorphismes}
  de tous les automorphismes de $G$, muni de la composition, est un
  groupe appelé le \emph{groupe des automorphismes} de $G$.
\end{definition}

\begin{definition}[Automorphisme intérieur]\label{def:inn}
  Pour $g\in G$, l'application
  $\iota_g\colon G\to G$ définie par $\iota_g(x) = gxg^{-1}$
  est un automorphisme de $G$ appelé \emph{automorphisme
  intérieur}\index{automorphisme!intérieur} défini par $g$.
  L'ensemble $\Inn(G) = \{\iota_g\mid g\in G\}$ est le
  \emph{groupe des automorphismes intérieurs}.
\end{definition}

\begin{theoreme}\label{thm:inn-normal-aut}
  Soit $G$ un groupe. Alors~:
  \begin{enumerate}[label=(\roman*)]
    \item $\Inn(G)\trianglelefteq\Aut(G)$\,;
    \item $\Inn(G)\cong G/Z(G)$, où $Z(G)$ désigne le centre de $G$.
  \end{enumerate}
\end{theoreme}

\begin{proof}
  \textbf{(ii)} Définissons $\Phi\colon G\to\Aut(G)$ par
  $\Phi(g) = \iota_g$. C'est un homomorphisme car
  \[
    \Phi(gh)(x) = (gh)x(gh)^{-1} = g(hxh^{-1})g^{-1}
    = \iota_g(\iota_h(x)) = (\iota_g\circ\iota_h)(x)
    = (\Phi(g)\circ\Phi(h))(x).
  \]
  Le noyau de $\Phi$ est
  $\Ker\Phi = \{g\in G\mid \iota_g=\Id_G\}
  = \{g\in G\mid gxg^{-1}=x,\;\forall x\in G\} = Z(G)$.
  L'image de $\Phi$ est $\Inn(G)$ par définition.
  Le premier théorème d'isomorphisme donne alors
  $G/Z(G)\cong\Inn(G)$.

  \medskip

  \textbf{(i)} Pour montrer que $\Inn(G)\trianglelefteq\Aut(G)$,
  soit $\varphi\in\Aut(G)$ et $\iota_g\in\Inn(G)$. Pour tout $x\in G$,
  \[
    (\varphi\circ\iota_g\circ\varphi^{-1})(x)
    = \varphi(g\,\varphi^{-1}(x)\,g^{-1})
    = \varphi(g)\cdot x\cdot\varphi(g)^{-1}
    = \iota_{\varphi(g)}(x).
  \]
  Donc $\varphi\circ\iota_g\circ\varphi^{-1} = \iota_{\varphi(g)}\in\Inn(G)$.
\end{proof}

\begin{exemple}\label{ex:aut-Zn}
  Pour $n\geq 1$, un automorphisme $\varphi$ de $\Z/n\Z$ est entièrement
  déterminé par $\varphi(\bar{1})$, et $\varphi(\bar{1})$ doit être un
  générateur de $\Z/n\Z$. Les générateurs de $\Z/n\Z$ sont les
  $\bar{k}$ avec $\gcd(k,n)=1$, donc
  \[
    \Aut(\Z/n\Z) \;\cong\; (\Z/n\Z)^*,
  \]
  le groupe des unités de l'anneau $\Z/n\Z$, qui est d'ordre
  $\varphi(n)$ (indicatrice d'Euler)\index{indicatrice d'Euler}.
\end{exemple}

\begin{exemple}\label{ex:aut-S3}
  On a $\Aut(\Sym_3)\cong\Inn(\Sym_3)\cong\Sym_3/Z(\Sym_3)\cong\Sym_3$,
  car $Z(\Sym_3)=\{e\}$ (le centre du groupe symétrique $\Sym_n$ est
  trivial pour $n\geq 3$).
\end{exemple}

% ============================================================
\section{Exercices}\label{sec:exo-ch3}
% ============================================================

\begin{exercice}[$\star$]\label{exo:hom-trivial}
  Soit $f\colon G\to H$ un homomorphisme de groupes. Montrer que si
  $G$ est abélien et $f$ est surjectif, alors $H$ est abélien.
\end{exercice}

\begin{exercice}[$\star$]\label{exo:hom-ordre}
  Soit $f\colon G\to H$ un homomorphisme de groupes et $g\in G$
  d'ordre fini $n$. Montrer que $\ord(f(g))$ divise $n$.
  Montrer de plus que si $f$ est injectif, alors $\ord(f(g))=n$.
\end{exercice}

\begin{exercice}[$\star$]\label{exo:ker-image-fini}
  Soit $f\colon G\to H$ un homomorphisme de groupes finis.
  Montrer que $|G| = |\Ker f|\cdot|\Im f|$.
\end{exercice}

\begin{exercice}[$\star$]\label{exo:hom-cyclique}
  Montrer que tout homomorphisme $f\colon\Z/n\Z\to\Z/m\Z$ est
  déterminé par $f(\bar{1})$, et que $f(\bar{1})=\bar{k}$ définit un
  homomorphisme si et seulement si $m\mid kn$.
\end{exercice}

\begin{exercice}[$\star\star$]\label{exo:image-reciproque}
  Soit $f\colon G\to H$ un homomorphisme et $K\leq H$.
  Montrer que $f^{-1}(K)\leq G$.
  Montrer que si $K\trianglelefteq H$, alors $f^{-1}(K)\trianglelefteq G$.
\end{exercice}

\begin{exercice}[$\star\star$]\label{exo:produit-direct-quotient}
  Soient $G$ et $H$ deux groupes. Montrer que
  $G\times H / (G\times\{e_H\}) \cong H$
  en utilisant le premier théorème d'isomorphisme.
\end{exercice}

\begin{exercice}[$\star\star$]\label{exo:aut-produit}
  Soient $m,n\geq 2$ des entiers tels que $\gcd(m,n)=1$.
  Montrer que $\Aut(\Z/m\Z\times\Z/n\Z)\cong\Aut(\Z/m\Z)\times\Aut(\Z/n\Z)$.
  \textit{(Indication~: utiliser le fait que $\Z/mn\Z\cong\Z/m\Z\times\Z/n\Z$.)}
\end{exercice}

\begin{exercice}[$\star\star$]\label{exo:sous-groupes-quotient}
  Soit $G=\Z/12\Z$. Déterminer tous les sous-groupes de $G$ et
  vérifier le théorème de correspondance pour $N=\langle\bar{4}\rangle$.
\end{exercice}

\begin{exercice}[$\star\star$]\label{exo:troisieme-iso-application}
  Soit $G = \Z$, $N=6\Z$ et $H=2\Z$. On a $N\subseteq H\subseteq G$ et
  $N,H\trianglelefteq G$. Vérifier que
  $(G/N)/(H/N)\cong G/H$ en identifiant explicitement les deux
  membres.
\end{exercice}

\begin{exercice}[$\star\star\star$]\label{exo:aut-V4}
  Soit $V_4 = \Z/2\Z\times\Z/2\Z$ le groupe de Klein.
  Montrer que $\Aut(V_4)\cong\Sym_3$.
  \textit{(Indication~: $V_4$ possède exactement trois éléments
  d'ordre~$2$, et tout automorphisme permute ces trois éléments.)}
\end{exercice}

\begin{exercice}[$\star\star\star$]\label{exo:inn-cyclique}
  Montrer que si $G/Z(G)$ est cyclique, alors $G$ est abélien.
  En déduire que $\Inn(G)$ ne peut jamais être cyclique non trivial.
\end{exercice}

% ============================================================
\section*{Résumé du chapitre~3}
\addcontentsline{toc}{section}{Résumé du chapitre~3}
% ============================================================

\begin{center}
\fbox{\parbox{0.92\textwidth}{
\textbf{Concepts clés.}
\begin{itemize}
  \item Un \textbf{homomorphisme} $f\colon G\to H$ vérifie
    $f(ab)=f(a)f(b)$; il préserve le neutre et les inverses.
  \item $\Ker f\trianglelefteq G$ et $\Im f\leq H$.
    L'injectivité équivaut à $\Ker f=\{e\}$.
  \item \textbf{1\textsuperscript{er} théorème d'iso.}~:
    $G/\Ker f\cong\Im f$.
  \item \textbf{2\textsuperscript{e} théorème d'iso.}~:
    si $H\leq G$ et $N\trianglelefteq G$, alors $HN/N\cong H/(H\cap N)$.
  \item \textbf{3\textsuperscript{e} théorème d'iso.}~:
    si $N\subseteq H\trianglelefteq G$, alors
    $(G/N)/(H/N)\cong G/H$.
  \item \textbf{Correspondance}~: les sous-groupes de $G/N$ sont en
    bijection avec les sous-groupes de $G$ contenant $N$.
  \item $\Inn(G)\cong G/Z(G)$ et $\Inn(G)\trianglelefteq\Aut(G)$.
\end{itemize}
}}
\end{center}


% ============================================================
%
%  CHAPITRE 4 : ACTIONS DE GROUPES ET THÉORÈMES DE SYLOW
%
% ============================================================

\chapter{Actions de groupes et théorèmes de Sylow}\label{ch:actions-sylow}

\section*{Introduction}
\addcontentsline{toc}{section}{Introduction}

L'idée de \emph{groupe} est née historiquement de l'étude des
symétries~: symétries d'une équation polynomiale (Galois), symétries
d'une figure géométrique (Klein), transformations d'un espace
(Lie). Dans chacun de ces contextes, un groupe \emph{agit} sur un
ensemble --- sur les racines d'un polynôme, sur les sommets d'un
polyèdre, sur les points d'un espace.

\medskip

La formalisation de cette idée --- celle d'\emph{action de groupe} ---
est l'un des outils les plus puissants de l'algèbre. Elle permet de
transformer des questions abstraites sur la structure interne d'un groupe
en des questions concrètes de combinatoire (compter des orbites, des
points fixes, etc.), et réciproquement.

\medskip

Ce chapitre développe la théorie des actions de groupes, en
commençant par la formule orbite-stabilisateur et l'équation aux
classes. Nous en déduisons des résultats structurels importants
(théorème de Cauchy, $p$-groupes) avant de culminer avec les
\emph{théorèmes de Sylow}\index{Sylow}, qui constituent l'un des
sommets de la théorie des groupes finis.

% ============================================================
\section{Actions de groupes}\label{sec:actions}
% ============================================================

\subsection{Définition et exemples}\label{subsec:action-def}

\begin{definition}[Action de groupe]\label{def:action}
  Soit $G$ un groupe et $X$ un ensemble non vide. Une
  \emph{action}\index{action de groupe} (à gauche) de $G$ sur $X$
  est une application $G\times X\to X$, notée $(g,x)\mapsto g\cdot x$,
  satisfaisant~:
  \begin{enumerate}[label=(A\arabic*)]
    \item $e\cdot x = x$ pour tout $x\in X$\,;
    \item $g\cdot(h\cdot x) = (gh)\cdot x$ pour tous $g,h\in G$ et $x\in X$.
  \end{enumerate}
  On dit alors que $G$ \emph{agit sur} $X$, ou que $X$ est un
  \emph{$G$-ensemble}\index{G-ensemble@$G$-ensemble}.
\end{definition}

\begin{proposition}[Formulation comme homomorphisme]\label{prop:action-hom}
  Se donner une action de $G$ sur $X$ équivaut à se donner un
  homomorphisme de groupes $\rho\colon G\to\Sym(X)$, où $\Sym(X)$
  désigne le groupe des bijections de $X$ sur lui-même.
\end{proposition}

\begin{proof}
  Étant donné une action $G\times X\to X$, pour chaque $g\in G$,
  l'application $\sigma_g\colon X\to X$ définie par $\sigma_g(x)=g\cdot x$
  est une bijection de $X$ (d'inverse $\sigma_{g^{-1}}$), et
  l'application $\rho\colon g\mapsto\sigma_g$ est un homomorphisme de
  $G$ dans $\Sym(X)$ car $\sigma_{gh}=\sigma_g\circ\sigma_h$
  (axiome (A2)).

  Réciproquement, un homomorphisme $\rho\colon G\to\Sym(X)$ définit une
  action par $g\cdot x = \rho(g)(x)$.
\end{proof}

\begin{definition}[Noyau d'une action]\label{def:noyau-action}
  Le \emph{noyau}\index{noyau!d'une action} de l'action est
  $\Ker\rho = \{g\in G\mid g\cdot x = x,\;\forall x\in X\}$.
  L'action est dite \emph{fidèle}\index{action!fidèle} si
  $\Ker\rho = \{e\}$.
\end{definition}

\begin{exemple}\label{ex:action-translation}
  Tout groupe $G$ agit sur lui-même par \emph{translation à
  gauche}\index{translation}~: $g\cdot x = gx$. Cette action est
  fidèle. L'homomorphisme associé $\rho\colon G\to\Sym(G)$ est injectif~;
  c'est le \emph{théorème de Cayley}\index{Cayley (théorème de)}~:
  tout groupe est isomorphe à un sous-groupe d'un groupe symétrique.
\end{exemple}

\begin{exemple}\label{ex:action-conjugaison}
  Tout groupe $G$ agit sur lui-même par
  \emph{conjugaison}\index{conjugaison}~:
  $g\cdot x = gxg^{-1}$. Le noyau de cette action est le centre
  $Z(G)$.
\end{exemple}

\begin{exemple}\label{ex:action-sous-groupes}
  Le groupe $G$ agit par conjugaison sur l'ensemble de ses
  sous-groupes~: $g\cdot H = gHg^{-1}$. Le stabilisateur d'un
  sous-groupe $H$ sous cette action est le \emph{normalisateur}
  $N_G(H) = \{g\in G\mid gHg^{-1}=H\}$\index{normalisateur}.
\end{exemple}

% ============================================================
\subsection{Orbites et stabilisateurs}\label{subsec:orbites-stab}
% ============================================================

\begin{definition}[Orbite et stabilisateur]\label{def:orbite-stab}
  Soit $G$ un groupe agissant sur un ensemble $X$ et soit $x\in X$.
  \begin{itemize}
    \item L'\emph{orbite}\index{orbite} de $x$ est
      $\mathrm{Orb}_G(x) = G\cdot x = \{g\cdot x\mid g\in G\}$.
    \item Le \emph{stabilisateur}\index{stabilisateur} de $x$ est
      $\mathrm{Stab}_G(x) = G_x = \{g\in G\mid g\cdot x = x\}$.
  \end{itemize}
\end{definition}

\begin{proposition}\label{prop:stab-sous-groupe}
  Pour tout $x\in X$, le stabilisateur $\mathrm{Stab}_G(x)$ est un
  sous-groupe de $G$.
\end{proposition}

\begin{proof}
  On a $e\cdot x = x$, donc $e\in\mathrm{Stab}_G(x)$. Si
  $g,h\in\mathrm{Stab}_G(x)$, alors
  $(gh^{-1})\cdot x = g\cdot(h^{-1}\cdot x) = g\cdot x = x$
  (car $h^{-1}\cdot x = x$ puisque $h\cdot x=x$). Donc
  $\mathrm{Stab}_G(x)\leq G$.
\end{proof}

\begin{remarque}\label{rem:orbites-partition}
  La relation <<~$x$ et $y$ sont dans la même orbite~>> (c'est-à-dire
  $\exists\, g\in G,\; y = g\cdot x$) est une relation d'équivalence
  sur $X$. Les orbites forment donc une \emph{partition} de $X$.
  Si l'on note $X/G$ l'ensemble des orbites, on a
  $X = \bigsqcup_{\mathcal{O}\in X/G}\mathcal{O}$.
\end{remarque}

% ============================================================
\subsection{Formule orbite-stabilisateur}\label{subsec:orbite-stab}
% ============================================================

\begin{theoreme}[Formule orbite-stabilisateur]\label{thm:orbite-stab}
  Soit $G$ un groupe fini agissant sur un ensemble $X$ et soit $x\in X$.
  Alors\index{formule orbite-stabilisateur}
  \[
    |\mathrm{Orb}_G(x)| = [G:\mathrm{Stab}_G(x)]
    = \frac{|G|}{|\mathrm{Stab}_G(x)|}.
  \]
\end{theoreme}

\begin{proof}
  Posons $S = \mathrm{Stab}_G(x)$ et considérons l'application
  \[
    \Phi\colon G/S \to \mathrm{Orb}_G(x),\qquad gS\mapsto g\cdot x.
  \]

  \emph{Bonne définition.} Si $gS = g'S$, alors $g'^{-1}g\in S$, d'où
  $(g'^{-1}g)\cdot x = x$, soit $g\cdot x = g'\cdot x$. Donc
  $\Phi(gS) = \Phi(g'S)$.

  \emph{Injectivité.} Si $g\cdot x = g'\cdot x$, alors
  $(g'^{-1}g)\cdot x = x$, d'où $g'^{-1}g\in S$, soit $gS = g'S$.

  \emph{Surjectivité.} Par définition de l'orbite, tout élément de
  $\mathrm{Orb}_G(x)$ est de la forme $g\cdot x = \Phi(gS)$.

  Ainsi $\Phi$ est une bijection entre $G/S$ et $\mathrm{Orb}_G(x)$,
  d'où $|\mathrm{Orb}_G(x)| = |G/S| = [G:S]$.
\end{proof}

\begin{corollaire}\label{cor:orbite-divise}
  Si $G$ est fini, alors $|\mathrm{Orb}_G(x)|$ divise $|G|$ pour tout
  $x\in X$.
\end{corollaire}

% ============================================================
\section{Équation aux classes}\label{sec:equation-classes}
% ============================================================

\begin{theoreme}[Équation aux classes]\label{thm:equation-classes}
  Soit $G$ un groupe fini agissant sur un ensemble fini $X$. Soient
  $x_1,\ldots,x_r$ des représentants des orbites distinctes.\index{equation aux classes@équation aux classes}
  Alors
  \[
    |X| = \sum_{i=1}^{r} |\mathrm{Orb}_G(x_i)|
    = \sum_{i=1}^{r} [G:\mathrm{Stab}_G(x_i)].
  \]
\end{theoreme}

\begin{proof}
  Les orbites forment une partition de $X$
  (remarque~\ref{rem:orbites-partition}), donc
  $|X| = \sum_{i=1}^{r} |\mathrm{Orb}_G(x_i)|$.
  Le théorème orbite-stabilisateur~\ref{thm:orbite-stab} donne
  $|\mathrm{Orb}_G(x_i)| = [G:\mathrm{Stab}_G(x_i)]$.
\end{proof}

\begin{remarque}\label{rem:points-fixes}
  Notons $X^G = \{x\in X\mid g\cdot x = x,\;\forall g\in G\}$
  l'ensemble des \emph{points fixes}\index{point fixe} de l'action.
  Un point $x$ est fixe si et seulement si $\mathrm{Orb}_G(x) = \{x\}$.
  En séparant les orbites de cardinal $1$ (les points fixes) des autres,
  l'équation aux classes s'écrit~:
  \[
    |X| = |X^G| + \sum_{\substack{i=1\\|\mathrm{Orb}_G(x_i)|>1}}^{r}
    [G:\mathrm{Stab}_G(x_i)].
  \]
\end{remarque}

% ============================================================
\subsection{Lemme de Burnside}\label{subsec:burnside}
% ============================================================

\begin{theoreme}[Lemme de Burnside]\label{thm:burnside}
  Soit $G$ un groupe fini agissant sur un ensemble fini $X$. Le nombre
  d'orbites de cette action est\index{Burnside (lemme de)}
  \[
    |X/G| = \frac{1}{|G|}\sum_{g\in G} |X^g|,
  \]
  où $X^g = \{x\in X\mid g\cdot x = x\}$ désigne l'ensemble des
  points fixés par $g$.
\end{theoreme}

\begin{proof}
  Considérons l'ensemble
  $\Omega = \{(g,x)\in G\times X \mid g\cdot x = x\}$
  et comptons $|\Omega|$ de deux manières.

  \emph{En sommant sur $g$.} $|\Omega| = \sum_{g\in G}|X^g|$.

  \emph{En sommant sur $x$.}
  $|\Omega| = \sum_{x\in X}|\mathrm{Stab}_G(x)|$.
  Or, par le théorème orbite-stabilisateur,
  $|\mathrm{Stab}_G(x)| = |G|/|\mathrm{Orb}_G(x)|$.
  Donc
  \[
    |\Omega| = \sum_{x\in X}\frac{|G|}{|\mathrm{Orb}_G(x)|}
    = |G|\sum_{x\in X}\frac{1}{|\mathrm{Orb}_G(x)|}.
  \]
  Or, pour chaque orbite $\mathcal{O}$,
  $\sum_{x\in\mathcal{O}}\frac{1}{|\mathcal{O}|} = 1$, donc
  \[
    \sum_{x\in X}\frac{1}{|\mathrm{Orb}_G(x)|} = |X/G|.
  \]
  En égalant les deux expressions de $|\Omega|$,
  $\sum_{g\in G}|X^g| = |G|\cdot|X/G|$, d'où le résultat.
\end{proof}

\begin{exemple}[Comptage de colliers]\label{ex:colliers}
  Combien y a-t-il de colliers\index{collier} distincts formés de $4$
  perles, chacune pouvant être de $3$ couleurs, si deux colliers sont
  considérés identiques lorsqu'on peut passer de l'un à l'autre par
  rotation~?

  Le groupe agissant est $G=\Z/4\Z$ (rotations) et
  $X = \{1,2,3\}^4$ (coloriages).
  On a $|X|=81$ et $|G|=4$.
  \begin{itemize}
    \item $g=\bar{0}$ (identité)~: $|X^g| = 81$\,;
    \item $g=\bar{1}$ (rotation de $90°$)~: toutes les perles doivent
      être de même couleur, $|X^g|=3$\,;
    \item $g=\bar{2}$ (rotation de $180°$)~: perles opposées de même
      couleur, $|X^g|=3^2=9$\,;
    \item $g=\bar{3}$ (rotation de $270°$)~: comme $\bar{1}$,
      $|X^g|=3$.
  \end{itemize}
  Par Burnside, le nombre de colliers est
  $\frac{1}{4}(81+3+9+3) = \frac{96}{4} = 24$.
\end{exemple}

% ============================================================
\section{Conjugaison et équation aux classes d'un groupe}\label{sec:conjugaison}
% ============================================================

\begin{definition}[Classes de conjugaison]\label{def:classe-conjugaison}
  Soit $G$ un groupe. Deux éléments $x,y\in G$ sont dits
  \emph{conjugués}\index{conjugaison!classe de} s'il existe $g\in G$
  tel que $y=gxg^{-1}$. Les orbites de l'action de $G$ sur lui-même
  par conjugaison sont les \emph{classes de conjugaison} de $G$.
  La classe de $x$ est notée $\mathrm{Cl}(x)$.
\end{definition}

\begin{definition}[Centralisateur]\label{def:centralisateur}
  Le \emph{centralisateur}\index{centralisateur} d'un élément $x\in G$
  est $C_G(x) = \{g\in G\mid gxg^{-1}=x\} = \mathrm{Stab}_G(x)$
  pour l'action par conjugaison.
\end{definition}

\begin{proposition}[Équation aux classes d'un groupe]\label{prop:equation-classes-groupe}
  Soit $G$ un groupe fini. Soient $x_1,\ldots,x_r$ des représentants
  des classes de conjugaison non centrales ($|\mathrm{Cl}(x_i)|>1$).
  Alors\index{equation aux classes@équation aux classes!d'un groupe}
  \[
    |G| = |Z(G)| + \sum_{i=1}^{r} [G:C_G(x_i)].
  \]
\end{proposition}

\begin{proof}
  C'est un cas particulier de l'équation aux classes
  (théorème~\ref{thm:equation-classes}) appliquée à l'action de $G$ sur
  lui-même par conjugaison. Les points fixes sont exactement les
  éléments du centre~: $x\in G^G \iff gxg^{-1}=x$ pour tout
  $g\in G \iff x\in Z(G)$.
\end{proof}

% ============================================================
\subsection{Applications aux $p$-groupes}\label{subsec:p-groupes}
% ============================================================

\begin{definition}[$p$-groupe]\label{def:p-groupe}
  Soit $p$ un nombre premier. Un groupe fini $G$ est un
  \emph{$p$-groupe}\index{p-groupe@$p$-groupe} si $|G|=p^k$ pour un
  entier $k\geq 0$.
\end{definition}

\begin{theoreme}[Centre d'un $p$-groupe]\label{thm:centre-p-groupe}
  Soit $p$ un nombre premier et $G$ un $p$-groupe non trivial
  ($|G|=p^k$, $k\geq 1$). Alors $Z(G)\neq\{e\}$.
\end{theoreme}

\begin{proof}
  Appliquons l'équation aux classes de la
  proposition~\ref{prop:equation-classes-groupe}~:
  \[
    |G| = |Z(G)| + \sum_{i=1}^{r}[G:C_G(x_i)].
  \]
  Pour chaque $i$, $[G:C_G(x_i)] = |\mathrm{Cl}(x_i)|>1$, et
  $[G:C_G(x_i)]$ divise $|G|=p^k$, donc $p\mid [G:C_G(x_i)]$.
  Ainsi $p$ divise $\sum_{i=1}^{r}[G:C_G(x_i)]$ et $p$ divise $|G|$,
  d'où $p$ divise $|Z(G)|$. En particulier, $|Z(G)|\geq p > 1$.
\end{proof}

\begin{theoreme}[Groupes d'ordre $p^2$]\label{thm:ordre-p2}
  Tout groupe d'ordre $p^2$ ($p$ premier) est abélien.
\end{theoreme}

\begin{proof}
  Soit $|G|=p^2$. Par le théorème~\ref{thm:centre-p-groupe},
  $Z(G)\neq\{e\}$, donc $|Z(G)|\in\{p,p^2\}$ (car $|Z(G)|$ divise
  $|G|=p^2$ et $|Z(G)|>1$).

  Si $|Z(G)|=p^2$, alors $G = Z(G)$ et $G$ est abélien.

  Supposons $|Z(G)|=p$. Alors $|G/Z(G)| = p^2/p = p$, donc $G/Z(G)$
  est cyclique (tout groupe d'ordre premier est cyclique). Mais si
  $G/Z(G)$ est cyclique, alors $G$ est abélien (cf.\
  exercice~\ref{exo:inn-cyclique}), ce qui contredit $|Z(G)|=p<p^2$.
  Donc ce cas est impossible.
\end{proof}

% ============================================================
\section{Théorème de Cauchy}\label{sec:cauchy}
% ============================================================

\begin{theoreme}[Théorème de Cauchy]\label{thm:cauchy}
  Soit $G$ un groupe fini et $p$ un nombre premier divisant $|G|$.
  Alors $G$ possède un élément d'ordre $p$.\index{Cauchy (théorème de)}
\end{theoreme}

\begin{proof}
  On procède par récurrence forte sur $|G|$.

  \emph{Base.} Si $|G|=1$, aucun premier ne divise $|G|$, donc
  l'énoncé est trivialement vrai.

  \emph{Étape de récurrence.} Soit $|G|>1$ et supposons le résultat
  vrai pour tout groupe d'ordre strictement inférieur à $|G|$.
  Considérons l'équation aux classes
  (proposition~\ref{prop:equation-classes-groupe})~:
  \[
    |G| = |Z(G)| + \sum_{i=1}^{r}[G:C_G(x_i)].
  \]

  \textbf{Cas 1.} Il existe $i$ tel que $p\nmid [G:C_G(x_i)]$.
  Puisque $p\mid |G|$ et $|G| = [G:C_G(x_i)]\cdot|C_G(x_i)|$,
  on a $p\mid |C_G(x_i)|$. Or $C_G(x_i)\neq G$ (car
  $|\mathrm{Cl}(x_i)|>1$, donc $x_i\notin Z(G)$), d'où
  $|C_G(x_i)|<|G|$.
  Par hypothèse de récurrence, $C_G(x_i)$ possède un élément d'ordre
  $p$, qui est aussi un élément de $G$ d'ordre $p$.

  \textbf{Cas 2.} Pour tout $i$, $p\mid [G:C_G(x_i)]$.
  Alors $p$ divise $\sum_i [G:C_G(x_i)]$, et puisque $p\mid |G|$,
  on déduit $p\mid |Z(G)|$.
  Comme $Z(G)$ est abélien, on peut l'écrire (par le théorème des
  groupes abéliens finis) comme un produit de groupes cycliques. Puisque
  $p\mid |Z(G)|$, au moins un de ces facteurs cycliques est d'ordre
  divisible par $p$. Si $\Z/m\Z$ est un tel facteur avec $p\mid m$,
  alors $\Z/m\Z$ contient un élément d'ordre $p$ (l'élément
  $\overline{m/p}$). Cet élément est un élément d'ordre $p$ dans $G$.
\end{proof}

\begin{remarque}\label{rem:cauchy-elementaire}
  On peut donner une preuve plus élémentaire qui ne fait pas appel
  au théorème de structure des groupes abéliens finis. Si $Z(G)$ est
  abélien et $p\mid |Z(G)|$, il suffit de prendre un élément
  $a\in Z(G)$ d'ordre $n>1$. Si $p\mid n$, alors $a^{n/p}$ est
  d'ordre $p$. Sinon, considérer le quotient $Z(G)/\langle a\rangle$
  et raisonner par récurrence sur l'ordre.
\end{remarque}

% ============================================================
\section{Sous-groupes de Sylow}\label{sec:sylow}
% ============================================================

\begin{definition}[Sous-groupe de Sylow]\label{def:sylow}
  Soit $G$ un groupe fini d'ordre $|G|=p^k m$ avec $p$ premier et
  $\gcd(p,m)=1$ (c'est-à-dire $p^k$ est la plus grande puissance de $p$
  divisant $|G|$). Un \emph{$p$-sous-groupe de
  Sylow}\index{Sylow!sous-groupe de} (ou simplement \emph{Sylow
  $p$-sous-groupe}) de $G$ est un sous-groupe d'ordre $p^k$.
  On note $\mathrm{Syl}_p(G)$ l'ensemble des Sylow $p$-sous-groupes
  de $G$ et $n_p = |\mathrm{Syl}_p(G)|$ leur nombre.
\end{definition}

% ============================================================
\subsection{Premier théorème de Sylow}\label{subsec:sylow-1}
% ============================================================

\begin{theoreme}[Premier théorème de Sylow]\label{thm:sylow-1}
  Soit $G$ un groupe fini et $p$ un nombre premier. Si $p^k$ divise
  $|G|$, alors $G$ possède un sous-groupe d'ordre
  $p^k$.\index{Sylow!premier théorème}
  En particulier, les Sylow $p$-sous-groupes existent.
\end{theoreme}

\begin{proof}
  On procède par récurrence forte sur $|G|$. Le résultat est trivial
  pour $|G|=1$.

  Supposons $|G|>1$ et le résultat vrai pour tout groupe d'ordre
  strictement inférieur.

  \medskip

  \textbf{Cas 1~: $p\mid |Z(G)|$.}
  Par le théorème de Cauchy (théorème~\ref{thm:cauchy}),
  $Z(G)$ possède un élément $a$ d'ordre $p$. Posons
  $N=\langle a\rangle\trianglelefteq G$ (car $a\in Z(G)$).
  Le quotient $G/N$ est d'ordre $|G|/p$ et $p^{k-1}$ divise
  $|G/N|$. Par hypothèse de récurrence, $G/N$ possède un sous-groupe
  $\bar{H}$ d'ordre $p^{k-1}$. Par le théorème de correspondance
  (théorème~\ref{thm:correspondance}), $\bar{H}=H/N$ pour un
  sous-groupe $H$ de $G$ contenant $N$, et
  $|H| = |\bar{H}|\cdot|N| = p^{k-1}\cdot p = p^k$.

  \medskip

  \textbf{Cas 2~: $p\nmid |Z(G)|$.}
  Considérons l'équation aux classes~:
  $|G| = |Z(G)| + \sum_{i=1}^r [G:C_G(x_i)]$.
  Puisque $p\mid |G|$ et $p\nmid |Z(G)|$, il existe un $i$ tel que
  $p\nmid [G:C_G(x_i)]$.
  Posons $C=C_G(x_i)$; c'est un sous-groupe propre de $G$
  ($|\mathrm{Cl}(x_i)|>1$).
  Comme $|G|=|C|\cdot[G:C]$ et $p^k\mid|G|$ mais
  $p\nmid [G:C]$, on a $p^k\mid |C|$.
  Par hypothèse de récurrence (car $|C|<|G|$), $C$ possède un
  sous-groupe d'ordre $p^k$, qui est aussi un sous-groupe de $G$
  d'ordre $p^k$.
\end{proof}

% ============================================================
\subsection{Deuxième théorème de Sylow}\label{subsec:sylow-2}
% ============================================================

\begin{lemme}\label{lem:p-groupe-orbite}
  Soit $P$ un $p$-groupe agissant sur un ensemble fini $X$. Alors
  $|X|\equiv |X^P|\pmod{p}$,
  où $X^P = \{x\in X\mid g\cdot x = x,\;\forall g\in P\}$ est
  l'ensemble des points fixes.
\end{lemme}

\begin{proof}
  Par l'équation aux classes,
  $|X| = |X^P| + \sum_{j} |\mathrm{Orb}_P(x_j)|$,
  où la somme porte sur les orbites de cardinal $>1$. Pour chacune,
  $|\mathrm{Orb}_P(x_j)| = [P:\mathrm{Stab}_P(x_j)]$ divise
  $|P|=p^k$, et $|\mathrm{Orb}_P(x_j)|>1$, donc
  $p\mid|\mathrm{Orb}_P(x_j)|$. D'où
  $|X|\equiv |X^P|\pmod{p}$.
\end{proof}

\begin{theoreme}[Deuxième théorème de Sylow]\label{thm:sylow-2}
  Soit $G$ un groupe fini et $p$ un nombre premier. Tous les Sylow
  $p$-sous-groupes de $G$ sont conjugués~: si $P$ et $Q$ sont deux
  Sylow $p$-sous-groupes, il existe $g\in G$ tel que
  $Q = gPg^{-1}$.\index{Sylow!deuxième théorème}
\end{theoreme}

\begin{proof}
  Soit $P$ un Sylow $p$-sous-groupe de $G$ et posons
  $\Sigma = \{gPg^{-1}\mid g\in G\}$ l'ensemble des conjugués de $P$
  (qui sont tous des Sylow $p$-sous-groupes de $G$). Nous devons
  montrer que tout Sylow $p$-sous-groupe $Q$ appartient à $\Sigma$.

  Faisons agir $Q$ sur $\Sigma$ par conjugaison~:
  $q\cdot(gPg^{-1}) = (qg)P(qg)^{-1}$.
  Par le lemme~\ref{lem:p-groupe-orbite},
  $|\Sigma|\equiv|\Sigma^Q|\pmod{p}$.

  Un élément $gPg^{-1}\in\Sigma$ est fixé par $Q$ si et seulement si
  $Q\subseteq N_G(gPg^{-1})$, le normalisateur de $gPg^{-1}$.
  Si $gPg^{-1}\in\Sigma^Q$, alors $Q$ et $gPg^{-1}$ sont tous deux
  des Sylow $p$-sous-groupes du groupe $N_G(gPg^{-1})$, et
  $gPg^{-1}\trianglelefteq N_G(gPg^{-1})$.
  Considérons le quotient $N_G(gPg^{-1})/(gPg^{-1})$. L'image de $Q$
  dans ce quotient est $Q(gPg^{-1})/(gPg^{-1})$, qui est un
  $p$-groupe par le deuxième théorème d'isomorphisme
  (théorème~\ref{thm:deuxieme-iso}). Mais l'ordre de ce quotient
  divise $[N_G(gPg^{-1}):gPg^{-1}]$, qui n'est pas divisible par $p$
  (car $gPg^{-1}$ est déjà un Sylow $p$-sous-groupe). Donc
  $Q(gPg^{-1})/(gPg^{-1})$ est trivial, d'où $Q\subseteq gPg^{-1}$.
  Comme $|Q|=|gPg^{-1}|=p^k$, on conclut $Q = gPg^{-1}$.

  Ainsi, $\Sigma^Q = \{gPg^{-1}\}$ tel que $Q = gPg^{-1}$, donc
  $|\Sigma^Q|=1$. En particulier $Q\in\Sigma$.
\end{proof}

% ============================================================
\subsection{Troisième théorème de Sylow}\label{subsec:sylow-3}
% ============================================================

\begin{theoreme}[Troisième théorème de Sylow]\label{thm:sylow-3}
  Soit $G$ un groupe fini d'ordre $|G|=p^k m$ avec $\gcd(p,m)=1$.
  Alors\index{Sylow!troisième théorème}~:
  \begin{enumerate}[label=(\roman*)]
    \item $n_p\equiv 1\pmod{p}$\,;
    \item $n_p$ divise $m = |G|/p^k$.
  \end{enumerate}
\end{theoreme}

\begin{proof}
  \textbf{(i)} Soit $P$ un Sylow $p$-sous-groupe et
  $\Sigma = \mathrm{Syl}_p(G)$. Par le deuxième théorème de Sylow
  (théorème~\ref{thm:sylow-2}), $\Sigma$ est l'orbite de $P$ sous
  l'action de $G$ par conjugaison, donc $n_p = |\Sigma|$.

  Faisons agir $P$ sur $\Sigma$ par conjugaison. Par le
  lemme~\ref{lem:p-groupe-orbite},
  $|\Sigma|\equiv|\Sigma^P|\pmod{p}$.
  Dans la preuve du théorème~\ref{thm:sylow-2}, nous avons montré que
  si $Q\in\Sigma$ est fixé par un Sylow $p$-sous-groupe, alors
  $Q$ coïncide avec ce sous-groupe. Donc $\Sigma^P = \{P\}$,
  d'où $|\Sigma^P|=1$ et $n_p\equiv 1\pmod{p}$.

  \medskip

  \textbf{(ii)} Le groupe $G$ agit transitivement sur $\Sigma$ par
  conjugaison (théorème~\ref{thm:sylow-2}), donc
  $n_p = |\Sigma| = [G:N_G(P)]$, où $N_G(P)$ est le normalisateur
  de $P$ dans $G$. Puisque $P\leq N_G(P)$, on a $p^k\mid |N_G(P)|$,
  d'où
  \[
    n_p = \frac{|G|}{|N_G(P)|} \;\Bigg|\; \frac{|G|}{p^k} = m. \qedhere
  \]
\end{proof}

% ============================================================
\section{Applications des théorèmes de Sylow}\label{sec:applications-sylow}
% ============================================================

\subsection{Groupes d'ordre $pq$}\label{subsec:ordre-pq}

\begin{proposition}[Groupes d'ordre $pq$]\label{prop:ordre-pq}
  Soient $p<q$ deux nombres premiers tels que $q\not\equiv 1\pmod{p}$.
  Alors tout groupe d'ordre $pq$ est cyclique.
\end{proposition}

\begin{proof}
  Soit $|G|=pq$. Par le troisième théorème de Sylow
  (théorème~\ref{thm:sylow-3})~:
  \begin{itemize}
    \item $n_q\mid p$ et $n_q\equiv 1\pmod{q}$. Donc $n_q\in\{1,p\}$.
      Puisque $p<q$, $p\not\equiv 1\pmod{q}$, donc $n_q=1$.
    \item $n_p\mid q$ et $n_p\equiv 1\pmod{p}$. Donc $n_p\in\{1,q\}$.
      Puisque $q\not\equiv 1\pmod{p}$, on a $n_p=1$.
  \end{itemize}
  Soient $P$ et $Q$ les uniques Sylow $p$- et $q$-sous-groupes. Ils
  sont distingués ($n_p=n_q=1$), d'ordres $p$ et $q$ respectivement
  (donc cycliques). De plus, $P\cap Q=\{e\}$ (car $\gcd(p,q)=1$)
  et $|PQ| = |P|\cdot|Q|/|P\cap Q| = pq = |G|$, donc $G=PQ$.
  Comme $P$ et $Q$ sont distingués et $P\cap Q=\{e\}$, on a
  $G\cong P\times Q\cong\Z/p\Z\times\Z/q\Z\cong\Z/pq\Z$.
\end{proof}

\begin{corollaire}\label{cor:ordre-15}
  Tout groupe d'ordre $15$ est cyclique.
\end{corollaire}

\begin{proof}
  $15=3\times 5$. On a $p=3$, $q=5$ et $q=5\not\equiv 1\pmod{3}$. La
  proposition~\ref{prop:ordre-pq} s'applique.
\end{proof}

\subsection{Groupes d'ordre 12}\label{subsec:ordre-12}

\begin{proposition}\label{prop:pas-simple-12}
  Il n'existe pas de groupe simple d'ordre $12$.
\end{proposition}

\begin{proof}
  Soit $G$ un groupe d'ordre $12 = 2^2\cdot 3$.

  Par le troisième théorème de Sylow, $n_3\mid 4$ et
  $n_3\equiv 1\pmod{3}$, donc $n_3\in\{1,4\}$.

  \textbf{Cas 1~: $n_3=1$.} L'unique Sylow $3$-sous-groupe est
  distingué, donc $G$ n'est pas simple.

  \textbf{Cas 2~: $n_3=4$.} $G$ agit par conjugaison sur
  $\mathrm{Syl}_3(G)$ (qui a $4$ éléments), ce qui définit un
  homomorphisme $\varphi\colon G\to\Sym_4$. Le noyau
  $\Ker\varphi\trianglelefteq G$. Puisque l'action est transitive sur
  $\mathrm{Syl}_3(G)$, $\varphi$ n'est pas triviale, donc
  $\Ker\varphi\neq G$.

  Si $\Ker\varphi = \{e\}$, alors $G$ s'injecte dans $\Sym_4$
  ($|G|=12$ divise $|\Sym_4|=24$). Les $4$ Sylow $3$-sous-groupes
  fournissent $4\times 2 = 8$ éléments d'ordre $3$. Il reste $12-8=4$
  éléments dans $G$, qui doivent constituer un unique Sylow
  $2$-sous-groupe (d'ordre $4$). Ce Sylow $2$-sous-groupe est alors
  distingué, et $G$ n'est pas simple.

  Si $\{e\}\subsetneq\Ker\varphi\subsetneq G$, alors $G$ possède un
  sous-groupe distingué non trivial, et n'est pas simple.

  Dans tous les cas, $G$ n'est pas simple.
\end{proof}

\subsection{Diagramme de sous-groupes}\label{subsec:diagramme-sylow}

Le diagramme suivant illustre le treillis des sous-groupes d'un
groupe $G$ d'ordre $12$ avec $n_3=1$ et $n_2=1$ (cas
$G\cong\Z/12\Z$), en mettant en évidence les Sylow sous-groupes.

\begin{center}
\begin{tikzpicture}[
    every node/.style={draw, circle, minimum size=6mm, inner sep=1pt,
      font=\footnotesize},
    sylow/.style={fill=blue!15},
    level/.style={sibling distance=20mm},
    edge from parent/.style={draw,-}
  ]
  % Nodes
  \node (G) at (0,4) {$G$};
  \node[sylow] (P2) at (-2.5,2.5) {$P_2$};
  \node (6) at (0,2.5) {$\langle 2\rangle$};
  \node[sylow] (P3) at (2.5,2.5) {$P_3$};
  \node (2a) at (-2.5,1) {$\langle 6\rangle$};
  \node (3) at (0,1) {$\langle 4\rangle$};
  \node (e) at (0,-0.5) {$\{0\}$};
  % Edges
  \draw (G) -- (P2);
  \draw (G) -- (6);
  \draw (G) -- (P3);
  \draw (P2) -- (2a);
  \draw (P2) -- (3);
  \draw (6) -- (2a);
  \draw (6) -- (3);
  \draw (P3) -- (3);
  \draw (2a) -- (e);
  \draw (3) -- (e);
  % Labels
  \node[draw=none,right,font=\scriptsize\itshape] at (3,2.5) {Sylow $3$};
  \node[draw=none,left,font=\scriptsize\itshape] at (-3.8,2.5) {Sylow $2$};
\end{tikzpicture}
\captionof{figure}{Treillis des sous-groupes de $\Z/12\Z$, avec les
  Sylow sous-groupes en bleu.}
\label{fig:treillis-Z12}
\end{center}

% ============================================================
\section{Exercices}\label{sec:exo-ch4}
% ============================================================

\begin{exercice}[$\star$]\label{exo:action-verification}
  Vérifier que l'action de $\Sym_3$ sur $\{1,2,3\}$ par permutation
  naturelle satisfait les axiomes (A1) et (A2).
  Déterminer l'orbite et le stabilisateur de chaque élément.
\end{exercice}

\begin{exercice}[$\star$]\label{exo:orbite-stab-D4}
  Le groupe diédral $D_4$ (d'ordre $8$) agit sur les sommets d'un
  carré. Déterminer toutes les orbites et tous les stabilisateurs.
  Vérifier la formule orbite-stabilisateur.
\end{exercice}

\begin{exercice}[$\star$]\label{exo:conjugaison-S3}
  Déterminer les classes de conjugaison de $\Sym_3$ et vérifier
  l'équation aux classes.
\end{exercice}

\begin{exercice}[$\star$]\label{exo:centre-Dn}
  Déterminer le centre du groupe diédral $D_n$ pour tout $n\geq 3$.
\end{exercice}

\begin{exercice}[$\star\star$]\label{exo:burnside-cube}
  De combien de manières peut-on colorier les faces d'un cube avec $k$
  couleurs, si deux coloriages sont considérés identiques lorsqu'on peut
  passer de l'un à l'autre par une rotation~?
  \textit{(Le groupe des rotations du cube est isomorphe à $\Sym_4$,
  d'ordre $24$.)}
\end{exercice}

\begin{exercice}[$\star\star$]\label{exo:p-groupe-sous-groupe-distingue}
  Soit $G$ un $p$-groupe d'ordre $p^n$ ($n\geq 1$). Montrer que $G$
  possède un sous-groupe distingué d'ordre $p^k$ pour tout
  $0\leq k\leq n$.
  \textit{(Raisonner par récurrence en utilisant le fait que
  $Z(G)\neq\{e\}$.)}
\end{exercice}

\begin{exercice}[$\star\star$]\label{exo:normalizer-sylow}
  Soit $G$ un groupe fini et $P$ un Sylow $p$-sous-groupe de $G$.
  Montrer que $N_G(N_G(P)) = N_G(P)$.
\end{exercice}

\begin{exercice}[$\star\star$]\label{exo:unique-sylow-normal}
  Montrer que si $n_p=1$, alors l'unique Sylow $p$-sous-groupe est
  distingué dans $G$.
\end{exercice}

\begin{exercice}[$\star\star$]\label{exo:sylow-21}
  Déterminer les possibilités pour $n_3$ et $n_7$ dans un groupe
  d'ordre $21$. En déduire que tout groupe d'ordre $21$ possède un
  unique Sylow $7$-sous-groupe, qui est distingué.
\end{exercice}

\begin{exercice}[$\star\star\star$]\label{exo:pas-simple-36}
  Montrer qu'il n'existe pas de groupe simple d'ordre $36$.
  \textit{(Considérer $n_3$ et raisonner comme dans la
  proposition~\ref{prop:pas-simple-12}.)}
\end{exercice}

\begin{exercice}[$\star\star\star$]\label{exo:classification-p2q}
  Soient $p<q$ deux nombres premiers. Classifier (à isomorphisme
  près) les groupes d'ordre $p^2 q$ lorsque $q\not\equiv 1\pmod{p}$
  et $q\not\equiv 1\pmod{p^2}$.
\end{exercice}

\begin{exercice}[$\star\star\star$]\label{exo:sylow-intersection}
  Soit $G$ un groupe fini et $p$ un premier tel que $n_p>1$. Montrer
  qu'il existe deux Sylow $p$-sous-groupes distincts $P_1,P_2$ tels
  que $P_1\cap P_2$ est le plus grand $p$-sous-groupe distingué de $G$
  contenu dans tout Sylow $p$-sous-groupe, c'est-à-dire
  $P_1\cap P_2 = O_p(G) := \bigcap_{P\in\mathrm{Syl}_p(G)} P$.
\end{exercice}

\begin{exercice}[$\star\star\star$]\label{exo:simple-60}
  Soit $G$ un groupe simple d'ordre $60$. Montrer que $G\cong\Alt_5$.
  \textit{(Montrer que $n_5=6$, puis que l'action de $G$ sur
  $\mathrm{Syl}_5(G)$ donne un monomorphisme $G\hookrightarrow\Sym_6$
  dont l'image est contenue dans $\Alt_6$.)}
\end{exercice}

% ============================================================
\section*{Résumé du chapitre~4}
\addcontentsline{toc}{section}{Résumé du chapitre~4}
% ============================================================

\begin{center}
\fbox{\parbox{0.92\textwidth}{
\textbf{Concepts clés.}
\begin{itemize}
  \item Une \textbf{action} de $G$ sur $X$ est un homomorphisme
    $G\to\Sym(X)$.
  \item \textbf{Orbite-stabilisateur}~:
    $|\mathrm{Orb}(x)| = [G:\mathrm{Stab}(x)]$.
  \item \textbf{Équation aux classes}~:
    $|X| = |X^G| + \sum [G:\mathrm{Stab}(x_i)]$.
  \item \textbf{Burnside}~:
    $|X/G| = \frac{1}{|G|}\sum_{g\in G}|X^g|$.
  \item Le \textbf{centre d'un $p$-groupe} non trivial est non trivial.
  \item Tout groupe d'ordre $p^2$ est \textbf{abélien}.
  \item \textbf{Cauchy}~: si $p\mid|G|$, il existe un élément d'ordre $p$.
  \item \textbf{Sylow 1}~: les Sylow $p$-sous-groupes existent.
  \item \textbf{Sylow 2}~: ils sont tous conjugués.
  \item \textbf{Sylow 3}~:
    $n_p\equiv 1\pmod{p}$ et $n_p\mid |G|/p^k$.
\end{itemize}
}}
\end{center}
%%%%%%%%%%%%%%%%%%%%%%%%%%%%%%%%%%%%%%%%%%%%%%%%%%%%%%%%%%%%
%% ALGEBRE ABSTRAITE I — PARTIE 3 (CHAPITRES 5, 6, 7)
%% Anneaux, Idéaux, Quotients, Corps de fractions,
%% Anneaux principaux, euclidiens et factoriels
%%%%%%%%%%%%%%%%%%%%%%%%%%%%%%%%%%%%%%%%%%%%%%%%%%%%%%%%%%%%

%% ==========================================================
%%  CHAPITRE 5 : Anneaux — Définitions, Exemples, Idéaux
%% ==========================================================

\chapter{Anneaux --- Définitions, exemples, idéaux}\label{chap:anneaux}
\index{anneau}

\section*{Introduction}

L'étude des anneaux trouve son origine dans les travaux de
\textsc{Dedekind}\index{Dedekind, Richard}, qui introduisit la notion
d'\emph{ordre} dans les corps de nombres algébriques pour comprendre
la factorisation des entiers dans des extensions de~$\Z$.
\textsc{Hilbert}\index{Hilbert, David} systématisa l'usage du mot
\emph{Zahlring} (anneau de nombres) dans son \emph{Zahlbericht} (1897),
et c'est \textsc{Noether}\index{Noether, Emmy} qui, dans les
années~1920, dégagea la théorie abstraite des anneaux et de leurs
idéaux dans toute sa généralité, posant les fondements de l'algèbre
commutative moderne.

Le passage du groupe à l'anneau enrichit considérablement la structure~:
on dispose désormais de \emph{deux} opérations, l'addition et la
multiplication, liées par la distributivité. Cette structure apparaît
naturellement dans $\Z$, dans les anneaux de polynômes~$K[X]$, dans
les anneaux de matrices~$M_n(K)$, et dans bien d'autres contextes.
Le présent chapitre introduit les définitions fondamentales, étudie
les premières propriétés, puis développe la théorie des idéaux.

%% ----------------------------------------------------------
\section{Définition d'un anneau}\label{sec:def-anneau}
\index{anneau!définition}

\begin{definition}[Anneau]\label{def:anneau}
Un \emph{anneau}\index{anneau} est un triplet $(A, +, \cdot)$ où $A$ est
un ensemble muni de deux lois de composition internes vérifiant~:
\begin{enumerate}[label=\textup{(A\arabic*)}]
  \item\label{ax:A1} $(A, +)$ est un groupe abélien, d'élément neutre
        noté~$0_A$ (ou simplement~$0$).
  \item\label{ax:A2} La multiplication est associative~:
        $\forall\, a,b,c \in A,\; (a \cdot b) \cdot c = a \cdot (b \cdot c)$.
  \item\label{ax:A3} Il existe un élément $1_A \in A$ (ou simplement~$1$),
        appelé \emph{élément unité}\index{élément unité}, tel que
        $\forall\, a \in A,\; 1_A \cdot a = a \cdot 1_A = a$.
  \item\label{ax:A4} La multiplication est distributive par rapport à
        l'addition~:
        \[
          \forall\, a,b,c \in A,\quad
          a(b+c) = ab + ac \quad\text{et}\quad (a+b)c = ac + bc.
        \]
\end{enumerate}
Si de plus la multiplication est commutative, on dit que $A$ est un
\emph{anneau commutatif}\index{anneau!commutatif}.
\end{definition}

\begin{remarque}[Convention]\label{rem:convention-unite}
Dans cet ouvrage, sauf mention explicite du contraire, tous les anneaux
possèdent un élément unité~$1_A$, et les morphismes d'anneaux respectent
l'unité (cf.\ la section~\ref{sec:morphismes-anneaux}). Cette convention
suit l'usage de \textsc{Bourbaki}\index{Bourbaki}.
\end{remarque}

\begin{remarque}[Anneau trivial]\label{rem:anneau-trivial}
Si $1_A = 0_A$, alors pour tout $a \in A$,
$a = 1_A \cdot a = 0_A \cdot a = 0_A$ (cf.\ la
proposition~\ref{prop:proprietes-elementaires}), d'où $A = \{0\}$.
On appelle cet anneau l'\emph{anneau trivial}\index{anneau!trivial} ou
\emph{anneau nul}.
\end{remarque}

%% ----------------------------------------------------------
\section{Propriétés élémentaires}\label{sec:proprietes-elementaires}

\begin{proposition}[Propriétés élémentaires d'un anneau]%
\label{prop:proprietes-elementaires}
Soit $(A, +, \cdot)$ un anneau. Pour tous $a, b \in A$~:
\begin{enumerate}[label=\textup{(\roman*)}]
  \item\label{it:zero-abs} $0 \cdot a = a \cdot 0 = 0$ \quad
        ($0$ est \emph{absorbant}\index{élément absorbant}).
  \item\label{it:moins-un} $(-1) \cdot a = a \cdot (-1) = -a$.
  \item\label{it:moins-moins} $(-a)(-b) = ab$.
  \item\label{it:somme-produit}
        $\displaystyle\Bigl(\sum_{i=1}^{m} a_i\Bigr)\Bigl(\sum_{j=1}^{n} b_j\Bigr)
         = \sum_{i=1}^{m}\sum_{j=1}^{n} a_i b_j$.
\end{enumerate}
\end{proposition}

\begin{proof}
\ref{it:zero-abs}\; On a $0 \cdot a = (0 + 0) \cdot a = 0 \cdot a + 0 \cdot a$
par distributivité. En ajoutant $-(0 \cdot a)$ des deux côtés, on obtient
$0 = 0 \cdot a$. La relation $a \cdot 0 = 0$ se démontre de manière symétrique.

\ref{it:moins-un}\; On calcule $(-1) \cdot a + a = (-1) \cdot a + 1 \cdot a
= (-1 + 1) \cdot a = 0 \cdot a = 0$ par~\ref{it:zero-abs}. Donc
$(-1) \cdot a = -a$. L'égalité $a \cdot (-1) = -a$ s'obtient de façon analogue.

\ref{it:moins-moins}\; $(-a)(-b) = (-1 \cdot a)(-1 \cdot b) = (-1)(-1) \cdot ab$.
Or $(-1)(-1) + (-1) = (-1)((-1) + 1) = (-1) \cdot 0 = 0$, donc $(-1)(-1) = 1$.
D'où $(-a)(-b) = ab$.

\ref{it:somme-produit}\; Par récurrence sur $m$ et $n$, en utilisant la
distributivité à chaque étape.
\end{proof}

%% ----------------------------------------------------------
\section{Exemples fondamentaux}\label{sec:exemples-anneaux}
\index{anneau!exemples}

\begin{exemple}[Anneaux classiques]\label{ex:anneaux-classiques}
Les ensembles suivants, munis de l'addition et de la multiplication
usuelles, sont des anneaux commutatifs unitaires~:
\begin{enumerate}[label=\textup{(\alph*)}]
  \item $\Z$\index{$\Z$}, $\Q$\index{$\Q$}, $\R$\index{$\R$},
        $\C$\index{$\C$}.
  \item $\Z/n\Z$\index{$\Z/n\Z$} pour tout entier $n \geqslant 1$.
  \item L'anneau des entiers de Gauss\index{entiers de Gauss}
        $\Z[i] = \{a + bi \mid a, b \in \Z\}$.
  \item Pour $d \in \Z$ sans facteur carré,
        $\Z[\sqrt{d}] = \{a + b\sqrt{d} \mid a, b \in \Z\}$%
        \index{$\Z[\sqrt{d}]$}.
\end{enumerate}
\end{exemple}

\begin{exemple}[Anneau de polynômes]\label{ex:anneau-polynomes}
\index{anneau!de polynômes}\index{polynôme!anneau de}
Soit $A$ un anneau commutatif. L'\emph{anneau de polynômes}
$A[X]$\index{$A[X]$} est l'ensemble des suites $(a_n)_{n \geqslant 0}$
d'éléments de~$A$, presque toutes nulles, muni de l'addition terme à
terme et du produit de Cauchy~:
\[
  \Bigl(\sum_{k} a_k X^k\Bigr)\Bigl(\sum_{k} b_k X^k\Bigr)
  = \sum_{k} \Bigl(\sum_{i+j=k} a_i b_j\Bigr) X^k.
\]
On note $\deg(P)$\index{degré!d'un polynôme} le degré d'un polynôme
non nul~$P$. Par convention, $\deg(0) = -\infty$.
\end{exemple}

\begin{exemple}[Anneau de séries formelles]\label{ex:series-formelles}
\index{anneau!de séries formelles}\index{série formelle}
Soit $A$ un anneau commutatif. L'anneau $A[[X]]$\index{$A[[X]]$} des
\emph{séries formelles} à coefficients dans~$A$ est défini comme
l'ensemble des suites $(a_n)_{n \geqslant 0}$ d'éléments de~$A$
(sans condition de support fini), muni des mêmes opérations que $A[X]$.
L'anneau $A[X]$ est un sous-anneau de $A[[X]]$.
\end{exemple}

\begin{exemple}[Anneau de matrices]\label{ex:anneau-matrices}
\index{anneau!de matrices}\index{matrice!anneau de}
Soit $A$ un anneau et $n \geqslant 1$. L'ensemble $M_n(A)$%
\index{$M_n(A)$} des matrices carrées $n \times n$ à coefficients
dans~$A$, muni de l'addition et de la multiplication matricielles,
est un anneau unitaire d'élément unité $I_n$. Pour $n \geqslant 2$,
cet anneau est \emph{non commutatif}.
\end{exemple}

\begin{exemple}[Anneau de groupe]\label{ex:anneau-groupe}
\index{anneau!de groupe}
Soit $A$ un anneau commutatif et $G$ un groupe (fini pour simplifier).
L'\emph{anneau de groupe}\index{anneau!de groupe} $A[G]$ est le
$A$-module libre de base~$G$, c'est-à-dire l'ensemble des sommes
formelles $\sum_{g \in G} a_g \, g$ avec $a_g \in A$, muni de la
multiplication~:
\[
  \Bigl(\sum_{g \in G} a_g \, g\Bigr)\Bigl(\sum_{h \in G} b_h \, h\Bigr)
  = \sum_{g,h \in G} a_g b_h \, (gh).
\]
C'est un anneau unitaire (d'unité $1_A \cdot e_G$), commutatif si et
seulement si $G$ est abélien.
\end{exemple}

%% ----------------------------------------------------------
\section{Diviseurs de zéro, intégrité, éléments inversibles}%
\label{sec:diviseurs-zero}

\begin{definition}[Diviseur de zéro]\label{def:diviseur-zero}
\index{diviseur de zéro}
Soit $A$ un anneau. Un élément $a \in A \setminus \{0\}$ est un
\emph{diviseur de zéro} (à gauche) s'il existe $b \in A \setminus \{0\}$
tel que $ab = 0$. On définit de même les diviseurs de zéro à droite.
\end{definition}

\begin{definition}[Anneau intègre]\label{def:anneau-integre}
\index{anneau!intègre}\index{domaine d'intégrité|see{anneau intègre}}
Un anneau commutatif unitaire $A$ est \emph{intègre} (ou est un
\emph{domaine d'intégrité}) si $A \neq \{0\}$ et s'il ne possède
aucun diviseur de zéro, c'est-à-dire~:
\[
  \forall\, a, b \in A, \quad ab = 0 \implies a = 0 \text{ ou } b = 0.
\]
\end{definition}

\begin{exemple}\label{ex:integrite}
Les anneaux $\Z$, $\Q$, $\R$, $\C$, $K[X]$ (où $K$ est un corps),
$\Z[i]$ sont intègres. L'anneau $\Z/6\Z$ n'est pas intègre car
$\bar{2} \cdot \bar{3} = \bar{0}$.
\end{exemple}

\begin{definition}[Élément inversible, groupe des unités]%
\label{def:element-inversible}
\index{élément inversible}\index{unité|see{élément inversible}}
\index{groupe des unités}
Soit $A$ un anneau. Un élément $a \in A$ est \emph{inversible}
(ou est une \emph{unité}) s'il existe $b \in A$ tel que $ab = ba = 1$.
L'ensemble des éléments inversibles de~$A$ est noté $A^\times$ ou
$U(A)$\index{$U(A)$}~; c'est un groupe pour la multiplication,
appelé le \emph{groupe des unités} de~$A$.
\end{definition}

\begin{exemple}[Groupes des unités]\label{ex:groupes-unites}
\begin{enumerate}[label=\textup{(\alph*)}]
  \item $U(\Z) = \{-1, 1\}$.
  \item $U(\Z/n\Z) = \{\bar{a} \mid \pgcd(a,n) = 1\}
         = (\Z/n\Z)^\times$\index{$(\Z/n\Z)^\times$}.
  \item $U(K[X]) = K^\times = K \setminus \{0\}$ pour tout corps~$K$.
  \item $U(M_n(K)) = \mathrm{GL}_n(K)$\index{$\mathrm{GL}_n(K)$}.
  \item $U(\Z[i]) = \{1, -1, i, -i\}$.
  \item $U(K[[X]]) = \{f \in K[[X]] \mid f(0) \neq 0\}$.
\end{enumerate}
\end{exemple}

\begin{proposition}\label{prop:inversible-non-diviseur}
Soit $A$ un anneau et $a \in U(A)$. Alors $a$ n'est pas un diviseur
de zéro.
\end{proposition}

\begin{proof}
Supposons $ab = 0$ avec $a \in U(A)$. Alors
$b = 1 \cdot b = (a^{-1}a)b = a^{-1}(ab) = a^{-1} \cdot 0 = 0$.
Donc $a$ n'est pas un diviseur de zéro.
\end{proof}

%% ----------------------------------------------------------
\section{Sous-anneaux et morphismes d'anneaux}%
\label{sec:morphismes-anneaux}
\index{sous-anneau}\index{morphisme!d'anneaux}

\begin{definition}[Sous-anneau]\label{def:sous-anneau}
\index{sous-anneau}
Soit $A$ un anneau. Un sous-ensemble $B \subseteq A$ est un
\emph{sous-anneau} de~$A$ si~:
\begin{enumerate}[label=\textup{(\roman*)}]
  \item $1_A \in B$.
  \item $B$ est stable par soustraction~: $\forall\, a, b \in B,\;
        a - b \in B$.
  \item $B$ est stable par multiplication~: $\forall\, a, b \in B,\;
        ab \in B$.
\end{enumerate}
\end{definition}

\begin{definition}[Morphisme d'anneaux]\label{def:morphisme-anneaux}
\index{morphisme!d'anneaux}
Soient $A$ et $B$ deux anneaux. Une application $\varphi \colon A \to B$
est un \emph{morphisme d'anneaux} si~:
\begin{enumerate}[label=\textup{(\roman*)}]
  \item $\varphi(a + a') = \varphi(a) + \varphi(a')$ pour tous
        $a, a' \in A$.
  \item $\varphi(a \cdot a') = \varphi(a) \cdot \varphi(a')$ pour tous
        $a, a' \in A$.
  \item $\varphi(1_A) = 1_B$.
\end{enumerate}
Le \emph{noyau}\index{noyau!d'un morphisme d'anneaux} est
$\Ker(\varphi) = \varphi^{-1}(\{0_B\})$ et l'\emph{image}\index{image!d'un morphisme d'anneaux}
est $\Im(\varphi) = \varphi(A)$.
\end{definition}

\begin{proposition}\label{prop:noyau-ideal}
Soit $\varphi \colon A \to B$ un morphisme d'anneaux.
\begin{enumerate}[label=\textup{(\roman*)}]
  \item $\Im(\varphi)$ est un sous-anneau de~$B$.
  \item $\Ker(\varphi)$ est un idéal bilatère de~$A$
        (cf.\ définition~\ref{def:ideal}).
  \item $\varphi$ est injectif si et seulement si
        $\Ker(\varphi) = \{0\}$.
\end{enumerate}
\end{proposition}

\begin{proof}
(i) L'image contient $\varphi(1_A) = 1_B$, et elle est stable par
soustraction et multiplication par fonctorialité de~$\varphi$.

(ii) Le noyau $I = \Ker(\varphi)$ est un sous-groupe de $(A,+)$ car
$\varphi$ est un morphisme de groupes additifs. Soient $a \in A$ et
$x \in I$. Alors $\varphi(ax) = \varphi(a)\varphi(x) = \varphi(a) \cdot 0 = 0$,
donc $ax \in I$. De même $xa \in I$. Ainsi $I$ est un idéal bilatère.

(iii) Si $\varphi$ est injectif, le seul antécédent de~$0$ est~$0$, donc
$\Ker(\varphi) = \{0\}$. Réciproquement, si $\Ker(\varphi) = \{0\}$
et $\varphi(a) = \varphi(a')$, alors $\varphi(a - a') = 0$, d'où
$a - a' \in \Ker(\varphi) = \{0\}$ et $a = a'$.
\end{proof}

%% ----------------------------------------------------------
\section{Idéaux}\label{sec:ideaux}
\index{idéal}

\begin{definition}[Idéal bilatère]\label{def:ideal}
\index{idéal!bilatère}\index{idéal!à gauche}\index{idéal!à droite}
Soit $A$ un anneau. Un sous-ensemble $I \subseteq A$ est un \emph{idéal
à gauche} de~$A$ si~:
\begin{enumerate}[label=\textup{(\roman*)}]
  \item $(I, +)$ est un sous-groupe de $(A, +)$.
  \item $\forall\, a \in A,\;\forall\, x \in I,\; ax \in I$.
\end{enumerate}
On définit de même un \emph{idéal à droite} en remplaçant~(ii) par
$xa \in I$. Un idéal est dit \emph{bilatère} s'il est à la fois à
gauche et à droite. Dans un anneau commutatif, ces trois notions
coïncident~; on parle alors simplement d'\emph{idéal}.
\end{definition}

\begin{remarque}\label{rem:ideal-critere}
En pratique, pour vérifier que $I$ est un idéal d'un anneau commutatif~$A$,
il suffit de montrer que~:
\begin{enumerate}[label=\textup{(\alph*)}]
  \item $I \neq \varnothing$ (par exemple $0 \in I$).
  \item $\forall\, x, y \in I,\; x - y \in I$.
  \item $\forall\, a \in A,\;\forall\, x \in I,\; ax \in I$.
\end{enumerate}
\end{remarque}

\begin{definition}[Idéal principal]\label{def:ideal-principal}
\index{idéal!principal}
Soit $A$ un anneau commutatif et $a \in A$. L'\emph{idéal principal
engendré par~$a$} est~:
\[
  (a) = aA = \{ax \mid x \in A\}.
\]
Un idéal $I$ de~$A$ est dit \emph{principal} s'il existe $a \in A$ tel
que $I = (a)$.
\end{definition}

\begin{definition}[Idéal engendré]\label{def:ideal-engendre}
\index{idéal!engendré}
Soit $A$ un anneau commutatif et $S \subseteq A$. L'\emph{idéal engendré
par~$S$}, noté $(S)$, est le plus petit idéal de~$A$ contenant~$S$~:
\[
  (S) = \bigcap_{\substack{I \text{ idéal de } A \\ S \subseteq I}} I
  = \Bigl\{\sum_{i=1}^{n} a_i s_i \;\Big|\; n \geqslant 0,\;
    a_i \in A,\; s_i \in S\Bigr\}.
\]
Pour un nombre fini de générateurs, on note
$(a_1, \ldots, a_k) = a_1 A + \cdots + a_k A$.
\end{definition}

\begin{exemple}[Idéaux de~$\Z$]\label{ex:ideaux-Z}
\index{idéal!de $\Z$}
Tout idéal de~$\Z$ est principal~: si $I$ est un idéal non nul de~$\Z$,
soit $d$ le plus petit élément strictement positif de~$I$. Par division
euclidienne, tout élément de~$I$ est un multiple de~$d$, d'où $I = d\Z = (d)$.
\end{exemple}

\begin{exemple}[Idéaux de~{$K[X]$}]\label{ex:ideaux-KX}
\index{idéal!de $K[X]$}
De même, tout idéal de~$K[X]$ ($K$ corps) est principal. Si $I \neq \{0\}$,
le polynôme unitaire de degré minimal dans~$I$ engendre~$I$ (par division
euclidienne).
\end{exemple}

%% ----------------------------------------------------------
\section{Opérations sur les idéaux}\label{sec:operations-ideaux}
\index{idéal!opérations}

\begin{proposition}[Opérations sur les idéaux]\label{prop:operations-ideaux}
Soient $I$ et $J$ deux idéaux d'un anneau commutatif~$A$.
\begin{enumerate}[label=\textup{(\roman*)}]
  \item L'\emph{intersection}\index{idéal!intersection} $I \cap J$ est
        un idéal de~$A$.
  \item La \emph{somme}\index{idéal!somme}
        $I + J = \{a + b \mid a \in I,\; b \in J\}$ est un idéal de~$A$~;
        c'est le plus petit idéal contenant $I$ et~$J$.
  \item Le \emph{produit}\index{idéal!produit}
        $IJ = \bigl\{\sum_{k=1}^{n} a_k b_k \mid n \geqslant 1,\;
        a_k \in I,\; b_k \in J\bigr\}$ est un idéal de~$A$, et
        $IJ \subseteq I \cap J$.
\end{enumerate}
\end{proposition}

\begin{proof}
(i) L'intersection de sous-groupes est un sous-groupe, et si $a \in A$,
$x \in I \cap J$, alors $ax \in I$ et $ax \in J$, donc $ax \in I \cap J$.

(ii) La somme $I + J$ contient~$0 = 0 + 0$. Si $a + b,\, a' + b' \in I + J$
avec $a, a' \in I$ et $b, b' \in J$, alors $(a + b) - (a' + b') =
(a - a') + (b - b') \in I + J$. Pour $c \in A$,
$c(a + b) = ca + cb \in I + J$.

(iii) Vérifions que $IJ$ est un idéal. La stabilité par addition
résulte de la définition (on concatène les sommes). Pour $c \in A$ et
$\sum a_k b_k \in IJ$, on a $c \sum a_k b_k = \sum (ca_k) b_k \in IJ$
car $ca_k \in I$. Enfin, si $x \in IJ$, $x = \sum a_k b_k$ avec
$a_k b_k \in I$ (car $I$ est un idéal et $b_k \in A$) et
$a_k b_k \in J$ de même, donc $x \in I \cap J$.
\end{proof}

\begin{exemple}\label{ex:operations-ideaux-Z}
Dans $\Z$~: $(m) + (n) = (\pgcd(m,n))$, \quad $(m) \cap (n) = (\ppcm(m,n))$,
\quad $(m)(n) = (mn)$.
\end{exemple}

%% ----------------------------------------------------------
\section{Idéaux premiers et idéaux maximaux}\label{sec:premiers-maximaux}
\index{idéal!premier}\index{idéal!maximal}

\begin{definition}[Idéal premier]\label{def:ideal-premier}
\index{idéal!premier}
Un idéal $\mathfrak{p}$ d'un anneau commutatif~$A$ est \emph{premier}
si $\mathfrak{p} \neq A$ et~:
\[
  \forall\, a, b \in A, \quad ab \in \mathfrak{p}
  \implies a \in \mathfrak{p} \text{ ou } b \in \mathfrak{p}.
\]
\end{definition}

\begin{definition}[Idéal maximal]\label{def:ideal-maximal}
\index{idéal!maximal}
Un idéal $\mathfrak{m}$ d'un anneau commutatif~$A$ est \emph{maximal}
si $\mathfrak{m} \neq A$ et si les seuls idéaux contenant $\mathfrak{m}$
sont $\mathfrak{m}$ et~$A$~:
\[
  \forall\, I \text{ idéal de } A, \quad
  \mathfrak{m} \subseteq I \subseteq A
  \implies I = \mathfrak{m} \text{ ou } I = A.
\]
\end{definition}

\begin{theoreme}[Caractérisation des idéaux premiers]%
\label{thm:caract-premier}
\index{idéal!premier!caractérisation}
Soit $A$ un anneau commutatif et $\mathfrak{p}$ un idéal de~$A$.
Alors~:
\[
  \mathfrak{p} \text{ est premier}
  \iff
  A/\mathfrak{p} \text{ est intègre}.
\]
\end{theoreme}

\begin{proof}
$(\Rightarrow)$\; Supposons $\mathfrak{p}$ premier. Comme
$\mathfrak{p} \neq A$, on a $A/\mathfrak{p} \neq \{0\}$. Soient
$\bar{a}, \bar{b} \in A/\mathfrak{p}$ tels que $\bar{a}\bar{b} = \bar{0}$,
c'est-à-dire $ab \in \mathfrak{p}$. Comme $\mathfrak{p}$ est premier,
$a \in \mathfrak{p}$ ou $b \in \mathfrak{p}$, c'est-à-dire $\bar{a} = \bar{0}$
ou $\bar{b} = \bar{0}$. Donc $A/\mathfrak{p}$ est intègre.

$(\Leftarrow)$\; Supposons $A/\mathfrak{p}$ intègre. En particulier
$A/\mathfrak{p} \neq \{0\}$, donc $\mathfrak{p} \neq A$. Si
$ab \in \mathfrak{p}$, alors $\bar{a}\bar{b} = \overline{ab} = \bar{0}$
dans $A/\mathfrak{p}$. Par intégrité, $\bar{a} = \bar{0}$ ou
$\bar{b} = \bar{0}$, c'est-à-dire $a \in \mathfrak{p}$ ou
$b \in \mathfrak{p}$. Donc $\mathfrak{p}$ est premier.
\end{proof}

\begin{theoreme}[Caractérisation des idéaux maximaux]%
\label{thm:caract-maximal}
\index{idéal!maximal!caractérisation}
Soit $A$ un anneau commutatif et $\mathfrak{m}$ un idéal de~$A$.
Alors~:
\[
  \mathfrak{m} \text{ est maximal}
  \iff
  A/\mathfrak{m} \text{ est un corps}.
\]
\end{theoreme}

\begin{proof}
$(\Rightarrow)$\; Supposons $\mathfrak{m}$ maximal. Comme
$\mathfrak{m} \neq A$, l'anneau quotient $A/\mathfrak{m}$ est non nul.
Soit $\bar{a} \in A/\mathfrak{m}$, $\bar{a} \neq \bar{0}$, c'est-à-dire
$a \notin \mathfrak{m}$. Considérons l'idéal $\mathfrak{m} + (a)$.
Comme $a \notin \mathfrak{m}$, on a
$\mathfrak{m} \subsetneq \mathfrak{m} + (a)$. Par maximalité,
$\mathfrak{m} + (a) = A$. Il existe donc $m \in \mathfrak{m}$ et $b \in A$
tels que $m + ab = 1$. En passant au quotient, $\bar{a}\bar{b} = \bar{1}$,
ce qui montre que $\bar{a}$ est inversible. Ainsi $A/\mathfrak{m}$ est
un corps.

$(\Leftarrow)$\; Supposons que $A/\mathfrak{m}$ est un corps. En
particulier $A/\mathfrak{m} \neq \{0\}$, donc $\mathfrak{m} \neq A$.
Soit $I$ un idéal de~$A$ tel que $\mathfrak{m} \subsetneq I$. Il existe
$a \in I \setminus \mathfrak{m}$, donc $\bar{a} \neq \bar{0}$ dans
$A/\mathfrak{m}$. Comme $A/\mathfrak{m}$ est un corps, il existe
$\bar{b}$ tel que $\bar{a}\bar{b} = \bar{1}$, c'est-à-dire
$ab - 1 \in \mathfrak{m} \subseteq I$. Comme $a \in I$ et $ab \in I$,
on obtient $1 = ab - (ab - 1) \in I$, d'où $I = A$. Ainsi $\mathfrak{m}$
est maximal.
\end{proof}

\begin{corollaire}[Maximal implique premier]\label{cor:maximal-premier}
\index{idéal!maximal implique premier}
Tout idéal maximal est premier.
\end{corollaire}

\begin{proof}
Si $\mathfrak{m}$ est maximal, alors $A/\mathfrak{m}$ est un corps par
le théorème~\ref{thm:caract-maximal}. Or tout corps est intègre (un
élément non nul est inversible, donc n'est pas diviseur de zéro par la
proposition~\ref{prop:inversible-non-diviseur}). Donc $A/\mathfrak{m}$
est intègre, et $\mathfrak{m}$ est premier par le
théorème~\ref{thm:caract-premier}.
\end{proof}

\begin{remarque}\label{rem:premier-pas-maximal}
La réciproque est fausse en général~: l'idéal $(0)$ de~$\Z$ est premier
(car $\Z$ est intègre) mais n'est pas maximal (il est strictement contenu
dans~$(2)$, par exemple).
\end{remarque}

\begin{exemple}[Idéaux de $\Z/n\Z$]\label{ex:ideaux-ZnZ}
Les idéaux de $\Z/n\Z$ sont de la forme $d\Z/n\Z$ où $d \mid n$.
L'idéal $d\Z/n\Z$ est premier (resp.\ maximal) si et seulement si $d$ est
un nombre premier divisant~$n$ (auquel cas il est aussi maximal, car
$(\Z/n\Z)/(d\Z/n\Z) \simeq \Z/d\Z$ est un corps).
\end{exemple}

%% ----------------------------------------------------------
\section{Treillis des idéaux de $\Z/12\Z$}\label{sec:treillis-Z12Z}

\begin{figure}[ht]
\centering
\begin{tikzpicture}[
  node distance=1.5cm,
  every node/.style={draw, circle, inner sep=2pt, minimum size=6mm,
                     font=\small}
]
  \node (Z12) at (0,4.5) {$\Z/12\Z$};
  \node (2)   at (-2,3)  {$(2)$};
  \node (3)   at (2,3)   {$(3)$};
  \node (4)   at (-3,1.5){$(4)$};
  \node (6)   at (0,1.5) {$(6)$};
  \node (0)   at (0,0)   {$(0)$};

  \draw (Z12) -- (2);
  \draw (Z12) -- (3);
  \draw (2)   -- (4);
  \draw (2)   -- (6);
  \draw (3)   -- (6);
  \draw (4)   -- (0);
  \draw (6)   -- (0);
\end{tikzpicture}
\caption{Treillis des idéaux de $\Z/12\Z$. Les idéaux maximaux
sont $(2)$ et $(3)$.}\label{fig:treillis-Z12Z}
\end{figure}

%% ----------------------------------------------------------
\section{Exercices}\label{sec:exo-anneaux}

\begin{exercice}\label{exo:ann-1}
Montrer que dans un anneau~$A$, l'élément unité est unique.
\end{exercice}

\begin{exercice}\label{exo:ann-2}
Montrer que $\Z[\sqrt{2}] = \{a + b\sqrt{2} \mid a, b \in \Z\}$ est un
sous-anneau de~$\R$. Déterminer $U(\Z[\sqrt{2}])$.
\end{exercice}

\begin{exercice}\label{exo:ann-3}
Soit $A$ un anneau commutatif. Montrer que $A$ est intègre si et seulement
si pour tous $a, b, c \in A$ avec $a \neq 0$, l'égalité $ab = ac$ entraîne
$b = c$ (propriété de simplification).
\end{exercice}

\begin{exercice}\label{exo:ann-4}
Trouver tous les diviseurs de zéro dans $\Z/12\Z$.
\end{exercice}

\begin{exercice}\label{exo:ann-5}
Soit $A$ un anneau fini intègre. Montrer que $A$ est un corps.
\emph{Indication~: considérer l'application $x \mapsto ax$ pour
$a \neq 0$.}
\end{exercice}

\begin{exercice}\label{exo:ann-6}
\index{nilpotent}
Un élément $a$ d'un anneau~$A$ est dit \emph{nilpotent} s'il existe
$n \geqslant 1$ tel que $a^n = 0$. Montrer que si $a$ est nilpotent,
alors $1 - a$ est inversible. \emph{Indication~: série géométrique finie.}
\end{exercice}

\begin{exercice}\label{exo:ann-7}
Soit $\varphi \colon A \to B$ un morphisme d'anneaux. Montrer que si
$\mathfrak{p}$ est un idéal premier de~$B$, alors
$\varphi^{-1}(\mathfrak{p})$ est un idéal premier de~$A$.
\end{exercice}

\begin{exercice}\label{exo:ann-8}
Montrer que l'intersection d'une famille quelconque d'idéaux est un idéal.
L'union de deux idéaux est-elle toujours un idéal~?
\end{exercice}

\begin{exercice}\label{exo:ann-9}
Déterminer tous les idéaux premiers et maximaux de $\Z/30\Z$.
\end{exercice}

\begin{exercice}\label{exo:ann-10}
Soit $A$ un anneau commutatif. Montrer que l'ensemble des éléments
nilpotents de~$A$, noté $\mathfrak{N}(A)$\index{nilradical}, est un
idéal de~$A$ (le \emph{nilradical}).
\end{exercice}

\begin{exercice}\label{exo:ann-11}
Soit $K$ un corps. Montrer que $K[X]/(X^2)$ est un anneau local (un seul
idéal maximal). Déterminer les éléments inversibles et les éléments
nilpotents.
\end{exercice}

%% ----------------------------------------------------------
\section*{Résumé du chapitre~\thechapter}

\begin{itemize}
  \item Un \emph{anneau} $(A,+,\cdot)$ est un groupe abélien pour
        l'addition, muni d'une multiplication associative, unitaire et
        distributive.
  \item Un \emph{anneau intègre} est un anneau commutatif non nul sans
        diviseurs de zéro. Le \emph{groupe des unités} $U(A)$ est le
        groupe des éléments inversibles.
  \item Un \emph{idéal} $I$ d'un anneau commutatif~$A$ est un
        sous-groupe additif stable par multiplication par les éléments
        de~$A$. Un idéal \emph{principal} est engendré par un seul
        élément.
  \item Un idéal $\mathfrak{p}$ est \emph{premier} si et seulement si
        $A/\mathfrak{p}$ est intègre. Un idéal $\mathfrak{m}$ est
        \emph{maximal} si et seulement si $A/\mathfrak{m}$ est un corps.
  \item Tout idéal maximal est premier. La réciproque est fausse en
        général.
\end{itemize}


%% ==========================================================
%%  CHAPITRE 6 : Anneaux quotients, domaines d'intégrité,
%%               corps de fractions
%% ==========================================================

\chapter{Anneaux quotients, domaines d'intégrité, corps de fractions}%
\label{chap:quotients-fractions}
\index{anneau quotient}\index{corps de fractions}

\section*{Introduction}

Le passage au quotient, que nous avons rencontré dans le cadre des
groupes, se transpose naturellement aux anneaux~: pour former l'anneau
quotient $A/I$, il faut et il suffit que $I$ soit un idéal bilatère.
Ce chapitre développe cette construction, établit les théorèmes
d'isomorphisme pour les anneaux, puis étudie le théorème des restes
chinois. Nous examinons ensuite la notion de caractéristique et
construisons, pour tout anneau intègre, son corps de fractions --- la
plus petite extension dans laquelle tout élément non nul devient
inversible.

%% ----------------------------------------------------------
\section{Anneau quotient}\label{sec:anneau-quotient}
\index{anneau quotient!construction}

\begin{theoreme}[Construction de l'anneau quotient]%
\label{thm:anneau-quotient}
Soit $A$ un anneau (commutatif) et $I$ un idéal de~$A$. L'ensemble
quotient $A/I = \{\bar{a} = a + I \mid a \in A\}$, muni des opérations~:
\[
  \bar{a} + \bar{b} = \overline{a + b}, \qquad
  \bar{a} \cdot \bar{b} = \overline{ab},
\]
est un anneau (commutatif), d'élément neutre additif $\bar{0}$ et
d'élément unité $\bar{1}$. La surjection canonique
$\pi \colon A \twoheadrightarrow A/I$, $a \mapsto \bar{a}$, est un
morphisme d'anneaux de noyau~$I$.
\end{theoreme}

\begin{proof}
Il faut d'abord vérifier que la multiplication est \emph{bien définie}.
Soient $a, a', b, b' \in A$ tels que $\bar{a} = \bar{a}'$ et
$\bar{b} = \bar{b}'$, c'est-à-dire $a - a' \in I$ et $b - b' \in I$.
Alors~:
\[
  ab - a'b' = ab - a'b + a'b - a'b' = (a - a')b + a'(b - b').
\]
Comme $I$ est un idéal, $(a - a')b \in I$ et $a'(b - b') \in I$, donc
$ab - a'b' \in I$, c'est-à-dire $\overline{ab} = \overline{a'b'}$.

La vérification des axiomes d'anneau pour $A/I$ est immédiate, car chaque
axiome dans $A/I$ se déduit de l'axiome correspondant dans~$A$ par passage
au quotient. Enfin, $\pi$ est clairement un morphisme d'anneaux surjectif,
et $\Ker(\pi) = \{a \in A \mid \bar{a} = \bar{0}\} = I$.
\end{proof}

\begin{theoreme}[Propriété universelle du quotient]%
\label{thm:prop-univ-quotient-anneaux}
\index{propriété universelle!anneau quotient}
Soient $A$ un anneau, $I$ un idéal de~$A$ et $\varphi \colon A \to B$
un morphisme d'anneaux tel que $I \subseteq \Ker(\varphi)$. Il existe
un unique morphisme d'anneaux $\bar{\varphi} \colon A/I \to B$ tel que
$\bar{\varphi} \circ \pi = \varphi$~:
\[
\begin{tikzcd}
  A \arrow[r, "\varphi"] \arrow[d, two heads, "\pi"'] & B \\
  A/I \arrow[ru, dashed, "\bar{\varphi}"']
\end{tikzcd}
\]
De plus, $\bar{\varphi}$ est injectif si et seulement si
$I = \Ker(\varphi)$, et surjectif si et seulement si $\varphi$ est
surjectif.
\end{theoreme}

\begin{proof}
Posons $\bar{\varphi}(\bar{a}) = \varphi(a)$. Cette définition est
licite~: si $\bar{a} = \bar{a}'$, alors $a - a' \in I \subseteq
\Ker(\varphi)$, donc $\varphi(a) = \varphi(a')$. La vérification que
$\bar{\varphi}$ est un morphisme d'anneaux est immédiate. L'unicité
découle de la surjectivité de~$\pi$.

On a $\Ker(\bar{\varphi}) = \{\bar{a} \mid \varphi(a) = 0\}
= \Ker(\varphi)/I$. Donc $\bar{\varphi}$ est injectif si et seulement si
$\Ker(\varphi) = I$. La surjectivité de~$\bar{\varphi}$ équivaut
clairement à celle de~$\varphi$.
\end{proof}

%% ----------------------------------------------------------
\section{Théorèmes d'isomorphisme pour les anneaux}%
\label{sec:iso-anneaux}
\index{théorème d'isomorphisme!pour les anneaux}

\begin{theoreme}[Premier théorème d'isomorphisme]%
\label{thm:premier-iso-anneaux}
\index{théorème d'isomorphisme!premier (anneaux)}
Soit $\varphi \colon A \to B$ un morphisme d'anneaux. Alors~:
\[
  A / \Ker(\varphi) \simeq \Im(\varphi).
\]
\end{theoreme}

\begin{proof}
On applique le théorème~\ref{thm:prop-univ-quotient-anneaux} avec
$I = \Ker(\varphi)$. Le morphisme induit
$\bar{\varphi} \colon A/\Ker(\varphi) \to B$ est injectif (car
$I = \Ker(\varphi)$) et son image est $\Im(\varphi)$.
\end{proof}

\begin{theoreme}[Deuxième théorème d'isomorphisme]%
\label{thm:deuxieme-iso-anneaux}
\index{théorème d'isomorphisme!deuxième (anneaux)}
Soit $A$ un anneau commutatif, $I$ et $J$ des idéaux de~$A$ avec
$I \subseteq J$. Alors $J/I$ est un idéal de $A/I$ et~:
\[
  (A/I)\big/(J/I) \simeq A/J.
\]
\end{theoreme}

\begin{proof}
Considérons le morphisme composé $A \xrightarrow{\pi_I} A/I
\xrightarrow{\pi_{J/I}} (A/I)/(J/I)$. C'est un morphisme d'anneaux
surjectif de noyau~$J$. Par le premier théorème d'isomorphisme,
$(A/I)/(J/I) \simeq A/J$.
\end{proof}

\begin{theoreme}[Troisième théorème d'isomorphisme]%
\label{thm:troisieme-iso-anneaux}
\index{théorème d'isomorphisme!troisième (anneaux)}
Soit $A$ un anneau commutatif, $I$ un idéal de~$A$ et $B$ un
sous-anneau de~$A$. Alors $I \cap B$ est un idéal de~$B$, $B + I$ est
un sous-anneau de~$A$ contenant~$I$ comme idéal, et~:
\[
  B / (I \cap B) \simeq (B + I) / I.
\]
\end{theoreme}

\begin{proof}
Le morphisme $B \hookrightarrow A \xrightarrow{\pi} A/I$ a pour noyau
$B \cap I$ et pour image $\{b + I \mid b \in B\} = (B + I)/I$. On
conclut par le premier théorème d'isomorphisme.
\end{proof}

%% ----------------------------------------------------------
\section{Théorème des restes chinois}\label{sec:CRT}
\index{théorème des restes chinois}

\begin{theoreme}[Théorème des restes chinois]\label{thm:CRT}
Soient $m, n \geqslant 1$ des entiers tels que $\pgcd(m,n) = 1$. Alors~:
\[
  \Z/mn\Z \simeq \Z/m\Z \times \Z/n\Z.
\]
\end{theoreme}

\begin{proof}
Considérons le morphisme d'anneaux~:
\[
  \varphi \colon \Z \longrightarrow \Z/m\Z \times \Z/n\Z, \qquad
  a \longmapsto (\bar{a}_m, \bar{a}_n),
\]
où $\bar{a}_m$ désigne la classe de~$a$ modulo~$m$.

\emph{Noyau.} On a $a \in \Ker(\varphi)$ si et seulement si $m \mid a$
et $n \mid a$. Comme $\pgcd(m,n) = 1$, cela équivaut à $mn \mid a$,
d'où $\Ker(\varphi) = mn\Z$.

\emph{Surjectivité.} Comme $\pgcd(m,n) = 1$, il existe $u, v \in \Z$
tels que $um + vn = 1$ (identité de Bézout\index{Bézout, identité de}).
Pour tout $(\bar{r}_m, \bar{s}_n) \in \Z/m\Z \times \Z/n\Z$, posons
$a = s \cdot um + r \cdot vn$. Alors~:
\[
  a \equiv r \cdot vn \equiv r \cdot (1 - um) \equiv r \pmod{m},
  \qquad
  a \equiv s \cdot um \equiv s \cdot (1 - vn) \equiv s \pmod{n}.
\]
Donc $\varphi(a) = (\bar{r}_m, \bar{s}_n)$, et $\varphi$ est surjectif.

Par le premier théorème d'isomorphisme
(théorème~\ref{thm:premier-iso-anneaux})~:
\[
  \Z/mn\Z = \Z/\Ker(\varphi) \simeq \Im(\varphi)
  = \Z/m\Z \times \Z/n\Z. \qedhere
\]
\end{proof}

\begin{theoreme}[Théorème des restes chinois --- version générale]%
\label{thm:CRT-general}
\index{théorème des restes chinois!version générale}
Soit $A$ un anneau commutatif et $I_1, \ldots, I_k$ des idéaux de~$A$
deux à deux \emph{comaximaux}\index{idéaux comaximaux}, c'est-à-dire
$I_i + I_j = A$ pour $i \neq j$. Alors~:
\[
  A\Big/\!\bigcap_{i=1}^{k} I_i \;\simeq\;
  \prod_{i=1}^{k} A/I_i,
\]
et de plus $\bigcap_{i=1}^{k} I_i = I_1 I_2 \cdots I_k$.
\end{theoreme}

\begin{proof}
Le morphisme canonique $\varphi \colon A \to \prod_{i=1}^{k} A/I_i$,
$a \mapsto (\bar{a}_1, \ldots, \bar{a}_k)$, a pour noyau
$\bigcap_{i} I_i$. Montrons la surjectivité par récurrence sur~$k$.

Pour $k = 2$~: comme $I_1 + I_2 = A$, il existe $a_1 \in I_1$,
$a_2 \in I_2$ avec $a_1 + a_2 = 1$. L'élément $e_1 = a_2$ satisfait
$e_1 \equiv 1 \pmod{I_1}$ et $e_1 \equiv 0 \pmod{I_2}$ (car
$e_1 = 1 - a_1$ et $a_1 \in I_1$; et $e_1 = a_2 \in I_2$).
Corrigeons~: $e_1 = a_2 \in I_2$ donc $e_1 \equiv 0 \pmod{I_2}$,
et $e_1 = 1 - a_1$ avec $a_1 \in I_1$ donc $e_1 \equiv 1 \pmod{I_1}$.
De même $e_2 = a_1$ satisfait $e_2 \equiv 0 \pmod{I_1}$ et
$e_2 \equiv 1 \pmod{I_2}$. Pour $(r_1, r_2) \in A/I_1 \times A/I_2$,
l'élément $x = r_1 e_1 + r_2 e_2$ satisfait $\varphi(x) = (r_1, r_2)$.

Montrons aussi que $I_1 \cap I_2 = I_1 I_2$. On a toujours
$I_1 I_2 \subseteq I_1 \cap I_2$. Réciproquement, si
$x \in I_1 \cap I_2$, alors $x = x \cdot 1 = x(a_1 + a_2)
= xa_1 + xa_2$. Or $xa_1 \in I_2 I_1$ (car $x \in I_2$) et
$xa_2 \in I_1 I_2$ (car $x \in I_1$), donc $x \in I_1 I_2$.

Le cas général $k \geqslant 3$ se traite par récurrence en montrant que
$I_1$ et $I_2 \cdots I_k$ sont comaximaux (exercice).
\end{proof}

%% ----------------------------------------------------------
\section{Anneau intègre et caractéristique}\label{sec:integre-caract}
\index{anneau!intègre}\index{caractéristique}

\begin{proposition}[Propriété de simplification]\label{prop:simplification}
\index{simplification (propriété de)}
Soit $A$ un anneau commutatif. Alors $A$ est intègre si et seulement si
pour tous $a, b, c \in A$ avec $a \neq 0$~:
\[
  ab = ac \implies b = c.
\]
\end{proposition}

\begin{proof}
$(\Rightarrow)$\; $ab = ac \implies a(b - c) = 0 \implies b - c = 0$
(car $a \neq 0$ et $A$ intègre).

$(\Leftarrow)$\; Si $ab = 0$ avec $a \neq 0$, alors $ab = a \cdot 0$,
d'où $b = 0$ par simplification.
\end{proof}

\begin{definition}[Caractéristique]\label{def:caracteristique}
\index{caractéristique!d'un anneau}
Soit $A$ un anneau. La \emph{caractéristique} de~$A$, notée
$\car(A)$\index{$\car(A)$}, est l'ordre de~$1_A$ dans le groupe additif
$(A,+)$~: c'est le plus petit entier $n \geqslant 1$ tel que
$n \cdot 1_A = 0$, ou $0$ si aucun tel entier n'existe.

De manière équivalente, l'unique morphisme d'anneaux $\Z \to A$,
$n \mapsto n \cdot 1_A$, a pour noyau $\car(A) \cdot \Z$.
\end{definition}

\begin{proposition}\label{prop:caract-integre}
La caractéristique d'un anneau intègre est $0$ ou un nombre premier.
\end{proposition}

\begin{proof}
Soit $A$ intègre avec $\car(A) = n > 0$. Supposons $n = ab$ avec
$1 < a, b < n$. Alors $0 = n \cdot 1_A = (a \cdot 1_A)(b \cdot 1_A)$
dans~$A$. Par intégrité, $a \cdot 1_A = 0$ ou $b \cdot 1_A = 0$, ce
qui contredit la minimalité de~$n$. Donc $n$ est premier.
\end{proof}

%% ----------------------------------------------------------
\section{Corps de fractions}\label{sec:corps-fractions}
\index{corps de fractions}

\subsection{Construction}\label{subsec:construction-fractions}

\begin{theoreme}[Corps de fractions]\label{thm:corps-fractions}
\index{corps de fractions!construction}
Soit $A$ un anneau commutatif intègre. Il existe un corps $K$, unique à
isomorphisme unique près, et un morphisme injectif d'anneaux
$\iota \colon A \hookrightarrow K$, tels que tout élément de~$K$ s'écrive
$\iota(a) \cdot \iota(b)^{-1}$ avec $a \in A$, $b \in A \setminus \{0\}$.
Le corps $K$ est appelé le \emph{corps de fractions} de~$A$ et noté
$\mathrm{Frac}(A)$.
\end{theoreme}

\begin{proof}
\emph{Construction.} Posons $S = A \times (A \setminus \{0\})$ et
définissons sur~$S$ la relation~:
\[
  (a,b) \sim (c,d) \iff ad = bc.
\]

\emph{\'Etape~1~: $\sim$ est une relation d'équivalence.}
La réflexivité ($ab = ba$) et la symétrie sont claires. Pour la
transitivité, supposons $(a,b) \sim (c,d)$ et $(c,d) \sim (e,f)$,
c'est-à-dire $ad = bc$ et $cf = de$. Alors~:
\[
  adf = bcf = bde, \quad\text{d'où}\quad d(af - be) = 0.
\]
Comme $d \neq 0$ et $A$ intègre, $af = be$, c'est-à-dire
$(a,b) \sim (e,f)$.

Notons $\dfrac{a}{b}$ la classe de $(a,b)$ et $K = S/{\sim}$.

\emph{\'Etape~2~: opérations sur~$K$.} Définissons~:
\[
  \frac{a}{b} + \frac{c}{d} = \frac{ad + bc}{bd}, \qquad
  \frac{a}{b} \cdot \frac{c}{d} = \frac{ac}{bd}.
\]
Notons que $bd \neq 0$ par intégrité de~$A$.

\emph{\'Etape~3~: bonne définition.}
Vérifions pour la multiplication. Soient $(a,b) \sim (a',b')$ et
$(c,d) \sim (c',d')$, c'est-à-dire $ab' = a'b$ et $cd' = c'd$.
Alors~:
\[
  (ac)(b'd') = (ab')(cd') = (a'b)(c'd) = (a'c')(bd),
\]
donc $(ac, bd) \sim (a'c', b'd')$. La vérification pour l'addition
est analogue~: $(ad + bc)(b'd') = adb'd' + bcb'd'
= a'dbd' + b \cdot c'd \cdot b' \cdot 1$. Développons plus
soigneusement~:
\begin{align*}
  (ad + bc)(b'd') &= adb'd' + bcb'd' \\
  (a'd' + b'c')(bd) &= a'd'bd + b'c'bd.
\end{align*}
Or $ab' = a'b$ donne $adb'd' = a'bdd' = a'd'bd$, et $cd' = c'd$
donne $bcb'd' = bb'cd' = bb'c'd = b'c'bd$. Donc les deux expressions
sont égales.

\emph{\'Etape~4~: $K$ est un corps.}
L'élément neutre additif est $\frac{0}{1}$, l'opposé de $\frac{a}{b}$
est $\frac{-a}{b}$, et l'élément unité est $\frac{1}{1}$. La
vérification de l'associativité, de la commutativité et de la
distributivité est directe (et fastidieuse). Si $\frac{a}{b} \neq
\frac{0}{1}$, c'est-à-dire $a \neq 0$, alors $\frac{b}{a} \in K$ et
$\frac{a}{b} \cdot \frac{b}{a} = \frac{ab}{ba} = \frac{1}{1}$. Donc
tout élément non nul est inversible~: $K$ est un corps.

\emph{\'Etape~5~: plongement de~$A$ dans~$K$.}
Le morphisme $\iota \colon A \to K$, $a \mapsto \frac{a}{1}$, est
un morphisme d'anneaux injectif. En effet, $\frac{a}{1} = \frac{0}{1}$
si et seulement si $a \cdot 1 = 0 \cdot 1$, c'est-à-dire $a = 0$.
\end{proof}

\begin{proposition}[Propriété universelle du corps de fractions]%
\label{prop:prop-univ-fractions}
\index{corps de fractions!propriété universelle}
Soit $A$ un anneau intègre et $\iota \colon A \hookrightarrow
\mathrm{Frac}(A)$ le plongement canonique. Pour tout morphisme
injectif d'anneaux $\varphi \colon A \hookrightarrow K'$ où $K'$ est
un corps, il existe un unique morphisme d'anneaux
$\tilde{\varphi} \colon \mathrm{Frac}(A) \to K'$ tel que
$\tilde{\varphi} \circ \iota = \varphi$~:
\[
\begin{tikzcd}
  A \arrow[r, hook, "\iota"] \arrow[d, hook, "\varphi"'] &
  \mathrm{Frac}(A) \arrow[dl, dashed, "\tilde{\varphi}"] \\
  K'
\end{tikzcd}
\]
\end{proposition}

\begin{proof}
Posons $\tilde{\varphi}\!\left(\frac{a}{b}\right) =
\varphi(a) \cdot \varphi(b)^{-1}$. Comme $\varphi$ est injectif et
$b \neq 0$, on a $\varphi(b) \neq 0$, donc $\varphi(b)^{-1}$ existe
dans le corps~$K'$.

\emph{Bonne définition.} Si $\frac{a}{b} = \frac{c}{d}$, alors
$ad = bc$, d'où $\varphi(a)\varphi(d) = \varphi(b)\varphi(c)$, d'où
$\varphi(a)\varphi(b)^{-1} = \varphi(c)\varphi(d)^{-1}$ (en multipliant
à droite par $\varphi(d)^{-1}$ et à gauche par $\varphi(b)^{-1}$).

La vérification que $\tilde{\varphi}$ est un morphisme d'anneaux est
directe. L'unicité résulte du fait que tout élément de $\mathrm{Frac}(A)$
s'écrit $\iota(a)\iota(b)^{-1}$, et un morphisme d'anneaux préserve les
inverses.
\end{proof}

\begin{exemple}[Corps de fractions classiques]\label{ex:corps-fractions}
\begin{enumerate}[label=\textup{(\alph*)}]
  \item $\mathrm{Frac}(\Z) = \Q$.
  \item $\mathrm{Frac}(K[X]) = K(X)$\index{$K(X)$}, le corps des
        fractions rationnelles à coefficients dans~$K$.
  \item $\mathrm{Frac}(\Z[i]) = \Q[i] = \{a + bi \mid a, b \in \Q\}
         = \Q(i)$.
\end{enumerate}
\end{exemple}

%% ----------------------------------------------------------
\section{Exercices}\label{sec:exo-quotients}

\begin{exercice}\label{exo:quot-1}
Montrer que $\Z[i]/(1+i) \simeq \Z/2\Z$.
\emph{Indication~: utiliser le morphisme $\Z \to \Z[i]/(1+i)$.}
\end{exercice}

\begin{exercice}\label{exo:quot-2}
Montrer que si $p$ est premier, l'anneau $\Z/p\Z$ n'a aucun idéal non
trivial.
\end{exercice}

\begin{exercice}\label{exo:quot-3}
Soit $K$ un corps et $a \in K$. Montrer que l'application d'évaluation
$\mathrm{ev}_a \colon K[X] \to K$, $P \mapsto P(a)$, est un morphisme
d'anneaux surjectif de noyau $(X - a)$. En déduire que
$K[X]/(X-a) \simeq K$.
\end{exercice}

\begin{exercice}\label{exo:quot-4}
Montrer que $\R[X]/(X^2 + 1) \simeq \C$.
\end{exercice}

\begin{exercice}\label{exo:quot-5}
Appliquer le théorème des restes chinois pour décrire la structure de
$\Z/60\Z$ comme produit d'anneaux.
\end{exercice}

\begin{exercice}\label{exo:quot-6}
Soit $A$ un anneau intègre de caractéristique~$p > 0$. Montrer que
l'application de Frobenius\index{Frobenius, morphisme de}
$F \colon A \to A$, $a \mapsto a^p$, est un morphisme d'anneaux
injectif.
\end{exercice}

\begin{exercice}\label{exo:quot-7}
Construire le corps de fractions de $\Z[\sqrt{-3}]$ et le décrire
explicitement.
\end{exercice}

\begin{exercice}\label{exo:quot-8}
Montrer que dans un anneau intègre, la caractéristique est $0$ ou un
nombre premier (redemonstration directe, sans utiliser la
proposition~\ref{prop:caract-integre}).
\end{exercice}

\begin{exercice}\label{exo:quot-9}
Soient $I$, $J$ des idéaux d'un anneau commutatif~$A$ avec $I + J = A$.
Montrer que $I^m + J^n = A$ pour tous $m, n \geqslant 1$.
\end{exercice}

\begin{exercice}\label{exo:quot-10}
Soit $n = p_1^{\alpha_1} \cdots p_k^{\alpha_k}$ la décomposition de~$n$
en facteurs premiers. Montrer que~:
\[
  \Z/n\Z \simeq \Z/p_1^{\alpha_1}\Z \times \cdots
  \times \Z/p_k^{\alpha_k}\Z.
\]
En déduire l'expression de $\varphi(n)$ (indicatrice d'Euler).
\end{exercice}

\begin{exercice}\label{exo:quot-11}
Montrer que l'anneau $\Z[\sqrt{-5}]$ est intègre. Déterminer son corps
de fractions.
\end{exercice}

%% ----------------------------------------------------------
\section*{Résumé du chapitre~\thechapter}

\begin{itemize}
  \item L'anneau quotient $A/I$ est bien défini lorsque $I$ est un idéal,
        et la surjection canonique $\pi \colon A \to A/I$ vérifie une
        propriété universelle.
  \item Les théorèmes d'isomorphisme pour les anneaux sont les analogues
        exacts de ceux pour les groupes.
  \item Le théorème des restes chinois~: si $\pgcd(m,n) = 1$, alors
        $\Z/mn\Z \simeq \Z/m\Z \times \Z/n\Z$.
  \item La caractéristique d'un anneau intègre est $0$ ou un nombre
        premier.
  \item Pour tout anneau intègre~$A$, le corps de fractions
        $\mathrm{Frac}(A)$ est le plus petit corps contenant~$A$. Sa
        construction repose sur des classes d'équivalence de couples
        $(a,b)$ avec $b \neq 0$.
\end{itemize}


%% ==========================================================
%%  CHAPITRE 7 : Anneaux principaux, euclidiens et factoriels
%% ==========================================================

\chapter{Anneaux principaux, euclidiens et factoriels}%
\label{chap:principaux-factoriels}
\index{anneau principal}\index{anneau euclidien}\index{anneau factoriel}

\section*{Introduction}

L'un des résultats les plus fondamentaux de l'arithmétique élémentaire
est le \emph{théorème fondamental de l'arithmétique}~: tout entier
$n \geqslant 2$ se décompose de manière essentiellement unique en produit
de nombres premiers. Ce résultat, tenu pour évident pendant des siècles,
s'est révélé d'une subtilité inattendue lorsque les mathématiciens ont
cherché à l'étendre à d'autres anneaux. \textsc{Kummer}\index{Kummer, Ernst}
et \textsc{Dedekind}\index{Dedekind, Richard} ont découvert que la
factorisation unique peut \emph{échouer} dans certains anneaux d'entiers
algébriques, ce qui a conduit à l'invention des idéaux et,
ultimement, à la théorie que nous développons ici.

Ce chapitre classifie les anneaux intègres selon la qualité de leur
arithmétique~: anneaux euclidiens, principaux, factoriels. Nous
établissons les implications~:
\[
  \text{corps} \subset \text{euclidien} \subset \text{principal}
  \subset \text{factoriel} \subset \text{intègre}
\]
et démontrons que chaque inclusion est stricte.

%% ----------------------------------------------------------
\section{Divisibilité dans un anneau intègre}\label{sec:divisibilite}
\index{divisibilité}

\begin{definition}[Divisibilité, éléments associés]\label{def:divisibilite}
\index{divisibilité!dans un anneau intègre}\index{éléments associés}
Soit $A$ un anneau commutatif intègre et $a, b \in A$.
\begin{enumerate}[label=\textup{(\roman*)}]
  \item On dit que $a$ \emph{divise}~$b$ (noté $a \mid b$) s'il existe
        $c \in A$ tel que $b = ac$.
  \item Les éléments $a$ et $b$ sont \emph{associés} (noté $a \sim b$)
        s'il existe $u \in U(A)$ tel que $b = ua$, ce qui équivaut à
        $a \mid b$ et $b \mid a$.
\end{enumerate}
\end{definition}

\begin{remarque}\label{rem:divisibilite-ideaux}
Dans un anneau commutatif intègre~$A$, on a
$a \mid b \iff (b) \subseteq (a)$, et $a \sim b \iff (a) = (b)$.
\end{remarque}

\begin{definition}[Élément irréductible]\label{def:irreductible}
\index{irréductible (élément)}
Soit $A$ un anneau intègre. Un élément $p \in A$ est \emph{irréductible}
si~:
\begin{enumerate}[label=\textup{(\roman*)}]
  \item $p \neq 0$ et $p \notin U(A)$.
  \item Si $p = ab$, alors $a \in U(A)$ ou $b \in U(A)$.
\end{enumerate}
\end{definition}

\begin{definition}[Élément premier]\label{def:premier}
\index{premier (élément)}
Soit $A$ un anneau intègre. Un élément $p \in A$ est \emph{premier}
si~:
\begin{enumerate}[label=\textup{(\roman*)}]
  \item $p \neq 0$ et $p \notin U(A)$.
  \item Si $p \mid ab$, alors $p \mid a$ ou $p \mid b$.
\end{enumerate}
De manière équivalente, $p$ est premier si et seulement si l'idéal
$(p)$ est un idéal premier non nul de~$A$.
\end{definition}

\begin{lemme}[Premier implique irréductible]\label{lem:premier-irred}
\index{premier implique irréductible}
Dans tout anneau intègre, un élément premier est irréductible.
\end{lemme}

\begin{proof}
Soit $p$ premier dans un anneau intègre~$A$. Supposons $p = ab$. Alors
$p \mid ab$, donc $p \mid a$ ou $p \mid b$.

\emph{Cas~1}~: $p \mid a$. Il existe $c \in A$ tel que $a = pc$. Alors
$p = ab = pcb$, d'où $p(1 - cb) = 0$. Comme $p \neq 0$ et $A$ intègre,
$cb = 1$, donc $b \in U(A)$.

\emph{Cas~2}~: $p \mid b$. Par un argument symétrique, $a \in U(A)$.

Dans les deux cas, $a \in U(A)$ ou $b \in U(A)$, donc $p$ est
irréductible.
\end{proof}

\begin{remarque}\label{rem:irred-pas-premier}
La réciproque est fausse en général. Dans
$\Z[\sqrt{-5}]$\index{$\Z[\sqrt{-5}]$}, l'élément~$2$ est irréductible
mais pas premier~: on a $2 \mid 6 = (1 + \sqrt{-5})(1 - \sqrt{-5})$,
mais $2 \nmid (1 + \sqrt{-5})$ et $2 \nmid (1 - \sqrt{-5})$.
\end{remarque}

%% ----------------------------------------------------------
\section{Anneaux principaux}\label{sec:anneaux-principaux}
\index{anneau principal}

\begin{definition}[Anneau principal]\label{def:anneau-principal}
\index{anneau principal!définition}
Un anneau commutatif intègre~$A$ est un \emph{anneau principal}
(ou un \emph{domaine d'idéaux principaux}, abrégé en PID\index{PID})
si tout idéal de~$A$ est principal, c'est-à-dire si pour tout
idéal~$I$ de~$A$, il existe $a \in A$ tel que $I = (a)$.
\end{definition}

\begin{exemple}\label{ex:anneaux-principaux}
Les anneaux $\Z$ et $K[X]$ (pour $K$ corps) sont des anneaux principaux
(exemples~\ref{ex:ideaux-Z} et~\ref{ex:ideaux-KX}).
\end{exemple}

\begin{theoreme}[Dans un PID, irréductible $\iff$ premier]%
\label{thm:irred-premier-PID}
\index{anneau principal!irréductible $\iff$ premier}
Soit $A$ un anneau principal et $p \in A$. Alors~:
\[
  p \text{ est irréductible} \iff p \text{ est premier}.
\]
\end{theoreme}

\begin{proof}
$(\Leftarrow)$\; C'est le lemme~\ref{lem:premier-irred}, valable dans
tout anneau intègre.

$(\Rightarrow)$\; Soit $p$ irréductible et supposons $p \mid ab$.
Comme $A$ est principal, l'idéal $(p, a)$ est principal~: il existe
$d \in A$ tel que $(p, a) = (d)$. En particulier $d \mid p$, donc
$p = du$ avec $u \in A$.

Comme $p$ est irréductible, $d \in U(A)$ ou $u \in U(A)$.

\emph{Cas~1}~: $u \in U(A)$, c'est-à-dire $d \sim p$. Alors
$(d) = (p)$, donc $a \in (p)$, c'est-à-dire $p \mid a$.

\emph{Cas~2}~: $d \in U(A)$. Alors $(d) = A$, donc $(p, a) = A$.
Il existe $x, y \in A$ tels que $px + ay = 1$. En multipliant par~$b$~:
$pbx + aby = b$. Or $p \mid ab$ donne $p \mid aby$, et $p \mid pbx$.
Donc $p \mid b$.

Dans les deux cas, $p \mid a$ ou $p \mid b$, donc $p$ est premier.
\end{proof}

%% ----------------------------------------------------------
\section{Anneaux euclidiens}\label{sec:anneaux-euclidiens}
\index{anneau euclidien}

\begin{definition}[Anneau euclidien]\label{def:anneau-euclidien}
\index{anneau euclidien!définition}\index{stathme}
Un anneau commutatif intègre~$A$ est \emph{euclidien} s'il existe une
application $v \colon A \setminus \{0\} \to \N$ (appelée
\emph{stathme}\index{stathme} ou \emph{fonction euclidienne}) telle que
pour tous $a \in A$ et $b \in A \setminus \{0\}$, il existe
$q, r \in A$ vérifiant~:
\[
  a = bq + r, \qquad
  r = 0 \text{ ou } v(r) < v(b).
\]
\end{definition}

\begin{theoreme}[Euclidien implique principal]%
\label{thm:euclidien-principal}
\index{anneau euclidien!implique principal}
Tout anneau euclidien est un anneau principal.
\end{theoreme}

\begin{proof}
Soit $A$ un anneau euclidien de stathme~$v$ et $I$ un idéal non nul
de~$A$. Choisissons $b \in I \setminus \{0\}$ tel que $v(b)$ soit
minimal parmi les valeurs de~$v$ sur $I \setminus \{0\}$. Montrons que
$I = (b)$.

L'inclusion $(b) \subseteq I$ est claire car $b \in I$ et $I$ est un
idéal. Réciproquement, soit $a \in I$. Par division euclidienne, il
existe $q, r \in A$ tels que $a = bq + r$ avec $r = 0$ ou $v(r) < v(b)$.
Or $r = a - bq \in I$ (car $a \in I$ et $bq \in I$). Si $r \neq 0$,
alors $v(r) < v(b)$, ce qui contredit la minimalité de~$v(b)$ dans
$I \setminus \{0\}$. Donc $r = 0$ et $a = bq \in (b)$.

Ainsi $I = (b)$ est principal.
\end{proof}

\begin{exemple}[Anneaux euclidiens classiques]\label{ex:anneaux-euclidiens}
\begin{enumerate}[label=\textup{(\alph*)}]
  \item $\Z$ est euclidien pour le stathme $v(n) = |n|$.
  \item $K[X]$ ($K$ corps) est euclidien pour le stathme
        $v(P) = \deg(P)$.
  \item Les entiers de Gauss $\Z[i]$ sont euclidiens pour le stathme
        $v(a + bi) = a^2 + b^2$ (la norme).
\end{enumerate}
\end{exemple}

\begin{theoreme}[{$\Z[i]$} est euclidien]\label{thm:Zi-euclidien}
\index{entiers de Gauss!anneau euclidien}
L'anneau $\Z[i]$ est euclidien pour le stathme
$N \colon \Z[i] \setminus \{0\} \to \N^*$ défini par
$N(a + bi) = a^2 + b^2$.
\end{theoreme}

\begin{proof}
Notons d'abord que $N$ est \emph{multiplicative}~:
$N(\alpha\beta) = N(\alpha)N(\beta)$ pour tous
$\alpha, \beta \in \Z[i]$ (vérification directe par calcul).

Soient $\alpha, \beta \in \Z[i]$ avec $\beta \neq 0$. Considérons le
quotient $\alpha/\beta \in \Q(i)$~:
\[
  \frac{\alpha}{\beta}
  = \frac{\alpha \bar{\beta}}{|\beta|^2}
  = u + vi
\]
avec $u, v \in \Q$. Choisissons $p, q \in \Z$ tels que $|u - p|
\leqslant \frac{1}{2}$ et $|v - q| \leqslant \frac{1}{2}$. Posons
$\gamma = p + qi \in \Z[i]$ et $\rho = \alpha - \beta\gamma \in \Z[i]$.
Alors~:
\[
  \frac{\rho}{\beta} = \frac{\alpha}{\beta} - \gamma
  = (u - p) + (v - q)i.
\]
Par conséquent~:
\[
  N(\rho) = N(\beta) \cdot N\!\left(\frac{\rho}{\beta}\right)
  = N(\beta) \bigl((u-p)^2 + (v-q)^2\bigr)
  \leqslant N(\beta) \left(\frac{1}{4} + \frac{1}{4}\right)
  = \frac{N(\beta)}{2} < N(\beta).
\]
Donc $\alpha = \beta\gamma + \rho$ avec $\rho = 0$ ou
$N(\rho) < N(\beta)$, ce qui montre que $\Z[i]$ est euclidien.
\end{proof}

%% ----------------------------------------------------------
\section{Échec de la factorisation unique~: $\Z[\sqrt{-5}]$}%
\label{sec:echec-factorisation}
\index{$\Z[\sqrt{-5}]$}\index{factorisation!échec}

\begin{proposition}\label{prop:Z-sqrt-5-non-UFD}
L'anneau $\Z[\sqrt{-5}]$ n'est pas factoriel.
\end{proposition}

\begin{proof}
Considérons la norme $N \colon \Z[\sqrt{-5}] \to \N$ définie par
$N(a + b\sqrt{-5}) = a^2 + 5b^2$. Cette norme est multiplicative~:
$N(\alpha\beta) = N(\alpha)N(\beta)$.

\emph{\'Etape~1~: les éléments $2$, $3$, $1 + \sqrt{-5}$,
$1 - \sqrt{-5}$ sont irréductibles.}

Pour~$2$~: $N(2) = 4$. Si $2 = \alpha\beta$, alors
$N(\alpha)N(\beta) = 4$. Les diviseurs possibles de~$4$ dans~$\N$
sont $1, 2, 4$. Or $N(a + b\sqrt{-5}) = 2$ exigerait $a^2 + 5b^2 = 2$,
ce qui n'a aucune solution dans~$\Z$ (si $b \neq 0$, $a^2 + 5b^2
\geqslant 5$~; si $b = 0$, $a^2 = 2$ est impossible dans~$\Z$). Donc
dans la décomposition $2 = \alpha\beta$, l'un des facteurs a norme~$1$
(donc est une unité, car $U(\Z[\sqrt{-5}]) = \{\pm 1\}$ puisque
$a^2 + 5b^2 = 1$ impose $b = 0$, $a = \pm 1$). Ainsi $2$ est
irréductible.

Le même argument s'applique à~$3$ ($N(3) = 9$, et $a^2 + 5b^2 = 3$ n'a
pas de solution) et à $1 \pm \sqrt{-5}$ ($N(1 \pm \sqrt{-5}) = 6$, et
$a^2 + 5b^2 = 2$ ou $3$ n'a pas de solution).

\emph{\'Etape~2~: deux factorisations distinctes de~$6$.}

On a~:
\[
  6 = 2 \cdot 3 = (1 + \sqrt{-5})(1 - \sqrt{-5}).
\]
Ce sont deux factorisations en irréductibles. Elles sont essentiellement
distinctes car $2$ n'est pas associé à $3$, $1 + \sqrt{-5}$ ou
$1 - \sqrt{-5}$ (les seules unités sont $\pm 1$, et par exemple
$2 \neq \pm(1 + \sqrt{-5})$).

Ainsi la factorisation n'est pas unique~: $\Z[\sqrt{-5}]$ n'est pas
factoriel.
\end{proof}

%% ----------------------------------------------------------
\section{Anneaux factoriels}\label{sec:anneaux-factoriels}
\index{anneau factoriel}

\begin{definition}[Anneau factoriel]\label{def:anneau-factoriel}
\index{anneau factoriel!définition}\index{UFD|see{anneau factoriel}}
Un anneau commutatif intègre~$A$ est \emph{factoriel} (ou est un
\emph{domaine de factorisation unique}, abrégé en UFD) si~:
\begin{enumerate}[label=\textup{(F\arabic*)}]
  \item\label{ax:F1} \textbf{Existence}~: tout élément
        $a \in A \setminus (\{0\} \cup U(A))$ s'écrit comme un produit
        fini d'éléments irréductibles~: $a = p_1 p_2 \cdots p_n$.
  \item\label{ax:F2} \textbf{Unicité}~: si $p_1 \cdots p_m
        = q_1 \cdots q_n$ sont deux factorisations en irréductibles,
        alors $m = n$ et, quitte à réordonner, $p_i \sim q_i$ pour tout~$i$.
\end{enumerate}
\end{definition}

\begin{theoreme}[PID implique UFD]\label{thm:PID-UFD}
\index{anneau principal!implique factoriel}
Tout anneau principal est factoriel.
\end{theoreme}

\begin{proof}
Soit $A$ un anneau principal. Nous devons démontrer l'existence et
l'unicité de la factorisation.

\medskip
\textbf{Existence de la factorisation.}
Supposons par l'absurde qu'il existe un élément
$a \in A \setminus (\{0\} \cup U(A))$ qui ne s'écrit pas comme produit
d'irréductibles. Comme $a$ n'est pas irréductible, on peut écrire
$a = a_1 b_1$ avec $a_1, b_1 \notin U(A)$ et $a_1, b_1 \neq 0$. Au
moins l'un des deux, disons~$a_1$, ne s'écrit pas comme produit
d'irréductibles (sinon $a = a_1 b_1$ serait un tel produit). De plus,
$(a) \subsetneq (a_1)$ car $a_1 \mid a$ mais $a_1 \nmid a$ impliquerait
$b_1 \in U(A)$, contradiction. Corrigeons~: on a $(a) \subsetneq (a_1)$
car $a = a_1 b_1 \in (a_1)$ donc $(a) \subseteq (a_1)$, et si
$(a) = (a_1)$, alors $a_1 = au$ pour un $u \in A$, d'où $a = a_1 b_1
= aub_1$, donc $ub_1 = 1$, c'est-à-dire $b_1 \in U(A)$, contradiction.

En itérant, on construit une chaîne strictement croissante d'idéaux~:
\[
  (a) \subsetneq (a_1) \subsetneq (a_2) \subsetneq \cdots
\]

Posons $J = \bigcup_{n \geqslant 0} (a_n)$ (avec $a_0 = a$). Alors $J$
est un idéal de~$A$~: si $x, y \in J$, il existe $n$ tel que
$x, y \in (a_n)$, donc $x - y \in (a_n) \subseteq J$~; et si $c \in A$,
$cx \in (a_n) \subseteq J$. Comme $A$ est principal, $J = (d)$ pour un
$d \in A$. Or $d \in J$, donc $d \in (a_N)$ pour un certain~$N$, d'où
$(d) \subseteq (a_N)$. Mais $(a_N) \subseteq J = (d)$, donc
$(a_N) = (d) = J$. Cela entraîne $(a_{N+1}) \subseteq J = (a_N)$,
ce qui contredit $(a_N) \subsetneq (a_{N+1})$.

Cette contradiction prouve que tout élément non nul et non inversible
admet une décomposition en produit d'irréductibles.

\medskip
\textbf{Unicité de la factorisation.}
Soient $p_1 \cdots p_m = q_1 \cdots q_n$ deux factorisations en
irréductibles. Comme $A$ est principal, tout irréductible est premier
(théorème~\ref{thm:irred-premier-PID}). Procédons par récurrence
sur~$m$.

\emph{Cas $m = 1$}~: $p_1 = q_1 \cdots q_n$. Comme $p_1$ est
irréductible, $n = 1$ et $p_1 = q_1$ (ils sont associés car le quotient
est une unité, mais ici l'égalité directe donne $p_1 \sim q_1$).

\emph{Pas de récurrence}~: supposons le résultat vrai pour $m - 1$.
L'élément $p_1$ divise $q_1 \cdots q_n$. Comme $p_1$ est premier, il
divise l'un des $q_j$. Quitte à réordonner, $p_1 \mid q_1$. Comme $q_1$
est irréductible et $p_1 \notin U(A)$, on a $q_1 = u p_1$ avec
$u \in U(A)$, c'est-à-dire $p_1 \sim q_1$.

Alors $p_1 p_2 \cdots p_m = u p_1 q_2 \cdots q_n$. Par simplification
(l'anneau est intègre), $p_2 \cdots p_m = u q_2 \cdots q_n$. En
remplaçant $q_2$ par $uq_2$ (qui est encore irréductible et associé
à~$q_2$), on obtient une factorisation $p_2 \cdots p_m = q_2' q_3
\cdots q_n$ en $m - 1$ et $n - 1$ irréductibles. Par hypothèse de
récurrence, $m - 1 = n - 1$ et, quitte à réordonner,
$p_i \sim q_i$ pour $i \geqslant 2$. Donc $m = n$ et
$p_i \sim q_i$ pour tout~$i$.
\end{proof}

%% ----------------------------------------------------------
\section{Lemme de Gauss et polynômes sur un UFD}%
\label{sec:gauss}
\index{Gauss, lemme de}\index{contenu d'un polynôme}

\begin{definition}[Contenu d'un polynôme]\label{def:contenu}
\index{contenu!d'un polynôme}
Soit $A$ un anneau factoriel et $f = a_0 + a_1 X + \cdots + a_n X^n
\in A[X]$ un polynôme non nul. Le \emph{contenu} de~$f$, noté $c(f)$,
est le pgcd des coefficients $a_0, a_1, \ldots, a_n$ (défini à
association près). Le polynôme $f$ est dit \emph{primitif}\index{primitif (polynôme)}
si $c(f) \sim 1$, c'est-à-dire si les coefficients de~$f$ sont
premiers entre eux dans leur ensemble.
\end{definition}

\begin{lemme}[Lemme de Gauss]\label{lem:gauss}
\index{Gauss, lemme de}
Soit $A$ un anneau factoriel et $f, g \in A[X]$ deux polynômes non nuls.
Alors~:
\[
  c(fg) \sim c(f) \cdot c(g).
\]
En particulier, le produit de deux polynômes primitifs est primitif.
\end{lemme}

\begin{proof}
En écrivant $f = c(f) f_0$ et $g = c(g) g_0$ avec $f_0$, $g_0$
primitifs, on a $fg = c(f)c(g) f_0 g_0$. Il suffit donc de montrer que
si $f$ et $g$ sont primitifs, alors $fg$ est primitif.

Supposons par l'absurde que $fg$ n'est pas primitif. Alors il existe un
élément irréductible $p \in A$ qui divise tous les coefficients de~$fg$.
Comme $A$ est factoriel (et en particulier un anneau intègre), l'idéal
$(p)$ est premier (car dans un UFD, tout irréductible est premier --- ce
que nous montrerons ci-dessous, ou utilisons directement le fait que $A$
est un UFD). Considérons la réduction modulo~$p$~:
$\bar{f}, \bar{g} \in (A/(p))[X]$, avec $\bar{f}\bar{g} = \bar{0}$.

Comme $p$ est premier, $A/(p)$ est intègre, donc $(A/(p))[X]$ est
intègre. Ainsi $\bar{f} = \bar{0}$ ou $\bar{g} = \bar{0}$,
c'est-à-dire $p$ divise tous les coefficients de~$f$ ou de~$g$. Cela
contredit le fait que $f$ et $g$ sont primitifs.
\end{proof}

\begin{remarque}\label{rem:irred-premier-UFD}
Dans un anneau factoriel, tout élément irréductible est premier. En
effet, soit $p$ irréductible et $p \mid ab$, disons $pc = ab$. En
décomposant $a$, $b$, $c$ en irréductibles et par unicité de la
factorisation, $p$ apparaît parmi les facteurs irréductibles de $ab$,
donc parmi ceux de~$a$ ou de~$b$, c'est-à-dire $p \mid a$ ou
$p \mid b$.
\end{remarque}

\begin{theoreme}[$A$ UFD implique {$A[X]$} UFD]\label{thm:AX-UFD}
\index{anneau factoriel!$A[X]$ factoriel}
Si $A$ est un anneau factoriel, alors $A[X]$ est un anneau factoriel.
\end{theoreme}

\begin{proof}[Esquisse de preuve]
Notons $K = \mathrm{Frac}(A)$. Comme $K$ est un corps, $K[X]$ est
euclidien, donc factoriel.

\emph{Existence.} Soit $f \in A[X]$ non nul, non inversible. On écrit
$f = c(f) \cdot f_0$ avec $f_0$ primitif. Le facteur $c(f) \in A$ se
décompose en irréductibles de~$A$ (qui restent irréductibles dans
$A[X]$). Pour~$f_0$, on le factorise dans~$K[X]$ en irréductibles~;
par le lemme de Gauss, on ramène ces facteurs à des polynômes primitifs
de~$A[X]$, qui sont irréductibles dans $A[X]$.

\emph{Unicité.} Si $p_1 \cdots p_m = q_1 \cdots q_n$ dans $A[X]$, on
sépare les facteurs constants (irréductibles de~$A$) des facteurs de
degré~$\geqslant 1$. L'unicité pour les facteurs constants provient de
celle dans~$A$~; pour les facteurs de degré~$\geqslant 1$, on passe
dans~$K[X]$ et on utilise l'unicité dans~$K[X]$ combinée au lemme de
Gauss.
\end{proof}

%% ----------------------------------------------------------
\section{Hiérarchie des anneaux}\label{sec:hierarchie}
\index{hiérarchie des anneaux}

Le diagramme suivant résume les inclusions entre les différentes classes
d'anneaux étudiées~:

\begin{figure}[ht]
\centering
\begin{tikzpicture}[
  box/.style={draw, rounded corners, minimum width=3.8cm,
              minimum height=0.9cm, font=\small, align=center},
  arrow/.style={-{Stealth}, thick},
  node distance=1.4cm
]
  \node[box] (corps)     {Corps};
  \node[box, below=of corps]     (eucl)  {Anneau euclidien};
  \node[box, below=of eucl]      (pid)   {Anneau principal (PID)};
  \node[box, below=of pid]       (ufd)   {Anneau factoriel (UFD)};
  \node[box, below=of ufd]       (int)   {Anneau intègre};

  \draw[arrow] (corps) -- (eucl)
    node[midway, right, font=\footnotesize] {trivial};
  \draw[arrow] (eucl)  -- (pid)
    node[midway, right, font=\footnotesize]
    {thm.~\ref{thm:euclidien-principal}};
  \draw[arrow] (pid)   -- (ufd)
    node[midway, right, font=\footnotesize]
    {thm.~\ref{thm:PID-UFD}};
  \draw[arrow] (ufd)   -- (int)
    node[midway, right, font=\footnotesize] {par définition};

  % Contre-exemples à gauche
  \node[font=\footnotesize, text width=4cm, align=right]
    at (-4.5, -1.4) {$\Z\bigl[\frac{1+\sqrt{-19}}{2}\bigr]$\\
    est PID, non euclidien};
  \node[font=\footnotesize, text width=4cm, align=right]
    at (-4.5, -4.2) {$\Z[X]$ est UFD,\\ non PID};
  \node[font=\footnotesize, text width=4cm, align=right]
    at (-4.5, -5.6) {$\Z[\sqrt{-5}]$ est intègre,\\ non UFD};
\end{tikzpicture}
\caption{Hiérarchie des anneaux commutatifs intègres. Chaque inclusion
est stricte.}\label{fig:hierarchie-anneaux}
\end{figure}

%% ----------------------------------------------------------
\section{Contre-exemple~: PID non euclidien}%
\label{sec:PID-non-euclidien}
\index{anneau principal!non euclidien}

\begin{theoreme}\label{thm:PID-non-euclidien}
L'anneau $\Z\!\left[\frac{1+\sqrt{-19}}{2}\right]$\index{$\Z[(1+\sqrt{-19})/2]$}
est un anneau principal qui n'est pas euclidien.
\end{theoreme}

\begin{remarque}\label{rem:PID-non-euclidien}
La preuve que cet anneau est principal mais non euclidien est
non triviale. La principalité découle de la théorie des formes
quadratiques binaires~: le nombre de classes de l'ordre
$\Z\!\left[\frac{1+\sqrt{-19}}{2}\right]$ (qui est l'anneau des entiers
de $\Q(\sqrt{-19})$) est~$1$, ce qui signifie que tout idéal est
principal. Pour la non-existence d'un stathme euclidien, on montre
qu'il n'existe pas de ``plus petit élément non inversible au sens de
tout stathme possible''. Nous renvoyons à \cite{DummitFoote} ou
\cite{Jacobson} pour les détails.
\end{remarque}

%% ----------------------------------------------------------
\section{Critère d'irréductibilité d'Eisenstein}%
\label{sec:eisenstein}
\index{Eisenstein, critère d'}

\begin{theoreme}[Critère d'Eisenstein]\label{thm:eisenstein}
\index{Eisenstein, critère d'}
Soit $A$ un anneau factoriel et
$f = a_n X^n + a_{n-1} X^{n-1} + \cdots + a_1 X + a_0 \in A[X]$
un polynôme de degré $n \geqslant 1$. S'il existe un élément
irréductible $p \in A$ tel que~:
\begin{enumerate}[label=\textup{(\roman*)}]
  \item $p \nmid a_n$,
  \item $p \mid a_i$ pour tout $0 \leqslant i \leqslant n-1$,
  \item $p^2 \nmid a_0$,
\end{enumerate}
alors $f$ est irréductible dans $A[X]$ (et dans $K[X]$ si $f$ est
primitif, où $K = \mathrm{Frac}(A)$).
\end{theoreme}

\begin{proof}
Supposons par l'absurde que $f = gh$ dans $A[X]$ avec
$\deg(g) = r \geqslant 1$ et $\deg(h) = s \geqslant 1$, $r + s = n$.
Écrivons $g = b_r X^r + \cdots + b_0$ et $h = c_s X^s + \cdots + c_0$.
Alors $a_0 = b_0 c_0$ et $a_n = b_r c_s$.

Réduisons modulo~$p$. Les conditions (i) et (ii) donnent
$\bar{f} = \bar{a}_n X^n$ dans $(A/(p))[X]$. Comme $p$ est irréductible
dans l'UFD~$A$, l'idéal $(p)$ est premier, et $A/(p)$ est intègre. De
$\bar{f} = \bar{g}\bar{h}$, on déduit $\bar{a}_n X^n = \bar{g} \cdot \bar{h}$
dans l'anneau intègre $(A/(p))[X]$.

Par unicité de la factorisation dans $(A/(p))[X]$ (c'est un domaine
intègre, et $X$ est irréductible --- en fait, on utilise simplement que
$(A/(p))[X]$ est intègre et que $X^n$ ne se factorise qu'en puissances
de~$X$ à unité près), on obtient $\bar{g} = \bar{b}_r X^r$ et
$\bar{h} = \bar{c}_s X^s$. En particulier, $\bar{b}_0 = 0$ et
$\bar{c}_0 = 0$, c'est-à-dire $p \mid b_0$ et $p \mid c_0$. Mais alors
$p^2 \mid b_0 c_0 = a_0$, ce qui contredit~(iii).
\end{proof}

\begin{exemple}\label{ex:eisenstein}
Le polynôme cyclotomique $\Phi_p(X) = X^{p-1} + X^{p-2} + \cdots + X + 1$
est irréductible sur~$\Z$ pour $p$ premier. En effet, le changement de
variable $Y = X - 1$ donne~:
\[
  \Phi_p(Y+1) = \frac{(Y+1)^p - 1}{Y}
  = Y^{p-1} + \binom{p}{1}Y^{p-2} + \cdots + \binom{p}{p-2}Y + p.
\]
Le critère d'Eisenstein\index{Eisenstein, critère d'} s'applique avec
l'irréductible~$p$~: le coefficient dominant est~$1$ ($p \nmid 1$),
tous les autres coefficients sont divisibles par~$p$
($p \mid \binom{p}{k}$ pour $1 \leqslant k \leqslant p-1$), et le terme
constant est~$p$ ($p^2 \nmid p$). Donc $\Phi_p(Y+1)$ est irréductible
dans $\Z[Y]$, et par conséquent $\Phi_p(X)$ est irréductible dans~$\Z[X]$.
\end{exemple}

%% ----------------------------------------------------------
\section{Exercices}\label{sec:exo-factoriels}

\begin{exercice}\label{exo:fact-1}
Déterminer les éléments inversibles et les éléments irréductibles
de~$\Z[i]$.
\end{exercice}

\begin{exercice}\label{exo:fact-2}
Montrer que $\Z[\sqrt{-3}]$ n'est pas factoriel.
\emph{Indication~: considérer $4 = 2 \cdot 2
= (1 + \sqrt{-3})(1 - \sqrt{-3})$.}
\end{exercice}

\begin{exercice}\label{exo:fact-3}
Montrer que dans un anneau principal, tout idéal premier non nul est
maximal.
\end{exercice}

\begin{exercice}\label{exo:fact-4}
Soit $K$ un corps. Montrer que $K[X,Y]$ n'est pas un anneau principal.
\emph{Indication~: considérer l'idéal $(X, Y)$.}
\end{exercice}

\begin{exercice}\label{exo:fact-5}
Appliquer le critère d'Eisenstein pour montrer que $X^4 + 2X^3 + 6X + 3$
est irréductible dans $\Z[X]$.
\end{exercice}

\begin{exercice}\label{exo:fact-6}
Soit $p$ un nombre premier et $p \equiv 1 \pmod{4}$. Montrer qu'il
existe $a, b \in \Z$ tels que $p = a^2 + b^2$.
\emph{Indication~: utiliser le fait que $-1$ est un carré modulo~$p$,
puis travailler dans $\Z[i]$ qui est principal.}
\end{exercice}

\begin{exercice}\label{exo:fact-7}
Montrer que $X^2 + Y^2 - 1$ est irréductible dans $\R[X,Y]$.
\end{exercice}

\begin{exercice}\label{exo:fact-8}
Soit $A$ un anneau principal et $a, b \in A$. Montrer qu'il existe
$d = \pgcd(a,b)$ et qu'on peut écrire $d = ua + vb$ pour certains
$u, v \in A$ (identité de Bézout).
\end{exercice}

\begin{exercice}\label{exo:fact-9}
Montrer que $\Z[X]$ est factoriel mais n'est pas principal.
\emph{Indication~: pour la non-principalité, considérer l'idéal
$(2, X)$.}
\end{exercice}

\begin{exercice}\label{exo:fact-10}
Soit $A$ un anneau factoriel. Montrer que $a \in A$ est irréductible
si et seulement si l'idéal $(a)$ est maximal parmi les idéaux
principaux propres de~$A$.
\end{exercice}

\begin{exercice}\label{exo:fact-11}
Factoriser $X^4 - 1$ dans $\Z[X]$, dans $\Q[X]$, dans $\R[X]$ et dans
$\C[X]$.
\end{exercice}

\begin{exercice}\label{exo:fact-12}
Montrer que pour tout entier $n \geqslant 2$, le polynôme
$f(X) = X^n - 2$ est irréductible dans $\Z[X]$. En déduire que
$[\Q(\sqrt[n]{2}) : \Q] = n$.
\end{exercice}

\begin{exercice}\label{exo:fact-13}
\index{anneau noethérien}
Un anneau commutatif~$A$ est dit \emph{noethérien}\index{noethérien}
si toute suite croissante d'idéaux $I_1 \subseteq I_2 \subseteq \cdots$
est stationnaire. Montrer que tout anneau principal est noethérien.
Donner un exemple d'anneau noethérien qui n'est pas principal.
\end{exercice}

%% ----------------------------------------------------------
\section*{Résumé du chapitre~\thechapter}

\begin{itemize}
  \item Dans un anneau intègre, tout élément premier est irréductible.
        La réciproque est vraie dans les anneaux principaux et les
        anneaux factoriels, mais fausse en général.
  \item Un anneau \emph{euclidien} possède une division euclidienne.
        Tout anneau euclidien est principal (les idéaux sont engendrés
        par leurs éléments de stathme minimal).
  \item Un anneau \emph{principal} (PID) est un anneau intègre dont tout
        idéal est principal. Tout PID est factoriel (argument par chaînes
        ascendantes pour l'existence, et primalité des irréductibles
        pour l'unicité).
  \item Un anneau \emph{factoriel} (UFD) possède la factorisation unique
        en irréductibles. Le lemme de Gauss montre que $A$ UFD entraîne
        $A[X]$ UFD.
  \item Hiérarchie stricte~: corps $\subsetneq$ euclidien $\subsetneq$
        PID $\subsetneq$ UFD $\subsetneq$ intègre.
  \item Le critère d'Eisenstein fournit une condition suffisante
        d'irréductibilité dans $A[X]$ pour $A$ factoriel.
\end{itemize}


%% ==========================================================
%%  BIBLIOGRAPHIE
%% ==========================================================

\begin{thebibliography}{99}

\bibitem{DummitFoote}
\textsc{Dummit}, D.\,S. et \textsc{Foote}, R.\,M.,
\emph{Abstract Algebra}, 3\textsuperscript{e} édition,
John Wiley \& Sons, Hoboken, 2004.

\bibitem{Lang}
\textsc{Lang}, S.,
\emph{Algebra}, 3\textsuperscript{e} édition,
Graduate Texts in Mathematics~\textbf{211},
Springer-Verlag, New York, 2002.

\bibitem{Artin}
\textsc{Artin}, M.,
\emph{Algebra}, 2\textsuperscript{e} édition,
Pearson, Boston, 2011.

\bibitem{Jacobson}
\textsc{Jacobson}, N.,
\emph{Basic Algebra~I}, 2\textsuperscript{e} édition,
W.\,H.\,Freeman and Company, New York, 1985.

\bibitem{Bourbaki}
\textsc{Bourbaki}, N.,
\emph{Algèbre}, Chapitres~1 à~3,
Springer-Verlag, Berlin, 2007 (réimpression).

\bibitem{Perrin}
\textsc{Perrin}, D.,
\emph{Algèbre~: cours de mathématiques de première année},
Ellipses, Paris, 1996.

\end{thebibliography}


\printindex

\end{document}
